%&preformat-disser
\RequirePackage[l2tabu,orthodox]{nag}  % Предупреждения об устаревших пакетах и командах
% Формат А4, 14pt (ГОСТ Р 7.0.11-2011, 5.3.6)
\documentclass[a4paper,14pt,oneside,openany]{memoir}

%%% Настройки и пакеты
%%%%%%%%%%%%%%%%%%%%%%%%%%%%%%%%%%%%%%%%%%%%%%%%%%%%%%%%%%%%%%%%%%%%%%%%%%%%%%%%
%%%% Общие настройки книги %%%%
%%%%%%%%%%%%%%%%%%%%%%%%%%%%%%%%%%%%%%%%%%%%%%%%%%%%%%%%%%%%%%%%%%%%%%%%%%%%%%%%

%%% Режим черновика %%%
\makeatletter
\@ifundefined{c@draft}{
  \newcounter{draft}
  \setcounter{draft}{0}  % 0 --- чистовик (максимальное соблюдение ГОСТ)
                         % 1 --- черновик (отклонения от ГОСТ, но быстрая
                         %       сборка итоговых PDF)
}{}
\makeatother

%%% Пометки в тексте %%%
\makeatletter
\@ifundefined{c@showmarkup}{
  \newcounter{showmarkup}
  \setcounter{showmarkup}{0}  % 0 --- скрыть пометки
                              % 1 --- показывать пометки
}{}
\makeatother

%%% Использование в pdflatex шрифтов не по-умолчанию %%%
\makeatletter
\@ifundefined{c@usealtfont}{
  \newcounter{usealtfont}
  \setcounter{usealtfont}{1}    % 0 --- шрифты на базе Computer Modern
                                % 1 --- использовать пакет pscyr, при его
                                %       наличии
                                % 2 --- использовать пакет XCharter, при наличии
                                %       подходящей версии
}{}
\makeatother

%%% Использование в xelatex и lualatex семейств шрифтов %%%
\makeatletter
\@ifundefined{c@fontfamily}{
  \newcounter{fontfamily}
  \setcounter{fontfamily}{1}  % 0 --- CMU семейство. Используется как fallback;
                              % 1 --- Шрифты от MS (Times New Roman и компания)
                              % 2 --- Семейство Liberation
}{}
\makeatother

%%% Библиография %%%
\makeatletter
\@ifundefined{c@bibliosel}{
  \newcounter{bibliosel}
  \setcounter{bibliosel}{1}   % 0 --- встроенная реализация с загрузкой файла
                              %       через движок bibtex8;
                              % 1 --- реализация пакетом biblatex через движок
                              %       biber
}{}
\makeatother

%%% Вывод типов ссылок в библиографии %%%
\makeatletter
\@ifundefined{c@mediadisplay}{
  \newcounter{mediadisplay}
  \setcounter{mediadisplay}{1}   % 0 --- не делать ничего; надписи [Текст] и
                                 %       [Эл. ресурс] будут выводиться только в ссылках с
                                 %       заполненным полем `media`;
                                 % 1 --- автоматически добавлять надпись [Текст] к ссылкам с
                                 %       незаполненным полем `media`; таким образом, у всех
                                 %       источников будет указан тип, что соответствует
                                 %       требованиям ГОСТ
                                 % 2 --- автоматически удалять надписи [Текст], [Эл. Ресурс] и др.;
                                 %       не соответствует ГОСТ
                                 % 3 --- автоматически удалять надпись [Текст];
                                 %       не соответствует ГОСТ
                                 % 4 --- автоматически удалять надпись [Эл. Ресурс];
                                 %       не соответствует ГОСТ
}{}
\makeatother

%%% Предкомпиляция tikz рисунков для ускорения работы %%%
\makeatletter
\@ifundefined{c@imgprecompile}{
  \newcounter{imgprecompile}
  \setcounter{imgprecompile}{0}   % 0 --- без предкомпиляции;
                                  % 1 --- пользоваться предварительно
                                  %       скомпилированными pdf вместо генерации
                                  %       заново из tikz
}{}
\makeatother
       % Общие настройки
%%% Проверка используемого TeX-движка %%%
\newif\ifxetexorluatex   % определяем новый условный оператор (http://tex.stackexchange.com/a/47579)
\ifxetex
    \xetexorluatextrue
\else
    \ifluatex
        \xetexorluatextrue
    \else
        \xetexorluatexfalse
    \fi
\fi

\usepackage{etoolbox}[2015/08/02]   % Для продвинутой проверки разных условий

\usepackage{comment}    % Позволяет убирать блоки текста (добавляет
                        % окружение comment и команду \excludecomment)

%%% Поля и разметка страницы %%%
\usepackage{pdflscape}  % Для включения альбомных страниц
\usepackage{geometry}   % Для последующего задания полей

%%% Математические пакеты %%%
\usepackage{amsthm,amsmath,amscd}   % Математические дополнения от AMS
\usepackage{amsfonts,amssymb}       % Математические дополнения от AMS
\usepackage{mathtools}              % Добавляет окружение multlined
\usepackage{xfrac}                  % Красивые дроби
\usepackage[
    locale = DE,
    list-separator       = {;\,},
    list-final-separator = {;\,},
    list-pair-separator  = {;\,},
    list-units           = single,
    range-units          = single,
    range-phrase={\text{\ensuremath{-}}},
    % quotient-mode        = fraction, % красивые дроби могут не соответствовать ГОСТ
    fraction-function    = \sfrac,
    separate-uncertainty,
    ]{siunitx}[=v2]                 % Размерности SI
\sisetup{inter-unit-product = \ensuremath{{}\cdot{}}}

% Кириллица в нумерации subequations
% Для правильной работы требуется выполнение сразу после загрузки пакетов
\patchcmd{\subequations}{\def\theequation{\theparentequation\alph{equation}}}
{\def\theequation{\theparentequation\asbuk{equation}}}
{\typeout{subequations patched}}{\typeout{subequations not patched}}

%%%% Установки для размера шрифта 14 pt %%%%
%% Формирование переменных и констант для сравнения (один раз для всех подключаемых файлов)%%
%% должно располагаться до вызова пакета fontspec или polyglossia, потому что они сбивают его работу
\newlength{\curtextsize}
\newlength{\bigtextsize}
\setlength{\bigtextsize}{13.9pt}

\makeatletter
%\show\f@size    % неплохо для отслеживания, но вызывает стопорение процесса,
                 % если документ компилируется без команды  -interaction=nonstopmode
\setlength{\curtextsize}{\f@size pt}
\makeatother

%%% Кодировки и шрифты %%%
\ifxetexorluatex
    \PassOptionsToPackage{no-math}{fontspec}    % https://tex.stackexchange.com/a/26295/104425
    \usepackage{polyglossia}[2014/05/21]        % Поддержка многоязычности
                                        % (fontspec подгружается автоматически)
\else
   %%% Решение проблемы копирования текста в буфер кракозябрами
    \ifnumequal{\value{usealtfont}}{0}{}{
        \input glyphtounicode.tex
        \input glyphtounicode-cmr.tex %from pdfx package
        \pdfgentounicode=1
    }
    \usepackage{cmap}   % Улучшенный поиск русских слов в полученном pdf-файле
    \ifnumequal{\value{usealtfont}}{2}{}{
        \defaulthyphenchar=127  % Если стоит до fontenc, то переносы
                                % не впишутся в выделяемый текст при
                                % копировании его в буфер обмена
    }
    \usepackage{textcomp}
    \usepackage[T1,T2A]{fontenc}                    % Поддержка русских букв
    \ifnumequal{\value{usealtfont}}{1}{% Используется pscyr, при наличии
        \IfFileExists{pscyr.sty}{\usepackage{pscyr}}{}  % Подключение pscyr
    }{}
    \usepackage[utf8]{inputenc}[2014/04/30]         % Кодировка utf8
    \usepackage[english, russian]{babel}[2014/03/24]% Языки: русский, английский
    \makeatletter\AtBeginDocument{\let\@elt\relax}\makeatother % babel 3.40 fix
    \ifnumequal{\value{usealtfont}}{2}{
        % http://dxdy.ru/post1238763.html#p1238763
        \usepackage[scaled=0.914]{XCharter}[2017/12/19] % Подключение русифицированных шрифтов XCharter
        \usepackage[charter, vvarbb, scaled=1.048]{newtxmath}[2017/12/14]
        \setDisplayskipStretch{-0.078}
    }{}
\fi

%%% Оформление абзацев %%%
\indentafterchapter     % Красная строка после заголовков типа chapter
\usepackage{indentfirst}

%%% Цвета %%%
\usepackage[dvipsnames, table]{xcolor} % Совместимо с tikz

%%% Таблицы %%%
\usepackage{longtable} % Длинные таблицы
\usepackage{multirow,makecell}   % Улучшенное форматирование таблиц
\usepackage{tabulary,tabularray} % Таблицы с автоматически подбирающейся
                                 % шириной столбцов
\UseTblrLibrary{booktabs}
\ExplSyntaxOn% define \IfTokenListEmpty to use \captionof with tabularray
\prg_generate_conditional_variant:Nnn \tl_if_empty:n { e } { TF }
\let \IfTokenListEmpty = \tl_if_empty:eTF
\ExplSyntaxOff

\usepackage{threeparttable}      % автоматический подгон ширины подписи таблицы

%%% Общее форматирование
%\usepackage{soul}% Поддержка переносоустойчивых подчёркиваний и зачёркиваний
\usepackage{icomma}  % Запятая в десятичных дробях

%%% Оптимизация расстановки переносов и длины последней строки абзаца
\IfFileExists{impnattypo.sty}{% проверка установленности пакета impnattypo
    \ifluatex
        \ifnumequal{\value{draft}}{1}{% Черновик
            \usepackage[hyphenation, lastparline, nosingleletter, homeoarchy,
            rivers, draft]{impnattypo}
        }{% Чистовик
            \usepackage[hyphenation, lastparline, nosingleletter]{impnattypo}
        }
    \else
        \usepackage[hyphenation, lastparline]{impnattypo}
    \fi
}{}

%% Векторная графика

\usepackage{tikz}                   % Продвинутый пакет векторной графики
\usetikzlibrary{chains}             % Для примера tikz рисунка
\usetikzlibrary{shapes.geometric}   % Для примера tikz рисунка
\usetikzlibrary{shapes.symbols}     % Для примера tikz рисунка
\usetikzlibrary{arrows}             % Для примера tikz рисунка

\usepackage[european,cuteinductors]{circuitikz} % Электрические схемы
\usepackage{pgfplots}                           % Графики
\pgfplotsset{compat=newest}
\usepgfplotslibrary{groupplots,units}
\pgfkeys{/pgf/number format/.cd,use comma,1000 sep={}} % форматирование чисел в графиках

%%% Гиперссылки %%%
\ifxetexorluatex
    \let\CYRDZE\relax
\fi
\usepackage{hyperref}[2012/11/06]

%%% Изображения %%%
\usepackage{graphicx}[2014/04/25]   % Подключаем пакет работы с графикой
\usepackage{caption}                % Подписи рисунков и таблиц
\usepackage{subcaption}             % Подписи подрисунков и подтаблиц
\usepackage{pdfpages}               % Добавление внешних pdf файлов

%%% Счётчики %%%
% В LaTeX с 2026-06-01 псевдонимы счётчиков поддерживаются ядром.
\ifdefined\newcounteralias\else
    \usepackage{aliascnt}
    \let\newcounteralias\newaliascnt
\fi
\usepackage[figure,table]{totalcount}   % Счётчик рисунков и таблиц
\usepackage{totcount}   % Пакет создания счётчиков на основе последнего номера
                        % подсчитываемого элемента (может требовать дважды
                        % компилировать документ)
\usepackage{totpages}   % Счётчик страниц, совместимый с hyperref (ссылается
                        % на номер последней страницы). Желательно ставить
                        % последним пакетом в преамбуле

%%% Продвинутое управление групповыми ссылками (пока только формулами) %%%
\usepackage[russian]{cleveref} % Язык задан явно, так как cleveref не всегда определяет его по babel.

% Добавление возможности использования пробелов в \labelcref
% https://tex.stackexchange.com/a/340502/104425
\usepackage{kvsetkeys}
\makeatletter
\let\org@@cref\@cref
\renewcommand*{\@cref}[2]{%
    \edef\process@me{%
        \noexpand\org@@cref{#1}{\zap@space#2 \@empty}%
    }\process@me
}
\makeatother

\usepackage{placeins} % для \FloatBarrier

\ifnumequal{\value{draft}}{1}{% Черновик
    \usepackage[firstpage]{draftwatermark}
    \SetWatermarkText{DRAFT}
    \SetWatermarkFontSize{14pt}
    \SetWatermarkScale{15}
    \SetWatermarkAngle{45}
}{}

% если текущий процесс запущен библиотекой tikz-external, то прекомпиляция должна быть включена
\ifdefined\tikzexternalrealjob
    \setcounter{imgprecompile}{1}
\fi

\ifnumequal{\value{imgprecompile}}{1}{% Только если у нас включена предкомпиляция
    \usetikzlibrary{external}   % подключение возможности предкомпиляции
    \tikzexternalize[prefix=images/cache/,optimize command away=\includepdf] % activate! % здесь можно указать отдельную папку для скомпилированных файлов
    \ifxetex
        \tikzset{external/up to date check={diff}}
    \fi
}{}
    % Общие пакеты
%%% Прикладные пакеты %%%
%\usepackage{calc}               % Пакет для расчётов параметров, например длины

%%% Для добавления Стр. над номерами страниц в оглавлении
%%% http://tex.stackexchange.com/a/306950
\usepackage{afterpage}

%%% Списки %%%
\usepackage{enumitem}

%%% Оформление списка обозначений
\usepackage[intoc,notocbasic]{nomencl}
\makenomenclature
\setlength{\nomitemsep}{-.8\parsep}
   % Пакеты книги
% Листинги с исходным кодом программ
\usepackage{fancyvrb}
\usepackage{listings}
\lccode`\~=0\relax %Без этого хака из-за особенностей пакета listings перестают работать конструкции с \MakeLowercase и т. п. в (xe|lua)latex

% Русская традиция начертания греческих букв
\usepackage{upgreek} % прямые греческие ради русской традиции

%%% Микротипографика
%\ifnumequal{\value{draft}}{0}{% Только если у нас режим чистовика
%    \usepackage[final, babel, shrink=45]{microtype}[2016/05/14] % улучшает представление букв и слов в строках, может помочь при наличии отдельно висящих слов
%}{}

% Отметка о версии черновика на каждой странице
% Чтобы работало надо в своей локальной копии по инструкции
% https://www.ctan.org/pkg/gitinfo2 создать небходимые файлы в папке
% ./git/hooks
% If you’re familiar with tweaking git, you can probably work it out for
% yourself. If not, I suggest you follow these steps:
% 1. First, you need a git repository and working tree. For this example,
% let’s suppose that the root of the working tree is in ~/compsci
% 2. Copy the file post-xxx-sample.txt (which is in the same folder of
% your TEX distribution as this pdf) into the git hooks directory in your
% working copy. In our example case, you should end up with a file called
% ~/compsci/.git/hooks/post-checkout
% 3. If you’re using a unix-like system, don’t forget to make the file executable.
% Just how you do this is outside the scope of this manual, but one
% possible way is with commands such as this:
% chmod g+x post-checkout.
% 4. Test your setup with “git checkout master” (or another suitable branch
% name). This should generate copies of gitHeadInfo.gin in the directories
% you intended.
% 5. Now make two more copies of this file in the same directory (hooks),
% calling them post-commit and post-merge, and you’re done. As before,
% users of unix-like systems should ensure these files are marked as
% executable.
\ifnumequal{\value{draft}}{1}{% Черновик
   \IfFileExists{.git/gitHeadInfo.gin}{
      \usepackage[mark,pcount]{gitinfo2}
      \renewcommand{\gitMark}{rev.\gitAbbrevHash\quad\gitCommitterEmail\quad\gitAuthorIsoDate}
      \renewcommand{\gitMarkFormat}{\rmfamily\color{Gray}\small\bfseries}
   }{}
}{}
  % Дополнительные пакеты
%%%%%%%%%%%%%%%%%%%%%%%%%%%%%%%%%%%%%%%%%%%%%%%%%%%%%%
%%%% Файл упрощённых настроек шаблона книги %%%%
%%%%%%%%%%%%%%%%%%%%%%%%%%%%%%%%%%%%%%%%%%%%%%%%%%%%%%

%%% Инициализирование переменных, не трогать!  %%%
\newcounter{intvl}
\newcounter{otstup}
\newcounter{contnumeq}
\newcounter{contnumfig}
\newcounter{contnumtab}
\newcounter{pgnum}
\newcounter{chapstyle}
\newcounter{headingdelim}
\newcounter{headingalign}
\newcounter{headingsize}
%%%%%%%%%%%%%%%%%%%%%%%%%%%%%%%%%%%%%%%%%%%%%%%%%%%%%%

%%% Область упрощённого управления оформлением %%%

%% Интервал между заголовками и между заголовком и текстом %%
% Заголовки отделяют от текста сверху и снизу
% тремя интервалами (ГОСТ Р 7.0.11-2011, 5.3.5)
\setcounter{intvl}{3}               % Коэффициент кратности к размеру шрифта

%% Отступы у заголовков в тексте %%
\setcounter{otstup}{0}              % 0 --- без отступа; 1 --- абзацный отступ

%% Нумерация формул, таблиц и рисунков %%
% Нумерация формул
\setcounter{contnumeq}{0}   % 0 --- пораздельно (во введении подряд,
                            %       без номера раздела);
                            % 1 --- сквозная нумерация по всей книге
% Нумерация рисунков
\setcounter{contnumfig}{0}  % 0 --- пораздельно (во введении подряд,
                            %       без номера раздела);
                            % 1 --- сквозная нумерация по всей книге
% Нумерация таблиц
\setcounter{contnumtab}{1}  % 0 --- пораздельно (во введении подряд,
                            %       без номера раздела);
                            % 1 --- сквозная нумерация по всей книге

%% Оглавление %%
\setcounter{pgnum}{1}       % 0 --- номера страниц никак не обозначены;
                            % 1 --- Стр. над номерами страниц (дважды
                            %       компилировать после изменения настройки)
\settocdepth{subsection}    % до какого уровня подразделов выносить в оглавление
\setsecnumdepth{subsection} % до какого уровня нумеровать подразделы


%% Текст и форматирование заголовков %%
\setcounter{chapstyle}{1}     % 0 --- разделы только под номером;
                              % 1 --- разделы с названием "Глава" перед номером
\setcounter{headingdelim}{1}  % 0 --- номер отделен пропуском в 1em или \quad;
                              % 1 --- номера разделов и приложений отделены
                              %       точкой с пробелом, подразделы пропуском
                              %       без точки;
                              % 2 --- номера разделов, подразделов и приложений
                              %       отделены точкой с пробелом.

%% Выравнивание заголовков в тексте %%
\setcounter{headingalign}{0}  % 0 --- по центру;
                              % 1 --- по левому краю

%% Размеры заголовков в тексте %%
\setcounter{headingsize}{0}   % 0 --- по ГОСТ, все всегда 14 пт;
                              % 1 --- пропорционально изменяющийся размер
                              %       в зависимости от базового шрифта

%% Подпись таблиц %%

% Смещение строк подписи после первой строки
\newcommand{\tabindent}{0cm}

% Тип форматирования заголовка таблицы:
% plain --- название и текст в одной строке
% split --- название и текст в разных строках
\newcommand{\tabformat}{plain}

%%% Настройки форматирования таблицы `plain`

% Выравнивание по центру подписи, состоящей из одной строки:
% true  --- выравнивать
% false --- не выравнивать
\newcommand{\tabsinglecenter}{false}

% Выравнивание подписи таблиц:
% justified   --- выравнивать как обычный текст («по ширине»)
% centering   --- выравнивать по центру
% centerlast  --- выравнивать по центру только последнюю строку
% centerfirst --- выравнивать по центру только первую строку (не рекомендуется)
% raggedleft  --- выравнивать по правому краю
% raggedright --- выравнивать по левому краю
\newcommand{\tabjust}{justified}

% Разделитель записи «Таблица #» и названия таблицы
\newcommand{\tablabelsep}{~\cyrdash\ }

%%% Настройки форматирования таблицы `split`

% Положение названия таблицы:
% \centering   --- выравнивать по центру
% \raggedleft  --- выравнивать по правому краю
% \raggedright --- выравнивать по левому краю
\newcommand{\splitformatlabel}{\raggedleft}

% Положение текста подписи:
% \centering   --- выравнивать по центру
% \raggedleft  --- выравнивать по правому краю
% \raggedright --- выравнивать по левому краю
\newcommand{\splitformattext}{\raggedright}

%% Подпись рисунков %%
%Разделитель записи «Рисунок #» и названия рисунка
\newcommand{\figlabelsep}{~\cyrdash\ }  % (ГОСТ 2.105, 4.3.1)
                                        % "--- здесь не работает

%%% Цвета гиперссылок %%%
% Latex color definitions: http://latexcolor.com/
\definecolor{linkcolor}{rgb}{0.9,0,0}
\definecolor{citecolor}{rgb}{0,0.6,0}
\definecolor{urlcolor}{rgb}{0,0,1}
%\definecolor{linkcolor}{rgb}{0,0,0} %black
%\definecolor{citecolor}{rgb}{0,0,0} %black
%\definecolor{urlcolor}{rgb}{0,0,0} %black
         % Настройки книги

%%% Данные издания и оформление
% Новые переменные, которые могут использоваться во всём проекте
%%% Заголовки библиографии:

\newcommand{\bibtitleauthor}{Публикации автора по теме книги}

% для стиля библиографии `\insertbiblioauthorgrouped`
\newcommand{\bibtitleauthorvak}{В изданиях из списка ВАК РФ}
\newcommand{\bibtitleauthorscopus}{В изданиях, входящих в международную базу цитирования Scopus}
\newcommand{\bibtitleauthorwos}{В изданиях, входящих в международную базу цитирования Web of Science}
\newcommand{\bibtitleauthorother}{В прочих изданиях}
\newcommand{\bibtitleauthorconf}{В сборниках трудов конференций}
\newcommand{\bibtitleauthorpatent}{Зарегистрированные патенты}
\newcommand{\bibtitleauthorprogram}{Зарегистрированные программы для ЭВМ}

% для списка литературы:
\newcommand{\bibtitlefull}{Список литературы} % (ГОСТ Р 7.0.11-2011, 4)
    % Названия и подписи
%%% Основные сведения %%%
\newcommand{\thesisAuthorLastName}{Киселёв}
\newcommand{\thesisAuthorOtherNames}{Андрей Петрович}
\newcommand{\thesisAuthorInitials}{А.\,П.}
\newcommand{\thesisAuthor}             % Автор издания
{%
    \texorpdfstring{% \texorpdfstring takes two arguments and uses the first for (La)TeX and the second for pdf
        \thesisAuthorShort% запись на титульном листе
    }{%
        \thesisAuthorLastName, \thesisAuthorOtherNames% эта запись для свойств pdf-файла. В таком виде, если pdf будет обработан программами для сбора библиографических сведений, будет правильно представлена фамилия.
    }
}
\newcommand{\thesisAuthorShort}        % Книга, ФИО автора инициалами
{\thesisAuthorInitials~\thesisAuthorLastName}
%\newcommand{\thesisUdk}                % Книга, УДК
%{\fixme{xxx.xxx}}
\newcommand{\thesisTitleFormat}[2]{#1 #2}
\newcommand{\thesisTitle}              % Название издания
{\thesisTitleFormat{Элементарная}{геометрия}}
\newcommand{\thesisSpecialtyNumber}    % Номер специальности, если используется
{}
\newcommand{\thesisSpecialtyTitle}     % Раздел издания
{Стереометрия}
%% \newcommand{\thesisSpecialtyTwoNumber} % Книга, вторая специальность, номер
%% {\fixme{XX.XX.XX}}
%% \newcommand{\thesisSpecialtyTwoTitle}  % Книга, вторая специальность, название
%% {\fixme{Теория и~методика физического воспитания, спортивной тренировки,
%% оздоровительной и~адаптивной физической культуры}}
\newcommand{\thesisDegree}             % Название книги
{Прямые и плоскости}
\newcommand{\thesisDegreeShort}        % Номер книги
{Книга~I.}
\newcommand{\thesisCity}               % Город издания
{Москва}
\newcommand{\thesisYear}               % Год издания книги
{\the\year}
\newcommand{\thesisOrganization}       % Строка об исходном издании на титуле
{По 23-му изданию 1914 года}

%% \newcommand{\supervisorDead}{}           % Рисовать рамку вокруг фамилии
\newcommand{\supervisorFio}              % Автор современной переработки
{rocky\_versus}
\newcommand{\supervisorRegalia}          % Репозиторий проекта
{\href{https://github.com/rockyversus/kiselev-geometry-schoolbook}{\nolinkurl{github.com/rockyversus/kiselev-geometry-schoolbook}}}

%% \newcommand{\supervisorTwoDead}{}        % Рисовать рамку вокруг фамилии
%% \newcommand{\supervisorTwoFio}           % Второй научный руководитель, ФИО
%% {\fixme{Фамилия Имя Отчество}}
%% \newcommand{\supervisorTwoRegalia}       % Второй научный руководитель, регалии
%% {\fixme{уч. степень, уч. звание}}

\providecommand{\keywords}%            % Ключевые слова для метаданных PDF книги
{}
        % Данные для титула и метаданных
%%% Кодировки и шрифты %%%
\ifxetexorluatex
    % Язык по-умолчанию русский с поддержкой приятных команд пакета babel
    \setmainlanguage[babelshorthands=true]{russian}
    % Дополнительный язык = английский (в американской вариации по-умолчанию)
    \setotherlanguage{english}

    % Проверка существования шрифтов. Недоступна в pdflatex
    \ifnumequal{\value{fontfamily}}{1}{
        \IfFontExistsTF{Times New Roman}{}{\setcounter{fontfamily}{0}}
    }{}
    \ifnumequal{\value{fontfamily}}{2}{
        \IfFontExistsTF{LiberationSerif}{}{\setcounter{fontfamily}{0}}
    }{}

    \ifnumequal{\value{fontfamily}}{0}{% Семейство шрифтов CMU. Используется как fallback
        \setmonofont[
          BoldFont={cmuntb.otf},
          ItalicFont={cmunit.otf},
          BoldItalicFont={cmuntx.otf}
        ]{cmuntt.otf} % {CMU Typewriter Text} % моноширинный шрифт
        \newfontfamily\cyrillicfonttt[
          BoldFont={cmuntb.otf},
          ItalicFont={cmunit.otf},
          BoldItalicFont={cmuntx.otf}
        ]{cmuntt.otf} % {CMU Typewriter Text} % моноширинный шрифт для кириллицы
        \defaultfontfeatures{Ligatures=TeX}   % стандартные лигатуры TeX, замены нескольких дефисов на тире и т. п. Настройки моноширинного шрифта должны идти до этой строки, чтобы при врезках кода программ в коде не применялись лигатуры и замены дефисов
        \setmainfont[
          SlantedFont={cmunsl.otf},
          BoldSlantedFont={cmunbl.otf},
          BoldFont={cmunbx.otf},
          ItalicFont={cmunti.otf},
          BoldItalicFont={cmunbi.otf}
        ]{cmunrm.otf} % {CMU Serif}           % Шрифт с засечками
        \newfontfamily\cyrillicfont[
          SlantedFont={cmunsl.otf},
          BoldSlantedFont={cmunbl.otf},
          BoldFont={cmunbx.otf},
          ItalicFont={cmunti.otf},
          BoldItalicFont={cmunbi.otf}
        ]{cmunrm.otf}   % {CMU Serif}         % Шрифт с засечками для кириллицы
        \setsansfont[
          BoldFont={cmunsx.otf},
          ItalicFont={cmunsi.otf},
          BoldItalicFont={cmunso.otf}
        ]{cmunss.otf} % {CMU Sans Serif}      % Шрифт без засечек
        \newfontfamily\cyrillicfontsf[
          BoldFont={cmunsx.otf},
          ItalicFont={cmunsi.otf},
          BoldItalicFont={cmunso.otf}
        ]{cmunss.otf} % {CMU Sans Serif}      % Шрифт без засечек для кириллицы
    }

    \ifnumequal{\value{fontfamily}}{1}{                    % Семейство MS шрифтов
        \setmonofont{Courier New}                          % моноширинный шрифт
        \newfontfamily\cyrillicfonttt{Courier New}         % моноширинный шрифт для кириллицы
        \defaultfontfeatures{Ligatures=TeX}                % стандартные лигатуры TeX, замены нескольких дефисов на тире и т. п. Настройки моноширинного шрифта должны идти до этой строки, чтобы при врезках кода программ в коде не применялись лигатуры и замены дефисов
        \setmainfont{Times New Roman}                      % Шрифт с засечками
        \newfontfamily\cyrillicfont{Times New Roman}       % Шрифт с засечками для кириллицы
        \setsansfont{Arial}                                % Шрифт без засечек
        \newfontfamily\cyrillicfontsf{Arial}               % Шрифт без засечек для кириллицы
    }

    \ifnumequal{\value{fontfamily}}{2}{                    % Семейство шрифтов Liberation (https://pagure.io/liberation-fonts)
        \setmonofont{LiberationMono}[Scale=0.87] % моноширинный шрифт
        \newfontfamily\cyrillicfonttt{LiberationMono}[     % моноширинный шрифт для кириллицы
            Scale=0.87]
        \defaultfontfeatures{Ligatures=TeX}                % стандартные лигатуры TeX, замены нескольких дефисов на тире и т. п. Настройки моноширинного шрифта должны идти до этой строки, чтобы при врезках кода программ в коде не применялись лигатуры и замены дефисов
        \setmainfont{LiberationSerif}                      % Шрифт с засечками
        \newfontfamily\cyrillicfont{LiberationSerif}       % Шрифт с засечками для кириллицы
        \setsansfont{LiberationSans}                       % Шрифт без засечек
        \newfontfamily\cyrillicfontsf{LiberationSans}      % Шрифт без засечек для кириллицы
    }

\else
    \ifnumequal{\value{usealtfont}}{1}{% Используется pscyr, при наличии
        \IfFileExists{pscyr.sty}{\renewcommand{\rmdefault}{ftm}}{}
    }{}
\fi
       % Шрифты
%%% Шаблон %%%
\DeclareRobustCommand{\fixme}{\textcolor{red}}  % решаем проблему превращения
                                % названия цвета в результате \MakeUppercase,
                                % http://tex.stackexchange.com/a/187930,
                                % \DeclareRobustCommand protects \fixme
                                % from expanding inside \MakeUppercase
\AtBeginDocument{%
    \setlength{\parindent}{2.5em}                   % Абзацный отступ. Должен быть одинаковым по всему тексту и равен пяти знакам (ГОСТ Р 7.0.11-2011, 5.3.7).
}

%%% Таблицы %%%
\DeclareCaptionLabelSeparator{tabsep}{\tablabelsep} % нумерация таблиц
\DeclareCaptionFormat{split}{\splitformatlabel#1\par\splitformattext#3}

\captionsetup[table]{
        format=\tabformat,                % формат подписи (plain|hang)
        font=normal,                      % нормальные размер, цвет, стиль шрифта
        skip=.0pt,                        % отбивка под подписью
        parskip=.0pt,                     % отбивка между параграфами подписи
        position=above,                   % положение подписи
        justification=\tabjust,           % центровка
        indent=\tabindent,                % смещение строк после первой
        labelsep=tabsep,                  % разделитель
        singlelinecheck=\tabsinglecenter, % не выравнивать по центру, если умещается в одну строку
}

%%% Рисунки %%%
\DeclareCaptionLabelSeparator{figsep}{\figlabelsep} % нумерация рисунков

\captionsetup[figure]{
        format=plain,                     % формат подписи (plain|hang)
        font=normal,                      % нормальные размер, цвет, стиль шрифта
        skip=.0pt,                        % отбивка под подписью
        parskip=.0pt,                     % отбивка между параграфами подписи
        position=below,                   % положение подписи
        singlelinecheck=true,             % выравнивание по центру, если умещается в одну строку
        justification=centerlast,         % центровка
        labelsep=figsep,                  % разделитель
}

%%% Подписи подрисунков %%%
\DeclareCaptionSubType{figure}
\renewcommand\thesubfigure{\asbuk{subfigure}} % нумерация подрисунков
\DeclareCaptionFont{norm}{\fontsize{14pt}{16pt}\selectfont}
\newcommand{\subfigureskip}{0.pt}

\captionsetup[subfloat]{
        labelfont=norm,                 % нормальный размер подписей подрисунков
        textfont=norm,                  % нормальный размер подписей подрисунков
        labelsep=space,                 % разделитель
        labelformat=brace,              % одна скобка справа от номера
        justification=centering,        % центровка
        singlelinecheck=true,           % выравнивание по центру, если умещается в одну строку
        skip=\subfigureskip,            % отбивка над подписью
        parskip=.0pt,                   % отбивка между параграфами подписи
        position=below,                 % положение подписи
}

%%% Настройки ссылок на рисунки, таблицы и др. %%%
% Фикс проблемы со ссылками \subcaptionref в необновлённом TeXLive 2025 и memoir
\makeatletter
\@ifclasslater{memoir}{2025/09/22}{%
% Считаем, что версии memoir новее будут взаимодействовать с amsmath и label корректнее
}{
  \@ifclasslater{memoir}{2025/03/10}{%
  % Эта версия memoir некорректно взаимодействует с обновлённым amsmath
    \renewcommand{\memsubfig@captionpar}[2]{%
      \parbox[t]{#1}{\let\label=\memsub@label\@subcapsize\@contsubcstyle #2}}
  }{}
}
% Фикс проблемы со ссылками пакета subcaption и обновлённого amsmath
\@ifpackagelater{amsmath}{2025/01/11}{
% https://tex.stackexchange.com/a/746030
  \patchcmd\caption@subtypehook{\let\label\subcaption@label}
    {\let\label\subcaption@label\let\ltx@label\subcaption@label}{}{\fail}
}{}
\makeatother

% команды \cref...format отвечают за форматирование при помощи команды \cref
% команды \labelcref...format отвечают за форматирование при помощи команды \labelcref

\crefdefaultlabelformat{#2#1#3}

% Уравнение
\crefformat{equation}{(#2#1#3)} % одиночная ссылка с приставкой
\labelcrefformat{equation}{(#2#1#3)} % одиночная ссылка без приставки
\crefrangeformat{equation}{(#3#1#4) \cyrdash~(#5#2#6)} % диапазон ссылок с приставкой
\labelcrefrangeformat{equation}{(#3#1#4) \cyrdash~(#5#2#6)} % диапазон ссылок без приставки
\crefmultiformat{equation}{(#2#1#3)}{ и~(#2#1#3)}{, (#2#1#3)}{ и~(#2#1#3)} % перечисление ссылок с приставкой
\labelcrefmultiformat{equation}{(#2#1#3)}{ и~(#2#1#3)}{, (#2#1#3)}{ и~(#2#1#3)} % перечисление без приставки

% Подуравнение
\crefformat{subequation}{(#2#1#3)} % одиночная ссылка с приставкой
\labelcrefformat{subequation}{(#2#1#3)} % одиночная ссылка без приставки
\crefrangeformat{subequation}{(#3#1#4) \cyrdash~(#5#2#6)} % диапазон ссылок с приставкой
\labelcrefrangeformat{subequation}{(#3#1#4) \cyrdash~(#5#2#6)} % диапазон ссылок без приставки
\crefmultiformat{subequation}{(#2#1#3)}{ и~(#2#1#3)}{, (#2#1#3)}{ и~(#2#1#3)} % перечисление ссылок с приставкой
\labelcrefmultiformat{subequation}{(#2#1#3)}{ и~(#2#1#3)}{, (#2#1#3)}{ и~(#2#1#3)} % перечисление без приставки

% Глава
\crefformat{chapter}{#2#1#3} % одиночная ссылка с приставкой
\labelcrefformat{chapter}{#2#1#3} % одиночная ссылка без приставки
\crefrangeformat{chapter}{#3#1#4 \cyrdash~#5#2#6} % диапазон ссылок с приставкой
\labelcrefrangeformat{chapter}{#3#1#4 \cyrdash~#5#2#6} % диапазон ссылок без приставки
\crefmultiformat{chapter}{#2#1#3}{ и~#2#1#3}{, #2#1#3}{ и~#2#1#3} % перечисление ссылок с приставкой
\labelcrefmultiformat{chapter}{#2#1#3}{ и~#2#1#3}{, #2#1#3}{ и~#2#1#3} % перечисление без приставки

% Параграф
\crefformat{section}{#2#1#3} % одиночная ссылка с приставкой
\labelcrefformat{section}{#2#1#3} % одиночная ссылка без приставки
\crefrangeformat{section}{#3#1#4 \cyrdash~#5#2#6} % диапазон ссылок с приставкой
\labelcrefrangeformat{section}{#3#1#4 \cyrdash~#5#2#6} % диапазон ссылок без приставки
\crefmultiformat{section}{#2#1#3}{ и~#2#1#3}{, #2#1#3}{ и~#2#1#3} % перечисление ссылок с приставкой
\labelcrefmultiformat{section}{#2#1#3}{ и~#2#1#3}{, #2#1#3}{ и~#2#1#3} % перечисление без приставки

% Приложение
\crefformat{appendix}{#2#1#3} % одиночная ссылка с приставкой
\labelcrefformat{appendix}{#2#1#3} % одиночная ссылка без приставки
\crefrangeformat{appendix}{#3#1#4 \cyrdash~#5#2#6} % диапазон ссылок с приставкой
\labelcrefrangeformat{appendix}{#3#1#4 \cyrdash~#5#2#6} % диапазон ссылок без приставки
\crefmultiformat{appendix}{#2#1#3}{ и~#2#1#3}{, #2#1#3}{ и~#2#1#3} % перечисление ссылок с приставкой
\labelcrefmultiformat{appendix}{#2#1#3}{ и~#2#1#3}{, #2#1#3}{ и~#2#1#3} % перечисление без приставки

% Рисунок
\crefformat{figure}{#2#1#3} % одиночная ссылка с приставкой
\labelcrefformat{figure}{#2#1#3} % одиночная ссылка без приставки
\crefrangeformat{figure}{#3#1#4 \cyrdash~#5#2#6} % диапазон ссылок с приставкой
\labelcrefrangeformat{figure}{#3#1#4 \cyrdash~#5#2#6} % диапазон ссылок без приставки
\crefmultiformat{figure}{#2#1#3}{ и~#2#1#3}{, #2#1#3}{ и~#2#1#3} % перечисление ссылок с приставкой
\labelcrefmultiformat{figure}{#2#1#3}{ и~#2#1#3}{, #2#1#3}{ и~#2#1#3} % перечисление без приставки

% Таблица
\crefformat{table}{#2#1#3} % одиночная ссылка с приставкой
\labelcrefformat{table}{#2#1#3} % одиночная ссылка без приставки
\crefrangeformat{table}{#3#1#4 \cyrdash~#5#2#6} % диапазон ссылок с приставкой
\labelcrefrangeformat{table}{#3#1#4 \cyrdash~#5#2#6} % диапазон ссылок без приставки
\crefmultiformat{table}{#2#1#3}{ и~#2#1#3}{, #2#1#3}{ и~#2#1#3} % перечисление ссылок с приставкой
\labelcrefmultiformat{table}{#2#1#3}{ и~#2#1#3}{, #2#1#3}{ и~#2#1#3} % перечисление без приставки

% Листинг
\crefformat{lstlisting}{#2#1#3} % одиночная ссылка с приставкой
\labelcrefformat{lstlisting}{#2#1#3} % одиночная ссылка без приставки
\crefrangeformat{lstlisting}{#3#1#4 \cyrdash~#5#2#6} % диапазон ссылок с приставкой
\labelcrefrangeformat{lstlisting}{#3#1#4 \cyrdash~#5#2#6} % диапазон ссылок без приставки
\crefmultiformat{lstlisting}{#2#1#3}{ и~#2#1#3}{, #2#1#3}{ и~#2#1#3} % перечисление ссылок с приставкой
\labelcrefmultiformat{lstlisting}{#2#1#3}{ и~#2#1#3}{, #2#1#3}{ и~#2#1#3} % перечисление без приставки

% Листинг
\crefformat{ListingEnv}{#2#1#3} % одиночная ссылка с приставкой
\labelcrefformat{ListingEnv}{#2#1#3} % одиночная ссылка без приставки
\crefrangeformat{ListingEnv}{#3#1#4 \cyrdash~#5#2#6} % диапазон ссылок с приставкой
\labelcrefrangeformat{ListingEnv}{#3#1#4 \cyrdash~#5#2#6} % диапазон ссылок без приставки
\crefmultiformat{ListingEnv}{#2#1#3}{ и~#2#1#3}{, #2#1#3}{ и~#2#1#3} % перечисление ссылок с приставкой
\labelcrefmultiformat{ListingEnv}{#2#1#3}{ и~#2#1#3}{, #2#1#3}{ и~#2#1#3} % перечисление без приставки

%%% Настройки гиперссылок %%%
\ifluatex
    \hypersetup{
        unicode,                % Unicode encoded PDF strings
    }
\fi

\hypersetup{
    linktocpage=true,           % ссылки с номера страницы в оглавлении, списке таблиц и списке рисунков
%    linktoc=all,                % both the section and page part are links
%    pdfpagelabels=false,        % set PDF page labels (true|false)
    plainpages=false,           % Forces page anchors to be named by the Arabic form  of the page number, rather than the formatted form
    colorlinks,                 % ссылки отображаются раскрашенным текстом, а не раскрашенным прямоугольником, вокруг текста
    linkcolor={linkcolor},      % цвет ссылок типа ref, eqref и подобных
    citecolor={citecolor},      % цвет ссылок-цитат
    urlcolor={urlcolor},        % цвет гиперссылок
%    hidelinks,                  % Hide links (removing color and border)
    pdftitle={\thesisTitle},    % Заголовок
    pdfauthor={\thesisAuthor},  % Автор
    pdfsubject={\thesisSpecialtyNumber\ \thesisSpecialtyTitle},      % Тема
%    pdfcreator={Создатель},     % Создатель, Приложение
%    pdfproducer={Производитель},% Производитель, Производитель PDF
    pdfkeywords={\keywords},    % Ключевые слова
    pdflang={ru},
}
\ifnumequal{\value{draft}}{1}{% Черновик
    \hypersetup{
        draft,
    }
}{}

%%% Списки %%%
% Используем короткое тире (endash) для ненумерованных списков (ГОСТ 2.105-95, пункт 4.1.7, требует дефиса, но так лучше смотрится)
\renewcommand{\labelitemi}{\normalfont\bfseries{--}}

% Перечисление строчными буквами латинского алфавита (ГОСТ 2.105-95, 4.1.7)
%\renewcommand{\theenumi}{\alph{enumi}}
%\renewcommand{\labelenumi}{\theenumi)}

% Перечисление строчными буквами русского алфавита (ГОСТ 2.105-95, 4.1.7)
\makeatletter
\AddEnumerateCounter{\asbuk}{\russian@alph}{щ}      % Управляем списками/перечислениями через пакет enumitem, а он 'не знает' про asbuk, потому 'учим' его
\makeatother
%\renewcommand{\theenumi}{\asbuk{enumi}} %первый уровень нумерации
%\renewcommand{\labelenumi}{\theenumi)} %первый уровень нумерации
\renewcommand{\theenumii}{\asbuk{enumii}} %второй уровень нумерации
\renewcommand{\labelenumii}{\theenumii)} %второй уровень нумерации
\renewcommand{\theenumiii}{\arabic{enumiii}} %третий уровень нумерации
\renewcommand{\labelenumiii}{\theenumiii)} %третий уровень нумерации

\setlist{nosep,%                                    % Единый стиль для всех списков (пакет enumitem), без дополнительных интервалов.
    labelindent=\parindent,leftmargin=*%            % Каждый пункт, подпункт и перечисление записывают с абзацного отступа (ГОСТ 2.105-95, 4.1.8)
}

%%% Правильная нумерация приложений, рисунков и формул %%%
%% По ГОСТ 2.105, п. 4.3.8 Приложения обозначают заглавными буквами русского алфавита,
%% начиная с А, за исключением букв Ё, З, Й, О, Ч, Ь, Ы, Ъ.
%% Здесь также переделаны все нумерации русскими буквами.
\ifxetexorluatex
    \makeatletter
    \def\russian@Alph#1{\ifcase#1\or
       А\or Б\or В\or Г\or Д\or Е\or Ж\or
       И\or К\or Л\or М\or Н\or
       П\or Р\or С\or Т\or У\or Ф\or Х\or
       Ц\or Ш\or Щ\or Э\or Ю\or Я\else\xpg@ill@value{#1}{russian@Alph}\fi}
    \def\russian@alph#1{\ifcase#1\or
       а\or б\or в\or г\or д\or е\or ж\or
       и\or к\or л\or м\or н\or
       п\or р\or с\or т\or у\or ф\or х\or
       ц\or ш\or щ\or э\or ю\or я\else\xpg@ill@value{#1}{russian@alph}\fi}
    \def\cyr@Alph#1{\ifcase#1\or
        А\or Б\or В\or Г\or Д\or Е\or Ж\or
        И\or К\or Л\or М\or Н\or
        П\or Р\or С\or Т\or У\or Ф\or Х\or
        Ц\or Ш\or Щ\or Э\or Ю\or Я\else\xpg@ill@value{#1}{cyr@Alph}\fi}
    \def\cyr@alph#1{\ifcase#1\or
        а\or б\or в\or г\or д\or е\or ж\or
        и\or к\or л\or м\or н\or
        п\or р\or с\or т\or у\or ф\or х\or
        ц\or ш\or щ\or э\or ю\or я\else\xpg@ill@value{#1}{cyr@alph}\fi}
    \makeatother
\else
    \makeatletter
    \if@uni@ode
      \def\russian@Alph#1{\ifcase#1\or
        А\or Б\or В\or Г\or Д\or Е\or Ж\or
        И\or К\or Л\or М\or Н\or
        П\or Р\or С\or Т\or У\or Ф\or Х\or
        Ц\or Ш\or Щ\or Э\or Ю\or Я\else\@ctrerr\fi}
    \else
      \def\russian@Alph#1{\ifcase#1\or
        \CYRA\or\CYRB\or\CYRV\or\CYRG\or\CYRD\or\CYRE\or\CYRZH\or
        \CYRI\or\CYRK\or\CYRL\or\CYRM\or\CYRN\or
        \CYRP\or\CYRR\or\CYRS\or\CYRT\or\CYRU\or\CYRF\or\CYRH\or
        \CYRC\or\CYRSH\or\CYRSHCH\or\CYREREV\or\CYRYU\or
        \CYRYA\else\@ctrerr\fi}
    \fi
    \if@uni@ode
      \def\russian@alph#1{\ifcase#1\or
        а\or б\or в\or г\or д\or е\or ж\or
        и\or к\or л\or м\or н\or
        п\or р\or с\or т\or у\or ф\or х\or
        ц\or ш\or щ\or э\or ю\or я\else\@ctrerr\fi}
    \else
      \def\russian@alph#1{\ifcase#1\or
        \cyra\or\cyrb\or\cyrv\or\cyrg\or\cyrd\or\cyre\or\cyrzh\or
        \cyri\or\cyrk\or\cyrl\or\cyrm\or\cyrn\or
        \cyrp\or\cyrr\or\cyrs\or\cyrt\or\cyru\or\cyrf\or\cyrh\or
        \cyrc\or\cyrsh\or\cyrshch\or\cyrerev\or\cyryu\or
        \cyrya\else\@ctrerr\fi}
    \fi
    \makeatother
\fi


%%http://www.linux.org.ru/forum/general/6993203#comment-6994589 (используется totcount)
\makeatletter
\def\formtotal#1#2#3#4#5{%
    \newcount\@c
    \@c\totvalue{#1}\relax
    \newcount\@last
    \newcount\@pnul
    \@last\@c\relax
    \divide\@last 10
    \@pnul\@last\relax
    \divide\@pnul 10
    \multiply\@pnul-10
    \advance\@pnul\@last
    \multiply\@last-10
    \advance\@last\@c
    #2%
    \ifnum\@pnul=1#5\else%
    \ifcase\@last#5\or#3\or#4\or#4\or#4\else#5\fi
    \fi
}
\makeatother

\newcommand{\formbytotal}[5]{\total{#1}~\formtotal{#1}{#2}{#3}{#4}{#5}}

%%% Команды рецензирования %%%
\ifboolexpr{ (test {\ifnumequal{\value{draft}}{1}}) or (test {\ifnumequal{\value{showmarkup}}{1}})}{
        \newrobustcmd{\todo}[1]{\textcolor{red}{#1}}
        \newrobustcmd{\note}[2][]{\ifstrempty{#1}{#2}{\textcolor{#1}{#2}}}
        \newenvironment{commentbox}[1][]%
        {\ifstrempty{#1}{}{\color{#1}}}%
        {}
}{
        \newrobustcmd{\todo}[1]{}
        \newrobustcmd{\note}[2][]{}
        \excludecomment{commentbox}
}
      % Общие стили
%%% Переопределение именований, если иначе не сработает %%%
%\gappto\captionsrussian{
%    \renewcommand{\chaptername}{Глава}
%    \renewcommand{\appendixname}{Приложение} % (ГОСТ Р 7.0.11-2011, 5.7)
%}

%%% Изображения %%%
\graphicspath{{images/}{Book/images/}}         % Пути к изображениям

%%% Интервалы %%%
%% По ГОСТ Р 7.0.11-2011, пункту 5.3.6 требуется полуторный интервал
%% Реализация средствами класса (на основе setspace) ближе к типографской классике.
%% И правит сразу и в таблицах (если со звёздочкой)
%\DoubleSpacing*     % Двойной интервал
\OnehalfSpacing*    % Полуторный интервал
%\setSpacing{1.42}   % Полуторный интервал, подобный Ворду (возможно, стоит включать вместе с предыдущей строкой)

%%% Макет страницы %%%
% Выставляем значения полей (ГОСТ 7.0.11-2011, 5.3.7)
\geometry{a4paper, top=2cm, bottom=2cm, left=2.5cm, right=1cm, nofoot, nomarginpar} %, heightrounded, showframe
\setlength{\topskip}{0pt}   %размер дополнительного верхнего поля
\setlength{\footskip}{12.3pt} % снимет warning, согласно https://tex.stackexchange.com/a/334346

%%% Выравнивание и переносы %%%
%% http://tex.stackexchange.com/questions/241343/what-is-the-meaning-of-fussy-sloppy-emergencystretch-tolerance-hbadness
%% http://www.latex-community.org/forum/viewtopic.php?p=70342#p70342
\tolerance 1414
\hbadness 1414
\emergencystretch 1.5em % В случае проблем регулировать в первую очередь
\hfuzz 0.3pt
\vfuzz \hfuzz
%\raggedbottom
%\sloppy                 % Избавляемся от переполнений
\clubpenalty=10000      % Запрещаем разрыв страницы после первой строки абзаца
\widowpenalty=10000     % Запрещаем разрыв страницы после последней строки абзаца
\brokenpenalty=4991     % Ограничение на разрыв страницы, если строка заканчивается переносом

%%% Блок управления параметрами для выравнивания заголовков в тексте %%%
\newlength{\otstuplen}
\setlength{\otstuplen}{\theotstup\parindent}
\ifnumequal{\value{headingalign}}{0}{% выравнивание заголовков в тексте
    \newcommand{\hdngalign}{\centering}                % по центру
    \newcommand{\hdngaligni}{}% по центру
    \setlength{\otstuplen}{0pt}
}{%
    \newcommand{\hdngalign}{}                 % по левому краю
    \newcommand{\hdngaligni}{\hspace{\otstuplen}}      % по левому краю
} % В обоих случаях вроде бы без переноса, как и надо (ГОСТ Р 7.0.11-2011, 5.3.5)

%%% Оглавление %%%
\renewcommand{\cftchapterdotsep}{\cftdotsep}                % отбивка точками до номера страницы начала главы/раздела

%% Переносить слова в заголовке не допускается (ГОСТ Р 7.0.11-2011, 5.3.5). Заголовки в оглавлении должны точно повторять заголовки в тексте (ГОСТ Р 7.0.11-2011, 5.2.3). Прямого указания на запрет переносов в оглавлении нет, но по той же логике невнесения искажений в смысл, лучше в оглавлении не переносить:
\setrmarg{2.55em plus1fil}                             %To have the (sectional) titles in the ToC, etc., typeset ragged right with no hyphenation
\renewcommand{\cftchapterpagefont}{\normalfont}        % нежирные номера страниц у глав в оглавлении
\renewcommand{\cftchapterleader}{\cftdotfill{\cftchapterdotsep}}% нежирные точки до номеров страниц у глав в оглавлении
%\renewcommand{\cftchapterfont}{}                       % нежирные названия глав в оглавлении

\ifnumgreater{\value{headingdelim}}{0}{%
    \renewcommand\cftchapteraftersnum{.\space}       % добавляет точку с пробелом после номера раздела в оглавлении
}{}
\ifnumgreater{\value{headingdelim}}{1}{%
    \renewcommand\cftsectionaftersnum{.\space}       % добавляет точку с пробелом после номера подраздела в оглавлении
    \renewcommand\cftsubsectionaftersnum{.\space}    % добавляет точку с пробелом после номера подподраздела в оглавлении
    \renewcommand\cftsubsubsectionaftersnum{.\space} % добавляет точку с пробелом после номера подподподраздела в оглавлении
    \AfterEndPreamble{% без этого polyglossia сама всё переопределяет
        \setsecnumformat{\csname the#1\endcsname.\space}
    }
}{%
    \AfterEndPreamble{% без этого polyglossia сама всё переопределяет
        \setsecnumformat{\csname the#1\endcsname\quad}
    }
}

\renewcommand*{\cftappendixname}{\appendixname\space} % Слово Приложение в оглавлении

%%% Колонтитулы %%%
% Порядковый номер страницы печатают на середине верхнего поля страницы (ГОСТ Р 7.0.11-2011, 5.3.8)
\makeevenhead{plain}{}{\rmfamily\thepage}{}
\makeoddhead{plain}{}{\rmfamily\thepage}{}
\makeevenfoot{plain}{}{}{}
\makeoddfoot{plain}{}{}{}
\pagestyle{plain}

%%% добавить Стр. над номерами страниц в оглавлении
%%% http://tex.stackexchange.com/a/306950
\newif\ifendTOC

\newcommand*{\tocheader}{
\ifnumequal{\value{pgnum}}{1}{%
    \ifendTOC\else\hbox to \linewidth%
      {\noindent{}~\hfill{Стр.}}\par%
      \ifnumless{\value{page}}{3}{}{%
        \vspace{0.5\onelineskip}
      }
      \afterpage{\tocheader}
    \fi%
}{}%
}%

%%% Оформление заголовков глав, разделов, подразделов %%%
%% Работа должна быть выполнена ... размером шрифта 12-14 пунктов (ГОСТ Р 7.0.11-2011, 5.3.8). То есть не должно быть надписей шрифтом более 14. Так и поставим.
%% Эти установки будут давать одинаковый результат независимо от выбора базовым шрифтом 12 пт или 14 пт
\newcommand{\basegostsectionfont}{\fontsize{14pt}{16pt}\selectfont\bfseries}

\makechapterstyle{thesisgost}{%
    \chapterstyle{default}
    \setlength{\beforechapskip}{0pt}
    \setlength{\midchapskip}{0pt}
    \setlength{\afterchapskip}{\theintvl\curtextsize}
    \renewcommand*{\chapnamefont}{\basegostsectionfont}
    \renewcommand*{\chapnumfont}{\basegostsectionfont}
    \renewcommand*{\chaptitlefont}{\basegostsectionfont}
    \renewcommand*{\chapterheadstart}{}
    \ifnumgreater{\value{headingdelim}}{0}{%
        \renewcommand*{\afterchapternum}{.\space}   % добавляет точку с пробелом после номера раздела
    }{%
        \renewcommand*{\afterchapternum}{\quad}     % добавляет \quad после номера раздела
    }
    \renewcommand*{\printchapternum}{\hdngaligni\hdngalign\chapnumfont \thechapter}
    \renewcommand*{\printchaptername}{}
    \renewcommand*{\printchapternonum}{\hdngaligni\hdngalign}
}

\makeatletter
\makechapterstyle{thesisgostchapname}{%
    \chapterstyle{thesisgost}
    \renewcommand*{\printchapternum}{\chapnumfont \thechapter}
    \renewcommand*{\printchaptername}{\hdngaligni\hdngalign\chapnamefont \@chapapp} %
}
\makeatother

\chapterstyle{thesisgost}

\setsecheadstyle{\basegostsectionfont\hdngalign}
\setsecindent{\otstuplen}

\setsubsecheadstyle{\basegostsectionfont\hdngalign}
\setsubsecindent{\otstuplen}

\setsubsubsecheadstyle{\basegostsectionfont\hdngalign}
\setsubsubsecindent{\otstuplen}

\sethangfrom{\noindent #1} %все заголовки подразделов центрируются с учетом номера, как block

\ifnumequal{\value{chapstyle}}{1}{%
    \chapterstyle{thesisgostchapname}
    \renewcommand*{\cftchaptername}{\chaptername\space} % будет вписано слово Глава перед каждым номером раздела в оглавлении
}{}%

%%% Интервалы между заголовками
\setbeforesecskip{\theintvl\curtextsize}% Заголовки отделяют от текста сверху и снизу тремя интервалами (ГОСТ Р 7.0.11-2011, 5.3.5).
\setaftersecskip{\theintvl\curtextsize}
\setbeforesubsecskip{\theintvl\curtextsize}
\setaftersubsecskip{\theintvl\curtextsize}
\setbeforesubsubsecskip{\theintvl\curtextsize}
\setaftersubsubsecskip{\theintvl\curtextsize}

%%% Вертикальные интервалы глав (\chapter) в оглавлении как и у заголовков
% раскомментировать следующие 2
% \setlength{\cftbeforechapterskip}{0pt plus 0pt}   % ИЛИ эти 2 строки из учебника
% \renewcommand*{\insertchapterspace}{}
% или эту
% \renewcommand*{\cftbeforechapterskip}{0em}


%%% Блок дополнительного управления размерами заголовков
\ifnumequal{\value{headingsize}}{1}{% Пропорциональные заголовки и базовый шрифт 14 пт
    \renewcommand{\basegostsectionfont}{\large\bfseries}
    \renewcommand*{\chapnamefont}{\Large\bfseries}
    \renewcommand*{\chapnumfont}{\Large\bfseries}
    \renewcommand*{\chaptitlefont}{\Large\bfseries}
}{}

%%% Счётчики %%%

%% Упрощённые настройки шаблона книги: нумерация формул, таблиц, рисунков
\ifnumequal{\value{contnumeq}}{1}{%
    \counterwithout{equation}{chapter} % Убираем связанность номера формулы с номером главы/раздела
}{}
\ifnumequal{\value{contnumfig}}{1}{%
    \counterwithout{figure}{chapter}   % Убираем связанность номера рисунка с номером главы/раздела
}{}
\ifnumequal{\value{contnumtab}}{1}{%
    \counterwithout{table}{chapter}    % Убираем связанность номера таблицы с номером главы/раздела
}{}

\AfterEndPreamble{
%% регистрируем счётчики в системе totcounter
    \regtotcounter{totalcount@figure}
    \regtotcounter{totalcount@table}       % Если иным способом поставить в преамбуле то ошибка в числе таблиц
    \regtotcounter{TotPages}               % Если иным способом поставить в преамбуле то ошибка в числе страниц
    \newtotcounter{totalchapter}
}
     % Стили книги
% Заголовки нумерованных глав по изданию Киселёва 1914 года.
\makechapterstyle{kiselev}{%
    \chapterstyle{thesisgost}
    \setlength{\midchapskip}{0.25\onelineskip}
    \setlength{\afterchapskip}{\onelineskip}
    \renewcommand*{\chapnamefont}{\normalfont\normalsize}
    \renewcommand*{\chapnumfont}{\normalfont\normalsize}
    \renewcommand*{\chaptitlefont}{\normalfont\large\bfseries}
    \renewcommand*{\printchaptername}{%
        \centering\chapnamefont
        Г\kern.4em Л\kern.4em А\kern.4em В\kern.4em А\enspace
    }
    \renewcommand*{\printchapternum}{\chapnumfont\thechapter.}
    \renewcommand*{\afterchapternum}{\par\nobreak\vskip\midchapskip}
    \renewcommand*{\printchaptertitle}[1]{\centering\chaptitlefont ##1.}
}

% Нумерованные пункты Киселёва в тексте с метками для перекрёстных ссылок.
\newcounter{kiselevparagraph}
\newcommand*{\kiselevparagraph}[1]{%
    \leavevmode
    \refstepcounter{kiselevparagraph}\label{#1}%
    \textbf{\thekiselevparagraph.}%
}

% для вертикального центрирования ячеек в tabulary
\def\zz{\ifx\[$\else\aftergroup\zzz\fi}
%$ \] % <-- чиним подсветку синтаксиса в некоторых редакторах
\def\zzz{\setbox0\lastbox
\dimen0\dimexpr\extrarowheight + \ht0-\dp0\relax
\setbox0\hbox{\raise-.5\dimen0\box0}%
\ht0=\dimexpr\ht0+\extrarowheight\relax
\dp0=\dimexpr\dp0+\extrarowheight\relax
\box0
}

\lstdefinelanguage{Renhanced}%
{keywords={abbreviate,abline,abs,acos,acosh,action,add1,add,%
        aggregate,alias,Alias,alist,all,anova,any,aov,aperm,append,apply,%
        approx,approxfun,apropos,Arg,args,array,arrows,as,asin,asinh,%
        atan,atan2,atanh,attach,attr,attributes,autoload,autoloader,ave,%
        axis,backsolve,barplot,basename,besselI,besselJ,besselK,besselY,%
        beta,binomial,body,box,boxplot,break,browser,bug,builtins,bxp,by,%
        c,C,call,Call,case,cat,category,cbind,ceiling,character,char,%
        charmatch,check,chol,chol2inv,choose,chull,class,close,cm,codes,%
        coef,coefficients,co,col,colnames,colors,colours,commandArgs,%
        comment,complete,complex,conflicts,Conj,contents,contour,%
        contrasts,contr,control,helmert,contrib,convolve,cooks,coords,%
        distance,coplot,cor,cos,cosh,count,fields,cov,covratio,wt,CRAN,%
        create,crossprod,cummax,cummin,cumprod,cumsum,curve,cut,cycle,D,%
        data,dataentry,date,dbeta,dbinom,dcauchy,dchisq,de,debug,%
        debugger,Defunct,default,delay,delete,deltat,demo,de,density,%
        deparse,dependencies,Deprecated,deriv,description,detach,%
        dev2bitmap,dev,cur,deviance,off,prev,,dexp,df,dfbetas,dffits,%
        dgamma,dgeom,dget,dhyper,diag,diff,digamma,dim,dimnames,dir,%
        dirname,dlnorm,dlogis,dnbinom,dnchisq,dnorm,do,dotplot,double,%
        download,dpois,dput,drop,drop1,dsignrank,dt,dummy,dump,dunif,%
        duplicated,dweibull,dwilcox,dyn,edit,eff,effects,eigen,else,%
        emacs,end,environment,env,erase,eval,equal,evalq,example,exists,%
        exit,exp,expand,expression,External,extract,extractAIC,factor,%
        fail,family,fft,file,filled,find,fitted,fivenum,fix,floor,for,%
        For,formals,format,formatC,formula,Fortran,forwardsolve,frame,%
        frequency,ftable,ftable2table,function,gamma,Gamma,gammaCody,%
        gaussian,gc,gcinfo,gctorture,get,getenv,geterrmessage,getOption,%
        getwd,gl,glm,globalenv,gnome,GNOME,graphics,gray,grep,grey,grid,%
        gsub,hasTsp,hat,heat,help,hist,home,hsv,httpclient,I,identify,if,%
        ifelse,Im,image,\%in\%,index,influence,measures,inherits,install,%
        installed,integer,interaction,interactive,Internal,intersect,%
        inverse,invisible,IQR,is,jitter,kappa,kronecker,labels,lapply,%
        layout,lbeta,lchoose,lcm,legend,length,levels,lgamma,library,%
        licence,license,lines,list,lm,load,local,locator,log,log10,log1p,%
        log2,logical,loglin,lower,lowess,ls,lsfit,lsf,ls,machine,Machine,%
        mad,mahalanobis,make,link,margin,match,Math,matlines,mat,matplot,%
        matpoints,matrix,max,mean,median,memory,menu,merge,methods,min,%
        missing,Mod,mode,model,response,mosaicplot,mtext,mvfft,na,nan,%
        names,omit,nargs,nchar,ncol,NCOL,new,next,NextMethod,nextn,%
        nlevels,nlm,noquote,NotYetImplemented,NotYetUsed,nrow,NROW,null,%
        numeric,\%o\%,objects,offset,old,on,Ops,optim,optimise,optimize,%
        options,or,order,ordered,outer,package,packages,page,pairlist,%
        pairs,palette,panel,par,parent,parse,paste,path,pbeta,pbinom,%
        pcauchy,pchisq,pentagamma,persp,pexp,pf,pgamma,pgeom,phyper,pico,%
        pictex,piechart,Platform,plnorm,plogis,plot,pmatch,pmax,pmin,%
        pnbinom,pnchisq,pnorm,points,poisson,poly,polygon,polyroot,pos,%
        postscript,power,ppoints,ppois,predict,preplot,pretty,Primitive,%
        print,prmatrix,proc,prod,profile,proj,prompt,prop,provide,%
        psignrank,ps,pt,ptukey,punif,pweibull,pwilcox,q,qbeta,qbinom,%
        qcauchy,qchisq,qexp,qf,qgamma,qgeom,qhyper,qlnorm,qlogis,qnbinom,%
        qnchisq,qnorm,qpois,qqline,qqnorm,qqplot,qr,Q,qty,qy,qsignrank,%
        qt,qtukey,quantile,quasi,quit,qunif,quote,qweibull,qwilcox,%
        rainbow,range,rank,rbeta,rbind,rbinom,rcauchy,rchisq,Re,read,csv,%
        csv2,fwf,readline,socket,real,Recall,rect,reformulate,regexpr,%
        relevel,remove,rep,repeat,replace,replications,report,require,%
        resid,residuals,restart,return,rev,rexp,rf,rgamma,rgb,rgeom,R,%
        rhyper,rle,rlnorm,rlogis,rm,rnbinom,RNGkind,rnorm,round,row,%
        rownames,rowsum,rpois,rsignrank,rstandard,rstudent,rt,rug,runif,%
        rweibull,rwilcox,sample,sapply,save,scale,scan,scan,screen,sd,se,%
        search,searchpaths,segments,seq,sequence,setdiff,setequal,set,%
        setwd,show,sign,signif,sin,single,sinh,sink,solve,sort,source,%
        spline,splinefun,split,sqrt,stars,start,stat,stem,step,stop,%
        storage,strstrheight,stripplot,strsplit,structure,strwidth,sub,%
        subset,substitute,substr,substring,sum,summary,sunflowerplot,svd,%
        sweep,switch,symbol,symbols,symnum,sys,status,system,t,table,%
        tabulate,tan,tanh,tapply,tempfile,terms,terrain,tetragamma,text,%
        time,title,topo,trace,traceback,transform,tri,trigamma,trunc,try,%
        ts,tsp,typeof,unclass,undebug,undoc,union,unique,uniroot,unix,%
        unlink,unlist,unname,untrace,update,upper,url,UseMethod,var,%
        variable,vector,Version,vi,warning,warnings,weighted,weights,%
        which,while,window,write,\%x\%,x11,X11,xedit,xemacs,xinch,xor,%
        xpdrows,xy,xyinch,yinch,zapsmall,zip},%
    otherkeywords={!,!=,~,$,*,\%,\&,\%/\%,\%*\%,\%\%,<-,<<-},%$
    alsoother={._$},%$
    sensitive,%
    morecomment=[l]\#,%
    morestring=[d]",%
    morestring=[d]'% 2001 Robert Denham
}%

%решаем проблему с кириллицей в комментариях (в pdflatex) https://tex.stackexchange.com/a/103712
\lstset{extendedchars=true,keepspaces=true,literate={Ö}{{\"O}}1
    {Ä}{{\"A}}1
    {Ü}{{\"U}}1
    {ß}{{\ss}}1
    {ü}{{\"u}}1
    {ä}{{\"a}}1
    {ö}{{\"o}}1
    {~}{{\textasciitilde}}1
    {а}{{\selectfont\char224}}1
    {б}{{\selectfont\char225}}1
    {в}{{\selectfont\char226}}1
    {г}{{\selectfont\char227}}1
    {д}{{\selectfont\char228}}1
    {е}{{\selectfont\char229}}1
    {ё}{{\"e}}1
    {ж}{{\selectfont\char230}}1
    {з}{{\selectfont\char231}}1
    {и}{{\selectfont\char232}}1
    {й}{{\selectfont\char233}}1
    {к}{{\selectfont\char234}}1
    {л}{{\selectfont\char235}}1
    {м}{{\selectfont\char236}}1
    {н}{{\selectfont\char237}}1
    {о}{{\selectfont\char238}}1
    {п}{{\selectfont\char239}}1
    {р}{{\selectfont\char240}}1
    {с}{{\selectfont\char241}}1
    {т}{{\selectfont\char242}}1
    {у}{{\selectfont\char243}}1
    {ф}{{\selectfont\char244}}1
    {х}{{\selectfont\char245}}1
    {ц}{{\selectfont\char246}}1
    {ч}{{\selectfont\char247}}1
    {ш}{{\selectfont\char248}}1
    {щ}{{\selectfont\char249}}1
    {ъ}{{\selectfont\char250}}1
    {ы}{{\selectfont\char251}}1
    {ь}{{\selectfont\char252}}1
    {э}{{\selectfont\char253}}1
    {ю}{{\selectfont\char254}}1
    {я}{{\selectfont\char255}}1
    {А}{{\selectfont\char192}}1
    {Б}{{\selectfont\char193}}1
    {В}{{\selectfont\char194}}1
    {Г}{{\selectfont\char195}}1
    {Д}{{\selectfont\char196}}1
    {Е}{{\selectfont\char197}}1
    {Ё}{{\"E}}1
    {Ж}{{\selectfont\char198}}1
    {З}{{\selectfont\char199}}1
    {И}{{\selectfont\char200}}1
    {Й}{{\selectfont\char201}}1
    {К}{{\selectfont\char202}}1
    {Л}{{\selectfont\char203}}1
    {М}{{\selectfont\char204}}1
    {Н}{{\selectfont\char205}}1
    {О}{{\selectfont\char206}}1
    {П}{{\selectfont\char207}}1
    {Р}{{\selectfont\char208}}1
    {С}{{\selectfont\char209}}1
    {Т}{{\selectfont\char210}}1
    {У}{{\selectfont\char211}}1
    {Ф}{{\selectfont\char212}}1
    {Х}{{\selectfont\char213}}1
    {Ц}{{\selectfont\char214}}1
    {Ч}{{\selectfont\char215}}1
    {Ш}{{\selectfont\char216}}1
    {Щ}{{\selectfont\char217}}1
    {Ъ}{{\selectfont\char218}}1
    {Ы}{{\selectfont\char219}}1
    {Ь}{{\selectfont\char220}}1
    {Э}{{\selectfont\char221}}1
    {Ю}{{\selectfont\char222}}1
    {Я}{{\selectfont\char223}}1
    {і}{{\selectfont\char105}}1
    {ї}{{\selectfont\char168}}1
    {є}{{\selectfont\char185}}1
    {ґ}{{\selectfont\char160}}1
    {І}{{\selectfont\char73}}1
    {Ї}{{\selectfont\char136}}1
    {Є}{{\selectfont\char153}}1
    {Ґ}{{\selectfont\char128}}1
}

% Ширина текста минус ширина надписи 999
\newlength{\twless}
\newlength{\lmarg}
\setlength{\lmarg}{\widthof{999}}   % ширина надписи 999
\setlength{\twless}{\textwidth-\lmarg}

\lstset{ %
%    language=R,                     %  Язык указать здесь, если во всех листингах преимущественно один язык, в результате часть настроек может пойти только для этого языка
    numbers=left,                   % where to put the line-numbers
    numberstyle=\fontsize{12pt}{14pt}\selectfont\color{Gray},  % the style that is used for the line-numbers
    firstnumber=1,                  % в этой и следующей строках задаётся поведение нумерации 5, 10, 15...
    stepnumber=5,                   % the step between two line-numbers. If it's 1, each line will be numbered
    numbersep=5pt,                  % how far the line-numbers are from the code
    backgroundcolor=\color{white},  % choose the background color. You must add \usepackage{color}
    showspaces=false,               % show spaces adding particular underscores
    showstringspaces=false,         % underline spaces within strings
    showtabs=false,                 % show tabs within strings adding particular underscores
    frame=leftline,                 % adds a frame of different types around the code
    rulecolor=\color{black},        % if not set, the frame-color may be changed on line-breaks within not-black text (e.g. commens (green here))
    tabsize=2,                      % sets default tabsize to 2 spaces
    captionpos=t,                   % sets the caption-position to top
    breaklines=true,                % sets automatic line breaking
    breakatwhitespace=false,        % sets if automatic breaks should only happen at whitespace
%    title=\lstname,                 % show the filename of files included with \lstinputlisting;
    % also try caption instead of title
    basicstyle=\fontsize{12pt}{14pt}\selectfont\ttfamily,% the size of the fonts that are used for the code
%    keywordstyle=\color{blue},      % keyword style
    commentstyle=\color{ForestGreen}\emph,% comment style
    stringstyle=\color{Mahogany},   % string literal style
    escapeinside={\%*}{*)},         % if you want to add a comment within your code
    morekeywords={*,...},           % if you want to add more keywords to the set
    inputencoding=utf8,             % кодировка кода
    xleftmargin={\lmarg},           % Чтобы весь код и полоска с номерами строк была смещена влево, так чтобы цифры не вылезали за пределы текста слева
}

%http://tex.stackexchange.com/questions/26872/smaller-frame-with-listings
% Окружение, чтобы листинг был компактнее обведен рамкой, если она задается, а не на всю ширину текста
\makeatletter
\newenvironment{SmallListing}[1][]
{\lstset{#1}\VerbatimEnvironment\begin{VerbatimOut}{VerbEnv.tmp}}
{\end{VerbatimOut}\settowidth\@tempdima{%
        \lstinputlisting{VerbEnv.tmp}}
    \minipage{\@tempdima}\lstinputlisting{VerbEnv.tmp}\endminipage}
\makeatother

\DefineVerbatimEnvironment% с шрифтом 12 пт
{Verb}{Verbatim}
{fontsize=\fontsize{12pt}{14pt}\selectfont}

\newfloat[chapter]{ListingEnv}{lol}{Листинг}

\renewcommand{\lstlistingname}{Листинг}

%Общие счётчики окружений листингов
%http://tex.stackexchange.com/questions/145546/how-to-make-figure-and-listing-share-their-counter
% Если смешивать плавающие и не плавающие окружения, то могут быть проблемы с нумерацией
\makeatletter
\AfterEndPreamble{% https://tex.stackexchange.com/a/252682
    \let\c@ListingEnv\relax % drop existing counter "ListingEnv"
    \newcounteralias{ListingEnv}{lstlisting} % общий счётчик листингов
    \let\ftype@lstlisting\ftype@ListingEnv % give the floats the same precedence
}
\makeatother

% значок С++ — используйте команду \cpp
\newcommand{\cpp}{%
    C\nolinebreak\hspace{-.05em}%
    \raisebox{.2ex}{+}\nolinebreak\hspace{-.10em}%
    \raisebox{.2ex}{+}%
}

%%% Ради примера во второй главе
\let\originalepsilon\epsilon
\let\originalphi\phi
\let\originalkappa\kappa
\let\originalle\le
\let\originalleq\leq
\let\originalge\ge
\let\originalgeq\geq
\let\originalemptyset\emptyset
\let\originaltan\tan
\let\originalcot\cot
\let\originalcsc\csc

%%% Русская традиция начертания математических знаков
\renewcommand{\le}{\ensuremath{\leqslant}}
\renewcommand{\leq}{\ensuremath{\leqslant}}
\renewcommand{\ge}{\ensuremath{\geqslant}}
\renewcommand{\geq}{\ensuremath{\geqslant}}
\renewcommand{\emptyset}{\varnothing}

%%% Русская традиция начертания математических функций (на случай копирования из зарубежных источников)
\renewcommand{\tan}{\operatorname{tg}}
\renewcommand{\cot}{\operatorname{ctg}}
\renewcommand{\csc}{\operatorname{cosec}}

%%% Русская традиция начертания греческих букв (греческие буквы вертикальные, через пакет upgreek)
\renewcommand{\epsilon}{\ensuremath{\upvarepsilon}}   %  русская традиция записи
\renewcommand{\phi}{\ensuremath{\upvarphi}}
%\renewcommand{\kappa}{\ensuremath{\varkappa}}
\renewcommand{\alpha}{\upalpha}
\renewcommand{\beta}{\upbeta}
\renewcommand{\gamma}{\upgamma}
\renewcommand{\delta}{\updelta}
\renewcommand{\varepsilon}{\upvarepsilon}
\renewcommand{\zeta}{\upzeta}
\renewcommand{\eta}{\upeta}
\renewcommand{\theta}{\uptheta}
\renewcommand{\vartheta}{\upvartheta}
\renewcommand{\iota}{\upiota}
\renewcommand{\kappa}{\upkappa}
\renewcommand{\lambda}{\uplambda}
\renewcommand{\mu}{\upmu}
\renewcommand{\nu}{\upnu}
\renewcommand{\xi}{\upxi}
\renewcommand{\pi}{\uppi}
\renewcommand{\varpi}{\upvarpi}
\renewcommand{\rho}{\uprho}
%\renewcommand{\varrho}{\upvarrho}
\renewcommand{\sigma}{\upsigma}
%\renewcommand{\varsigma}{\upvarsigma}
\renewcommand{\tau}{\uptau}
\renewcommand{\upsilon}{\upupsilon}
\renewcommand{\varphi}{\upvarphi}
\renewcommand{\chi}{\upchi}
\renewcommand{\psi}{\uppsi}
\renewcommand{\omega}{\upomega}
    % Дополнительные стили

%%% Библиография
% Движок задаётся параметром bibliosel в common/setup.tex.
\ifnumequal{\value{bibliosel}}{0}{%
    %%% Реализация библиографии встроенными средствами посредством движка bibtex8 %%%

%%% Пакеты %%%
\usepackage{cite}                                   % Красивые ссылки на литературу


%%% Стили %%%
\bibliographystyle{BibTeX-Styles/utf8gost71u}    % Оформляем библиографию по ГОСТ 7.1 (ГОСТ Р 7.0.11-2011, 5.6.7)

\makeatletter
\renewcommand{\@biblabel}[1]{#1.}   % Заменяем библиографию с квадратных скобок на точку
\makeatother
%% Управление отступами между записями
%% требует etoolbox
%% http://tex.stackexchange.com/a/105642
%\patchcmd\thebibliography
% {\labelsep}
% {\labelsep\itemsep=5pt\parsep=0pt\relax}
% {}
% {\typeout{Couldn't patch the command}}

%%% Список литературы с красной строки (без висячего отступа) %%%
%\patchcmd{\thebibliography} %может потребовать включения пакета etoolbox
%  {\advance\leftmargin\labelsep}
%  {\leftmargin=0pt%
%   \setlength{\labelsep}{\widthof{\ }}% Управляет длиной отступа после точки
%   \itemindent=\parindent%
%   \addtolength{\itemindent}{\labelwidth}% Сдвигаем правее на величину номера с точкой
%   \advance\itemindent\labelsep%
%  }
%  {}{}

%%% Цитирование %%%
\renewcommand\citepunct{;\penalty\citepunctpenalty%
    \hskip.13emplus.1emminus.1em\relax}                % Разделение ; при перечислении ссылок (ГОСТ Р 7.0.5-2008)

\newcommand*{\autocite}[1]{}  % Чтобы примеры цитирования, рассчитанные на biblatex, не вызывали ошибок при компиляции в bibtex

%%% Создание команд для вывода списка литературы %%%
\newcommand*{\insertbibliofull}{
\bibliography{biblio/external,biblio/author}         % Подключаем BibTeX-базы % После запятых не должно быть лишних пробелов — он "думает", что это тоже имя пути
}

\newcommand*{\insertbiblioauthor}{
\bibliography{biblio/author}         % Подключаем BibTeX-базы % После запятых не должно быть лишних пробелов — он "думает", что это тоже имя пути
}

\newcommand*{\insertbiblioexternal}{
\bibliography{biblio/external}         % Подключаем BibTeX-базы
}


%% Счётчик использованных ссылок на литературу, обрабатывающий с учётом неоднократных ссылок
%% Требуется дважды компилировать, поскольку ему нужно считать актуальный внешний файл со списком литературы
\newtotcounter{citenum}
\def\oldcite{}
\let\oldcite=\bibcite
\def\bibcite{\stepcounter{citenum}\oldcite}
  % BibTeX8
}{%
    %%% Реализация библиографии пакетами biblatex и biblatex-gost с использованием движка biber %%%

\usepackage{csquotes} % biblatex рекомендует его подключать. Пакет для оформления сложных блоков цитирования.
%%% Загрузка пакета с основными настройками %%%
\makeatletter
\ifnumequal{\value{draft}}{0}{% Чистовик
\usepackage[%
backend=biber,% движок
bibencoding=utf8,% кодировка bib файла
sorting=none,% настройка сортировки списка литературы
style=gost-numeric,% стиль цитирования и библиографии (по ГОСТ)
language=autobib,% получение языка из babel/polyglossia, default: autobib % если ставить autocite или auto, то цитаты в тексте с указанием страницы, получат указание страницы на языке оригинала
autolang=other,% многоязычная библиография
clearlang=true,% внутренний сброс поля language, если он совпадает с языком из babel/polyglossia
defernumbers=true,% нумерация проставляется после двух компиляций, зато позволяет выцеплять библиографию по ключевым словам и нумеровать не из большего списка
sortcites=true,% сортировать номера затекстовых ссылок при цитировании (если в квадратных скобках несколько ссылок, то отображаться будут отсортированно, а не абы как)
doi=false,% Показывать или нет ссылки на DOI
isbn=false,% Показывать или нет ISBN, ISSN, ISRN
]{biblatex}[2016/09/17]
\ltx@iffilelater{biblatex-gost.def}{2017/05/03}%
{\toggletrue{bbx:gostbibliography}%
\renewcommand*{\revsdnamepunct}{\addcomma}}{}
}{%Черновик
\usepackage[%
backend=biber,% движок
bibencoding=utf8,% кодировка bib файла
sorting=none,% настройка сортировки списка литературы
% defernumbers=true, % откомментируйте, если требуется правильная нумерация ссылок на литературу в режиме черновика. Замедляет сборку
]{biblatex}[2016/09/17]%
}
\makeatother

\providebool{blxmc} % biblatex version needs and has MakeCapital workaround
\boolfalse{blxmc} % setting our new boolean flag to default false
\ifxetexorluatex
\else
% Исправление случая неподдержки знака номера в pdflatex
    \DefineBibliographyStrings{russian}{number={\textnumero}}

% Исправление случая отсутствия прописных букв в некоторых случаях
% https://github.com/plk/biblatex/issues/960#issuecomment-596658282
    \ifdefmacro{\ExplSyntaxOn}{}{\usepackage{expl3}}
    \makeatletter
    \ltx@ifpackagelater{biblatex}{2020/02/23}{
    % Assuming this version of biblatex defines MakeCapital correctly
    }{
        \ltx@ifpackagelater{biblatex}{2019/12/01}{
            % Assuming this version of biblatex defines MakeCapital incorrectly
            \usepackage{expl3}[2020/02/25]
            \@ifpackagelater{expl3}{2020/02/25}{
                \booltrue{blxmc} % setting our new boolean flag to true
            }{}
        }{}
    }
    \makeatother
    \ifblxmc
        \typeout{Assuming this version of biblatex defines MakeCapital
        incorrectly}
        \usepackage{xparse}
        \makeatletter
        \ExplSyntaxOn
        \NewDocumentCommand \blx@maketext@lowercase {m}
          {
            \text_lowercase:n {#1}
          }

        \NewDocumentCommand \blx@maketext@uppercase {m}
          {
            \text_uppercase:n {#1}
          }

        \RenewDocumentCommand \MakeCapital {m}
          {
            \text_titlecase_first:n {#1}
          }
        \ExplSyntaxOff

        \protected\def\blx@biblcstring#1#2#3{%
          \blx@begunit
          \blx@hyphenreset
          \blx@bibstringsimple
          \lowercase{\edef\blx@tempa{#3}}%
          \ifcsundef{#2@\blx@tempa}
            {\blx@warn@nostring\blx@tempa
             \blx@endnounit}
            {#1{\blx@maketext@lowercase{\csuse{#2@\blx@tempa}}}%
             \blx@endunit}}

        \protected\def\blx@bibucstring#1#2#3{%
          \blx@begunit
          \blx@hyphenreset
          \blx@bibstringsimple
          \lowercase{\edef\blx@tempa{#3}}%
          \ifcsundef{#2@\blx@tempa}
            {\blx@warn@nostring\blx@tempa
             \blx@endnounit}
            {#1{\blx@maketext@uppercase{\csuse{#2@\blx@tempa}}}%
             \blx@endunit}}
        \makeatother
    \fi
\fi

%%% Подключение файлов bib %%%
\addbibresource[label=bl-external]{biblio/external.bib}
\addbibresource[label=bl-author]{biblio/author.bib}
\addbibresource[label=bl-registered]{biblio/registered.bib}

%http://tex.stackexchange.com/a/141831/79756
%There is a way to automatically map the language field to the langid field. The following lines in the preamble should be enough to do that.
%This command will copy the language field into the langid field and will then delete the contents of the language field. The language field will only be deleted if it was successfully copied into the langid field.
\DeclareSourcemap{ %модификация bib файла перед тем, как им займётся biblatex
    \maps{
        \map{% перекидываем значения полей language в поля langid, которыми пользуется biblatex
            \step[fieldsource=language, fieldset=langid, origfieldval, final]
            \step[fieldset=language, null]
        }
        \map{% перекидываем значения полей numpages в поля pagetotal, которыми пользуется biblatex
            \step[fieldsource=numpages, fieldset=pagetotal, origfieldval, final]
            \step[fieldset=numpages, null]
        }
        \map{% перекидываем значения полей pagestotal в поля pagetotal, которыми пользуется biblatex
            \step[fieldsource=pagestotal, fieldset=pagetotal, origfieldval, final]
            \step[fieldset=pagestotal, null]
        }
        \map[overwrite]{% перекидываем значения полей shortjournal, если они есть, в поля journal, которыми пользуется biblatex
            \step[fieldsource=shortjournal, final]
            \step[fieldset=journal, origfieldval]
            \step[fieldset=shortjournal, null]
        }
        \map[overwrite]{% перекидываем значения полей shortbooktitle, если они есть, в поля booktitle, которыми пользуется biblatex
            \step[fieldsource=shortbooktitle, final]
            \step[fieldset=booktitle, origfieldval]
            \step[fieldset=shortbooktitle, null]
        }
        \map{% если в поле medium написано "Электронный ресурс", то устанавливаем поле media, которым пользуется biblatex, в значение eresource.
            \step[fieldsource=medium,
            match=\regexp{Электронный\s+ресурс},
            final]
            \step[fieldset=media, fieldvalue=eresource]
            \step[fieldset=medium, null]
        }
        \map[overwrite]{% стираем значения всех полей issn
            \step[fieldset=issn, null]
        }
        \map[overwrite]{% стираем значения всех полей abstract, поскольку ими не пользуемся, а там бывают "неприятные" латеху символы
            \step[fieldsource=abstract]
            \step[fieldset=abstract,null]
        }
        \map[overwrite]{ % переделка формата записи даты
            \step[fieldsource=urldate,
            match=\regexp{([0-9]{2})\.([0-9]{2})\.([0-9]{4})},
            replace={$3-$2-$1$4}, % $4 вставлен исключительно ради нормальной работы программ подсветки синтаксиса, которые некорректно обрабатывают $ в таких конструкциях
            final]
        }
        \map[overwrite]{ % стираем ключевые слова
            \step[fieldsource=keywords]
            \step[fieldset=keywords,null]
        }
        % реализация foreach различается для biblatex v3.12 и v3.13.
        % Для версии v3.13 эта конструкция заменяет последующие 7 структур map
        % \map[overwrite,foreach={authorvak,authorscopus,authorwos,authorconf,authorother,authorparent,authorprogram}]{ % записываем информацию о типе публикации в ключевые слова
        %     \step[fieldsource=$MAPLOOP,final=true]
        %     \step[fieldset=keywords,fieldvalue={,biblio$MAPLOOP},append=true]
        % }
        \map[overwrite]{ % записываем информацию о типе публикации в ключевые слова
            \step[fieldsource=authorvak,final=true]
            \step[fieldset=keywords,fieldvalue={,biblioauthorvak},append=true]
        }
        \map[overwrite]{ % записываем информацию о типе публикации в ключевые слова
            \step[fieldsource=authorscopus,final=true]
            \step[fieldset=keywords,fieldvalue={,biblioauthorscopus},append=true]
        }
        \map[overwrite]{ % записываем информацию о типе публикации в ключевые слова
            \step[fieldsource=authorwos,final=true]
            \step[fieldset=keywords,fieldvalue={,biblioauthorwos},append=true]
        }
        \map[overwrite]{ % записываем информацию о типе публикации в ключевые слова
            \step[fieldsource=authorconf,final=true]
            \step[fieldset=keywords,fieldvalue={,biblioauthorconf},append=true]
        }
        \map[overwrite]{ % записываем информацию о типе публикации в ключевые слова
            \step[fieldsource=authorother,final=true]
            \step[fieldset=keywords,fieldvalue={,biblioauthorother},append=true]
        }
        \map[overwrite]{ % записываем информацию о типе публикации в ключевые слова
            \step[fieldsource=authorpatent,final=true]
            \step[fieldset=keywords,fieldvalue={,biblioauthorpatent},append=true]
        }
        \map[overwrite]{ % записываем информацию о типе публикации в ключевые слова
            \step[fieldsource=authorprogram,final=true]
            \step[fieldset=keywords,fieldvalue={,biblioauthorprogram},append=true]
        }
        \map[overwrite]{ % добавляем ключевые слова, чтобы различать источники
            \perdatasource{biblio/external.bib}
            \step[fieldset=keywords, fieldvalue={,biblioexternal},append=true]
        }
        \map[overwrite]{ % добавляем ключевые слова, чтобы различать источники
            \perdatasource{biblio/author.bib}
            \step[fieldset=keywords, fieldvalue={,biblioauthor},append=true]
        }
        \map[overwrite]{ % добавляем ключевые слова, чтобы различать источники
            \perdatasource{biblio/registered.bib}
            \step[fieldset=keywords, fieldvalue={,biblioregistered},append=true]
        }
        \map[overwrite]{ % добавляем ключевые слова, чтобы различать источники
            \step[fieldset=keywords, fieldvalue={,bibliofull},append=true]
        }
%        \map[overwrite]{% стираем значения всех полей series
%            \step[fieldset=series, null]
%        }
        \map[overwrite]{% перекидываем значения полей howpublished в поля organization для типа online
            \step[typesource=online, typetarget=online, final]
            \step[fieldsource=howpublished, fieldset=organization, origfieldval]
            \step[fieldset=howpublished, null]
        }
    }
}

\ifnumequal{\value{mediadisplay}}{1}{
    \DeclareSourcemap{
        \maps{%
            \map{% использование media=text по умолчанию
                \step[fieldset=media, fieldvalue=text]
            }
        }
    }
}{}
\ifnumequal{\value{mediadisplay}}{2}{
    \DeclareSourcemap{
        \maps{%
            \map[overwrite]{% удаление всех записей media
                \step[fieldset=media, null]
            }
        }
    }
}{}
\ifnumequal{\value{mediadisplay}}{3}{
    \DeclareSourcemap{
        \maps{
            \map[overwrite]{% стираем значения всех полей media=text
                \step[fieldsource=media,match={text},final]
                \step[fieldset=media, null]
            }
        }
    }
}{}
\ifnumequal{\value{mediadisplay}}{4}{
    \DeclareSourcemap{
        \maps{
            \map[overwrite]{% стираем значения всех полей media=eresource
                \step[fieldsource=media,match={eresource},final]
                \step[fieldset=media, null]
            }
        }
    }
}{}

\DeclareSourcemap{ %модификация bib файла перед тем, как им займётся biblatex
    \maps{
        \map[overwrite]{% стираем значения всех полей addendum
            \perdatasource{biblio/author.bib}
            \step[fieldset=addendum, null] % не выводим сведения об объёме авторских статей
        }
    }
}

\defbibfilter{vakscopuswos}{%
    keyword=biblioauthorvak or keyword=biblioauthorscopus or keyword=biblioauthorwos
}

\defbibfilter{scopuswos}{%
    keyword=biblioauthorscopus or keyword=biblioauthorwos
}

\defbibfilter{papersregistered}{%
    keyword=biblioauthor or keyword=biblioregistered
}

%%% Убираем неразрывные пробелы перед двоеточием и точкой с запятой %%%
%\makeatletter
%\ifnumequal{\value{draft}}{0}{% Чистовик
%    \renewcommand*{\addcolondelim}{%
%      \begingroup%
%      \def\abx@colon{%
%        \ifdim\lastkern>\z@\unkern\fi%
%        \abx@puncthook{:}\space}%
%      \addcolon%
%      \endgroup}
%
%    \renewcommand*{\addsemicolondelim}{%
%      \begingroup%
%      \def\abx@semicolon{%
%        \ifdim\lastkern>\z@\unkern\fi%
%        \abx@puncthook{;}\space}%
%      \addsemicolon%
%      \endgroup}
%}{}
%\makeatother

%%% Правка записей типа thesis, чтобы дважды не писался автор
%\ifnumequal{\value{draft}}{0}{% Чистовик
%\DeclareBibliographyDriver{thesis}{%
%  \usebibmacro{bibindex}%
%  \usebibmacro{begentry}%
%  \usebibmacro{heading}%
%  \newunit
%  \usebibmacro{author}%
%  \setunit*{\labelnamepunct}%
%  \usebibmacro{thesistitle}%
%  \setunit{\respdelim}%
%  %\printnames[last-first:full]{author}%Вот эту строчку нужно убрать, чтобы автор книги не дублировался
%  \newunit\newblock
%  \printlist[semicolondelim]{specdata}%
%  \newunit
%  \usebibmacro{institution+location+date}%
%  \newunit\newblock
%  \usebibmacro{chapter+pages}%
%  \newunit
%  \printfield{pagetotal}%
%  \newunit\newblock
%  \usebibmacro{doi+eprint+url+note}%
%  \newunit\newblock
%  \usebibmacro{addendum+pubstate}%
%  \setunit{\bibpagerefpunct}\newblock
%  \usebibmacro{pageref}%
%  \newunit\newblock
%  \usebibmacro{related:init}%
%  \usebibmacro{related}%
%  \usebibmacro{finentry}}
%}{}

%\newbibmacro{string+doi}[1]{% новая макрокоманда на простановку ссылки на doi
%    \iffieldundef{doi}{#1}{\href{http://dx.doi.org/\thefield{doi}}{#1}}}

%\ifnumequal{\value{draft}}{0}{% Чистовик
%\renewcommand*{\mkgostheading}[1]{\usebibmacro{string+doi}{#1}} % ссылка на doi с авторов. стоящих впереди записи
%\renewcommand*{\mkgostheading}[1]{#1} % только лишь убираем курсив с авторов
%}{}
%\DeclareFieldFormat{title}{\usebibmacro{string+doi}{#1}} % ссылка на doi с названия работы
%\DeclareFieldFormat{journaltitle}{\usebibmacro{string+doi}{#1}} % ссылка на doi с названия журнала
%%% Тире как разделитель в библиографии традиционной руской длины:
\renewcommand*{\newblockpunct}{\addperiod\addnbspace\cyrdash\space\bibsentence}
%%% Убрать тире из разделителей элементов в библиографии:
%\renewcommand*{\newblockpunct}{%
%    \addperiod\space\bibsentence}%block punct.,\bibsentence is for vol,etc.
%%% Изменение точки с запятой на запятую в перечислении библиографических
%%% ссылок:
%\renewcommand*{\multicitedelim}{\addcomma\space}

%%% Возвращаем запись «Режим доступа» %%%
%\DefineBibliographyStrings{english}{%
%    urlfrom = {Mode of access}
%}
%\DeclareFieldFormat{url}{\bibstring{urlfrom}\addcolon\space\url{#1}}

%%% В списке литературы обозначение одной буквой диапазона страниц англоязычного источника %%%
\DefineBibliographyStrings{english}{%
    pages = {p\adddot} %заглавность буквы затем по месту определяется работой самого biblatex
}

%%% В ссылке на источник в основном тексте с указанием конкретной страницы обозначение одной большой буквой %%%
%\DefineBibliographyStrings{russian}{%
%    page = {C\adddot}
%}

%%% Исправление длины тире в диапазонах %%%
% \cyrdash --- тире «русской» длины, \textendash --- en-dash
\DefineBibliographyExtras{russian}{%
  \protected\def\bibrangedash{%
    \cyrdash\penalty\value{abbrvpenalty}}% almost unbreakable dash
  \protected\def\bibdaterangesep{\bibrangedash}%тире для дат
}
\DefineBibliographyExtras{english}{%
  \protected\def\bibrangedash{%
    \cyrdash\penalty\value{abbrvpenalty}}% almost unbreakable dash
  \protected\def\bibdaterangesep{\bibrangedash}%тире для дат
}

%Set higher penalty for breaking in number, dates and pages ranges
\setcounter{abbrvpenalty}{10000} % default is \hyphenpenalty which is 12

%Set higher penalty for breaking in names
\setcounter{highnamepenalty}{10000} % If you prefer the traditional BibTeX behavior (no linebreaks at highnamepenalty breakpoints), set it to ‘infinite’ (10 000 or higher).
\setcounter{lownamepenalty}{10000}

%%% Set low penalties for breaks at uppercase letters and lowercase letters
%\setcounter{biburllcpenalty}{500} %управляет разрывами ссылок после маленьких букв RTFM biburllcpenalty
%\setcounter{biburlucpenalty}{3000} %управляет разрывами ссылок после больших букв, RTFM biburlucpenalty

%%% Список литературы с красной строки (без висячего отступа) %%%
%\defbibenvironment{bibliography} % переопределяем окружение библиографии из gost-numeric.bbx пакета biblatex-gost
%  {\list
%     {\printtext[labelnumberwidth]{%
%       \printfield{prefixnumber}%
%       \printfield{labelnumber}}}
%     {%
%      \setlength{\labelwidth}{\labelnumberwidth}%
%      \setlength{\leftmargin}{0pt}% default is \labelwidth
%      \setlength{\labelsep}{\widthof{\ }}% Управляет длиной отступа после точки % default is \biblabelsep
%      \setlength{\itemsep}{\bibitemsep}% Управление дополнительным вертикальным разрывом между записями. \bibitemsep по умолчанию соответствует \itemsep списков в документе.
%      \setlength{\itemindent}{\bibhang}% Пользуемся тем, что \bibhang по умолчанию принимает значение \parindent (абзацного отступа), который переназначен в styles.tex
%      \addtolength{\itemindent}{\labelwidth}% Сдвигаем правее на величину номера с точкой
%      \addtolength{\itemindent}{\labelsep}% Сдвигаем ещё правее на отступ после точки
%      \setlength{\parsep}{\bibparsep}%
%     }%
%      \renewcommand*{\makelabel}[1]{\hss##1}%
%  }
%  {\endlist}
%  {\item}

%%% Макросы автоматического подсчёта количества авторских публикаций.
% Печатают невидимую (пустую) библиографию, считая количество источников.
% http://tex.stackexchange.com/a/66851/79756
%
\makeatletter
    \newtotcounter{citenum}
    \defbibenvironment{counter}
        {\setcounter{citenum}{0}\renewcommand{\blx@driver}[1]{}} % begin code: убирает весь выводимый текст
        {} % end code
        {\stepcounter{citenum}} % item code: cчитает "печатаемые в библиографию" источники

    \newtotcounter{citeauthorvak}
    \defbibenvironment{countauthorvak}
        {\setcounter{citeauthorvak}{0}\renewcommand{\blx@driver}[1]{}}
        {}
        {\stepcounter{citeauthorvak}}

    \newtotcounter{citeauthorscopus}
    \defbibenvironment{countauthorscopus}
        {\setcounter{citeauthorscopus}{0}\renewcommand{\blx@driver}[1]{}}
        {}
        {\stepcounter{citeauthorscopus}}

    \newtotcounter{citeauthorwos}
    \defbibenvironment{countauthorwos}
        {\setcounter{citeauthorwos}{0}\renewcommand{\blx@driver}[1]{}}
        {}
        {\stepcounter{citeauthorwos}}

    \newtotcounter{citeauthorother}
    \defbibenvironment{countauthorother}
        {\setcounter{citeauthorother}{0}\renewcommand{\blx@driver}[1]{}}
        {}
        {\stepcounter{citeauthorother}}

    \newtotcounter{citeauthorconf}
    \defbibenvironment{countauthorconf}
        {\setcounter{citeauthorconf}{0}\renewcommand{\blx@driver}[1]{}}
        {}
        {\stepcounter{citeauthorconf}}

    \newtotcounter{citeauthor}
    \defbibenvironment{countauthor}
        {\setcounter{citeauthor}{0}\renewcommand{\blx@driver}[1]{}}
        {}
        {\stepcounter{citeauthor}}

    \newtotcounter{citeauthorvakscopuswos}
    \defbibenvironment{countauthorvakscopuswos}
        {\setcounter{citeauthorvakscopuswos}{0}\renewcommand{\blx@driver}[1]{}}
        {}
        {\stepcounter{citeauthorvakscopuswos}}

    \newtotcounter{citeauthorscopuswos}
    \defbibenvironment{countauthorscopuswos}
        {\setcounter{citeauthorscopuswos}{0}\renewcommand{\blx@driver}[1]{}}
        {}
        {\stepcounter{citeauthorscopuswos}}

    \newtotcounter{citeregistered}
    \defbibenvironment{countregistered}
        {\setcounter{citeregistered}{0}\renewcommand{\blx@driver}[1]{}}
        {}
        {\stepcounter{citeregistered}}

    \newtotcounter{citeauthorpatent}
    \defbibenvironment{countauthorpatent}
        {\setcounter{citeauthorpatent}{0}\renewcommand{\blx@driver}[1]{}}
        {}
        {\stepcounter{citeauthorpatent}}

    \newtotcounter{citeauthorprogram}
    \defbibenvironment{countauthorprogram}
        {\setcounter{citeauthorprogram}{0}\renewcommand{\blx@driver}[1]{}}
        {}
        {\stepcounter{citeauthorprogram}}

    \newtotcounter{citeexternal}
    \defbibenvironment{countexternal}
        {\setcounter{citeexternal}{0}\renewcommand{\blx@driver}[1]{}}
        {}
        {\stepcounter{citeexternal}}
\makeatother

\defbibheading{nobibheading}{} % пустой заголовок, для подсчёта публикаций с помощью невидимой библиографии
\defbibheading{pubgroup}{\section*{#1}} % обычный стиль, заголовок-секция
\defbibheading{pubsubgroup}{\noindent\textbf{#1}} % для подразделов "по типу источника"

%%%Сортировка списка литературы Русский-Английский (предварительно удалить dissertation.bbl) (начало)
%%%Источник: https://github.com/odomanov/biblatex-gost/wiki/%D0%9A%D0%B0%D0%BA-%D1%81%D0%B4%D0%B5%D0%BB%D0%B0%D1%82%D1%8C,-%D1%87%D1%82%D0%BE%D0%B1%D1%8B-%D1%80%D1%83%D1%81%D1%81%D0%BA%D0%BE%D1%8F%D0%B7%D1%8B%D1%87%D0%BD%D1%8B%D0%B5-%D0%B8%D1%81%D1%82%D0%BE%D1%87%D0%BD%D0%B8%D0%BA%D0%B8-%D0%BF%D1%80%D0%B5%D0%B4%D1%88%D0%B5%D1%81%D1%82%D0%B2%D0%BE%D0%B2%D0%B0%D0%BB%D0%B8-%D0%BE%D1%81%D1%82%D0%B0%D0%BB%D1%8C%D0%BD%D1%8B%D0%BC
%\DeclareSourcemap{
%    \maps[datatype=bibtex]{
%        \map{
%            \step[fieldset=langid, fieldvalue={tempruorder}]
%        }
%        \map[overwrite]{
%            \step[fieldsource=langid, match=russian, final]
%            \step[fieldsource=presort,
%            match=\regexp{(.+)},
%            replace=\regexp{aa$1}]
%        }
%        \map{
%            \step[fieldsource=langid, match=russian, final]
%            \step[fieldset=presort, fieldvalue={az}]
%        }
%        \map[overwrite]{
%            \step[fieldsource=langid, notmatch=russian, final]
%            \step[fieldsource=presort,
%            match=\regexp{(.+)},
%            replace=\regexp{za$1}]
%        }
%        \map{
%            \step[fieldsource=langid, notmatch=russian, final]
%            \step[fieldset=presort, fieldvalue={zz}]
%        }
%        \map{
%            \step[fieldsource=langid, match={tempruorder}, final]
%            \step[fieldset=langid, null]
%        }
%    }
%}
%Сортировка списка литературы (конец)

%%% Создание команд для вывода списка литературы %%%
\newcommand*{\insertbibliofull}{
    \printbibliography[keyword=bibliofull,section=0,title=\bibtitlefull]
    \ifnumequal{\value{draft}}{0}{
      \printbibliography[heading=nobibheading,env=counter,keyword=bibliofull,section=0]
    }{}
}
\newcommand*{\insertbiblioauthor}{
    \printbibliography[heading=pubgroup, section=0, filter=papersregistered, title=\bibtitleauthor]
}
% Вариант вывода печатных работ автора, с группировкой по типу источника.
\newcommand*{\insertbiblioauthorgrouped}{
    \section*{\bibtitleauthor}
    \printbibliography[heading=pubsubgroup, section=0, keyword=biblioauthorvak,    title=\bibtitleauthorvak,resetnumbers=false] % Работы автора из списка ВАК (сквозная нумерация)
    \printbibliography[heading=pubsubgroup, section=0, keyword=biblioauthorwos,    title=\bibtitleauthorwos,resetnumbers=false]% Работы автора, индексируемые Web of Science
    \printbibliography[heading=pubsubgroup, section=0, keyword=biblioauthorscopus, title=\bibtitleauthorscopus,resetnumbers=false]% Работы автора, индексируемые Scopus
    \printbibliography[heading=pubsubgroup, section=0, keyword=biblioauthorpatent, title=\bibtitleauthorpatent,resetnumbers=false]% Патенты
    \printbibliography[heading=pubsubgroup, section=0, keyword=biblioauthorprogram,title=\bibtitleauthorprogram,resetnumbers=false]% Программы для ЭВМ
    \printbibliography[heading=pubsubgroup, section=0, keyword=biblioauthorconf,   title=\bibtitleauthorconf,resetnumbers=false]% Тезисы конференций
    \printbibliography[heading=pubsubgroup, section=0, keyword=biblioauthorother,  title=\bibtitleauthorother,resetnumbers=false]% Прочие работы автора
}

\newcommand*{\insertbiblioexternal}{
    \printbibliography[heading=pubgroup,    section=0, keyword=biblioexternal,     title=\bibtitlefull]
}
    % BibLaTeX и Biber
}

%%% Диагностика сборки
\typeout{Selected options:}
\typeout{Draft mode: \arabic{draft}}
\typeout{Font: \arabic{fontfamily}}
\typeout{AltFont: \arabic{usealtfont}}
\typeout{Bibliography backend: \arabic{bibliosel}}
\typeout{Precompile images: \arabic{imgprecompile}}
\listfiles  % Версии пакетов в журнале сборки

\begin{document}

%%% Русские названия разделов и подписей для Babel и Polyglossia
% https://tex.stackexchange.com/a/156050
\gappto\captionsrussian{%%% Переопределение именований %%%
\renewcommand{\contentsname}{Оглавление}% (ГОСТ Р 7.0.11-2011, 4)
\renewcommand{\figurename}{Рисунок}% (ГОСТ Р 7.0.11-2011, 5.3.9)
\renewcommand{\tablename}{Таблица}% (ГОСТ Р 7.0.11-2011, 5.3.10)
\renewcommand{\listfigurename}{Список рисунков}%
\renewcommand{\listtablename}{Список таблиц}%
\renewcommand{\bibname}{\bibtitlefull}%}
%%% Переопределение именований %%%
\renewcommand{\contentsname}{Оглавление}% (ГОСТ Р 7.0.11-2011, 4)
\renewcommand{\figurename}{Рисунок}% (ГОСТ Р 7.0.11-2011, 5.3.9)
\renewcommand{\tablename}{Таблица}% (ГОСТ Р 7.0.11-2011, 5.3.10)
\renewcommand{\listfigurename}{Список рисунков}%
\renewcommand{\listtablename}{Список таблиц}%
\renewcommand{\bibname}{\bibtitlefull}%
\gappto\captionsrussian{% Переопределения названий для nomencl. Так как опция russian не для utf8
\renewcommand{\nomname}{Список сокращений и условных обозначений}%
\renewcommand{\eqdeclaration}[1]{, см.~(#1)}%
\renewcommand{\pagedeclaration}[1]{, стр.~#1}%
\renewcommand{\nomAname}{Латинские буквы}%
\renewcommand{\nomGname}{Греческие буквы}%
\renewcommand{\nomXname}{Верхние индексы}%
\renewcommand{\nomZname}{Индексы}%}
% Переопределения названий для nomencl. Так как опция russian не для utf8
\renewcommand{\nomname}{Список сокращений и условных обозначений}%
\renewcommand{\eqdeclaration}[1]{, см.~(#1)}%
\renewcommand{\pagedeclaration}[1]{, стр.~#1}%
\renewcommand{\nomAname}{Латинские буквы}%
\renewcommand{\nomGname}{Греческие буквы}%
\renewcommand{\nomXname}{Верхние индексы}%
\renewcommand{\nomZname}{Индексы}%

%%% Титул и оглавление
% Титульный лист (ГОСТ Р 7.0.11-2001, 5.1)
\thispagestyle{empty}
\hypersetup{
  pdfauthor={\thesisAuthor},
  pdftitle={\thesisTitle},
  pdfsubject={\thesisSpecialtyTitle. \thesisDegreeShort\ \thesisDegree},
}
\begin{center}
{\small \thesisOrganization\par}
\end{center}
%
\vspace{0pt plus4fill} %число перед fill = кратность относительно некоторого расстояния fill, кусками которого заполнены пустые места
\begingroup
\newcommand{\kiselevTitleNotice}{%
  \textbf{Открытый проект}\\
  Доступно бесплатно\\[0.5\baselineskip]
  Современная редакция\\
  и новые рисунки\\
  по лицензии \href{https://creativecommons.org/licenses/by/4.0/}{CC~BY~4.0}%
}
\IfFileExists{Book/images/rockyversus-github.jpg}{
  \begin{minipage}[b]{0.5\linewidth}
    \begin{flushleft}
      \href{https://github.com/rockyversus}{%
        \includegraphics[height=3.5cm]{Book/images/rockyversus-github.jpg}%
      }
    \end{flushleft}
  \end{minipage}%
  \begin{minipage}[b]{0.5\linewidth}
    \begin{flushright}
      \kiselevTitleNotice
    \end{flushright}
  \end{minipage}
}{
  \begin{flushright}
    \kiselevTitleNotice
  \end{flushright}
}
\endgroup
%
\vspace{0pt plus6fill} %число перед fill = кратность относительно некоторого расстояния fill, кусками которого заполнены пустые места
\begin{center}
{\Large\bfseries \thesisAuthor\par}
\end{center}
%
\vspace{0pt plus1fill} %число перед fill = кратность относительно некоторого расстояния fill, кусками которого заполнены пустые места
\begin{center}
\begingroup
  \renewcommand{\thesisTitleFormat}[2]{%
    \scalebox{1.65}{\large\MakeUppercase{#1}}\par
    \vspace{0.45\baselineskip}%
    \scalebox{3.4}{\large\bfseries\MakeUppercase{#2}}\par
  }
  \thesisTitle
\endgroup

\vspace{0pt plus2fill} %число перед fill = кратность относительно некоторого расстояния fill, кусками которого заполнены пустые места
\ifdefempty{\thesisSpecialtyNumber}
  {%
    {\small Раздел\par}
    {\Large\bfseries \MakeUppercase{\thesisSpecialtyTitle}\par}
  }
  {{\large Специальность \thesisSpecialtyNumber\ "---\par <<\thesisSpecialtyTitle>>\par}}
\ifdefined\thesisSpecialtyTwoNumber
  {\large Специальность \thesisSpecialtyTwoNumber\ "---\par <<\thesisSpecialtyTwoTitle>>\par}
\fi

\vspace{0pt plus2fill} %число перед fill = кратность относительно некоторого расстояния fill, кусками которого заполнены пустые места
{\small Подраздел\par}
{\large \MakeUppercase{\thesisDegreeShort}\par}
\smallskip
{\LARGE \MakeUppercase{\thesisDegree}\par}
\end{center}
%
\vspace{0pt plus4fill} %число перед fill = кратность относительно некоторого расстояния fill, кусками которого заполнены пустые места
\begin{flushright}
Переработка на современный лад\\
\ifdefined\supervisorDead
  \framebox{\supervisorFio}%
\else
  \supervisorFio
\fi\\
\supervisorRegalia
\ifdefined\supervisorTwoFio
  \par
  \ifdefined\supervisorTwoDead
    \framebox{\supervisorTwoFio}%
  \else
    \supervisorTwoFio
  \fi\\
  \supervisorTwoRegalia
\fi
\end{flushright}
%
\vspace{0pt plus4fill} %число перед fill = кратность относительно некоторого расстояния fill, кусками которого заполнены пустые места
{\centering\thesisCity\ "--- \thesisYear\par}

% Оглавление (ГОСТ Р 7.0.11-2011, 5.2)
\ifdefmacro{\microtypesetup}{\microtypesetup{protrusion=false}}{} % не рекомендуется применять пакет микротипографики к автоматически генерируемому оглавлению
\settowidth{\cftchapternumwidth}{\cftchapterfont III\cftchapteraftersnum\enspace}
\tableofcontents*
\addtocontents{toc}{\protect\tocheader}
\endTOCtrue
\ifdefmacro{\microtypesetup}{\microtypesetup{protrusion=true}}{}


%%% Введение
% Здесь номера рисунков и таблиц не зависят от номера главы.
\ifnumequal{\value{contnumfig}}{1}{}{%
    \counterwithout{figure}{chapter}%
}
\ifnumequal{\value{contnumtab}}{1}{}{%
    \counterwithout{table}{chapter}%
}
% А. П. Киселёв, «Элементарная геометрия», 23-е издание, 1914. Печатные с. 1–8.
\chapter*{Введение}
\addcontentsline{toc}{chapter}{Введение}

\section*{Математические предложения}
\addcontentsline{toc}{section}{Математические предложения}

% Печатная с. 1.
\kiselevparagraph{para:intro-mathematical-statements} В любой математической науке могут встретиться следующие предложения:

\textbf{Определения.} Так называют предложения, в которых разъясняется, какой смысл придают тому или другому выражению или названию. Например, в арифметике мы встречаем определения наименьшего кратного, общего наибольшего делителя и т.~п.

\textbf{Аксиомы.} Так называют истины, которые, вследствие своей очевидности, принимаются без доказательства. Таковы, например, предложения:

Если каждая из двух величин равна одной и той же третьей величине, то эти две величины равны между собой.

Если к равным величинам прибавим одну и ту же величину или вычтем из них одну и ту же величину, то равенство не нарушится.

Если к неравным величинам прибавим одну и ту же величину или вычтем из них одну и ту же величину, то направление неравенства не изменится, т.~е. большая величина останется большей.

\textbf{Теоремы.} Так называются предложения, истинность которых устанавливается только рассуждением (доказательством). Примером может служить арифметическая истина: «если сумма цифр делится на~9, то число делится на~9».

\textbf{Следствия.} Так называются предложения, которые составляют непосредственный вывод из аксиомы или теоремы. Например, из теоремы: «в пропорции произведение крайних членов равно произведению средних», выводится следствие: «крайний член пропорции равен произведению средних членов, делённому на другой крайний».

% Печатная с. 2.
\kiselevparagraph{para:intro-theorem-structure} \textbf{Состав теоремы.} В любой теореме можно выделить две части: условие и заключение. \emph{Условие} выражает то, что предполагается данным; \emph{заключение} --- то, что требуется доказать. Например, в теореме: «если сумма цифр делится на~9, то число делится на~9», условием служит первая часть теоремы: «\emph{если} сумма цифр делится на~9», а заключением --- вторая часть: «\emph{то} число делится на~9»; другими словами, нам дано, что сумма цифр некоторого числа делится на~9, а требуется доказать, что в таком случае и само число делится на~9.

Условие и заключение теоремы могут иногда состоять из нескольких отдельных условий и заключений; например, в теореме: «\emph{если число делится на~2 и на~3}, то оно разделится на~6», условие состоит из двух частей: если число делится на~2 и если число делится на~3.

Любую теорему можно сформулировать так, чтобы её условие начиналось со слова «если», а заключение --- со слова «то».

\kiselevparagraph{para:intro-converse-theorem} \textbf{Обратная теорема.} Теоремой, обратной данной теореме, называется такая, в которой условием служит заключение или часть заключения данной теоремы, а заключением --- условие или часть условия данной теоремы. Например, следующие две теоремы обратны друг другу:

\begin{center}
    \begin{tabular}{@{}p{0.44\linewidth}|p{0.44\linewidth}@{}}
        \emph{Если} сумма цифр делится на~9,\par \emph{то} число делится на~9. &
        \emph{Если} число делится на~9,\par \emph{то} сумма цифр делится на~9.
    \end{tabular}
\end{center}

Если одну из этих теорем назовём \emph{прямой}, то другую следует назвать \emph{обратной}.

В этом примере обе теоремы, и прямая и обратная, оказываются верными. Но так бывает не всегда. Например, теорема: «если каждое слагаемое делится на одно и то же число, то и сумма разделится на то же число» --- верна, но неверно обратное предложение: «если сумма делится на какое-нибудь число, то каждое слагаемое разделится на него».

\kiselevparagraph{para:intro-opposite-theorem} \textbf{Противоположная теорема.} Теоремой, противоположной данной теореме, называется такая, условие и заключение которой представляют \emph{отрицание} условия и заключения
% Печатная с. 3.
данной теоремы. Например, теореме: «если сумма цифр делится на~9, то число делится на~9» соответствует такая противоположная: «если сумма цифр \textbf{не} делится на~9, то число \textbf{не} делится на~9».

Здесь также следует заметить, что верность прямой теоремы ещё не служит признаком верности противоположной; например, противоположное предложение: «если каждое слагаемое \emph{не} делится на одно и то же число, то и сумма \emph{не} разделится на это число», --- неверно, тогда как прямое предложение верно.

\kiselevparagraph{para:intro-relations-between-theorems} \textbf{Зависимость между теоремами: прямой, обратной и противоположной.} Чтобы лучше понять эту зависимость, запишем теоремы в сокращённом виде:

\begin{enumerate}[label=\arabic*),ref=\arabic*]
    \item\label{item:direct-theorem} \emph{Прямая теорема:} если есть \(A\), то есть и \(B\).
    \item\label{item:converse-theorem} \emph{Обратная теорема:} если есть \(B\), то есть и \(A\).
    \item\label{item:opposite-direct-theorem} \emph{Противоположная прямой:} если нет \(A\), то нет и \(B\).
    \item\label{item:opposite-converse-theorem} \emph{Противоположная обратной:} если нет \(B\), то нет и \(A\).
\end{enumerate}

Легко заметить, что \hyperref[item:direct-theorem]{первое} и \hyperref[item:opposite-converse-theorem]{четвёртое} предложения \emph{следуют друг из друга}, как и \hyperref[item:converse-theorem]{второе} и \hyperref[item:opposite-direct-theorem]{третье}. Действительно, из предложения «если есть \(A\), то есть и \(B\)» непосредственно следует «если нет \(B\), то нет и \(A\)» (так как если бы \(A\) было, то, согласно \hyperref[item:direct-theorem]{первому предложению}, было бы и \(B\)). И наоборот, из предложения «если нет \(B\), то нет и \(A\)» выводим «если есть \(A\), то есть и \(B\)» (так как если бы \(B\) не было, то не было бы и \(A\)). Точно так же можно убедиться, что из \hyperref[item:converse-theorem]{второго предложения} следует \hyperref[item:opposite-direct-theorem]{третье}, и наоборот.

Поэтому, чтобы быть уверенным в верности всех четырёх теорем, достаточно доказать две из них: прямой и обратной или прямой и противоположной.

\section*{Прямая линия, плоскость. Понятие о геометрии}
\addcontentsline{toc}{section}{Прямая линия, плоскость. Понятие о геометрии}

\kiselevparagraph{para:intro-geometric-figures} \textbf{Геометрические фигуры.} Любая ограниченная часть пространства называется \emph{геометрическим телом}.

Геометрическое тело можно подразделять \emph{на части}; каждая часть геометрического тела также является геометрическим телом.

Граница геометрического тела, т.~е. то, чем оно отделяется от остального пространства, называется \emph{поверхностью}.

% Печатная с. 4.
Поверхность можно подразделять \emph{на части}; каждая часть поверхности также является поверхностью.

Граница поверхности называется \emph{линией}.

Линию можно также подразделять \emph{на части}; каждая часть линии также является линией.

Граница линии называется \emph{точкой}.

Геометрическое тело, поверхность, линия и точка не существуют по отдельности. Однако мысленно можно рассматривать поверхность независимо от геометрического тела, линию --- независимо от поверхности и точку --- независимо от линии. При этом поверхность мы должны представлять себе не имеющей толщины, линию --- не имеющей ни толщины, ни ширины, и точку --- не имеющей ни длины, ни ширины, ни толщины.

Любая линия содержит в себе бесчисленное множество точек. Принято говорить, что эти точки \emph{лежат} на линии, или что эта линия \emph{проходит} через эти точки. Их можно рассматривать как последовательные положения одной и той же точки, движущейся вдоль этой линии. Поэтому можно сказать, что \emph{линия --- след движения точки}. Если, например, мы ведём остриём карандаша по бумаге, получившийся след можно считать линией лишь приблизительно. Остриё карандаша не является геометрической точкой, поэтому след на бумаге имеет некоторую ширину (и даже толщину). Чем острее заточен карандаш, тем ближе его остриё к геометрической точке, а проведённая им линия --- к геометрической линии.

Аналогично поверхность можно рассматривать как \emph{след движения линии}, двигающейся в пространстве некоторым образом.

Совокупность каких бы то ни было точек, линий, поверхностей или тел, расположенных известным образом в пространстве, называется вообще геометрической \emph{фигурой}.

\kiselevparagraph{para:intro-geometry} \textbf{Геометрия.} Наука, рассматривающая свойства геометрических фигур, называется \emph{геометрией}, что в переводе с
% Печатная с. 5.
греческого языка означает \emph{землемерие}. Эта наука получила такое название потому, что в древности главной целью геометрии было измерение расстояний и площадей на земной поверхности.

\kiselevparagraph{para:intro-space-axiom} В начале изучения геометрии следует указать следующее общее свойство фигур:

\textbf{Аксиома пространства. Любую геометрическую фигуру можно перенести из одного места пространства в другое, не нарушая ни величины составляющих фигуру частей, ни их взаимного расположения.}

\kiselevparagraph{para:intro-straight-line} \textbf{Прямая линия.} Каждый знает, что такое \emph{прямая линия}, или просто \emph{прямая}, представление о которой нам даёт туго натянутая нить. \emph{Понятие о прямой элементарно}, т.~е. оно не может быть определено через другие, более простые понятия.

На рисунке прямую изображают в виде тонкой черты, проведённой от руки или с помощью чертёжной \emph{линейки}.

Прямая линия обладает следующими очевидными свойствами:

\textbf{Аксиомы прямой.}

\begin{enumerate}[label=\arabic*),ref=\arabic*]
    \item\label{item:line-through-two-points} \textbf{Через любые две точки пространства можно провести ровно одну прямую.}
    \item \textbf{Прямую можно продолжать без конца в обе стороны от каждой её точки.}
\end{enumerate}

Из \hyperref[item:line-through-two-points]{первой аксиомы} следует:

Если две прямые наложены одна на другую так, что две точки одной прямой совпадают с двумя точками другой прямой, то эти прямые \emph{совпадают} и во всех остальных точках (иначе через две точки можно было бы провести две различные прямые, что противоречит \hyperref[item:line-through-two-points]{первой аксиоме}).

По той же причине две прямые могут пересечься только \emph{в одной точке}.

\kiselevparagraph{para:intro-finite-and-infinite-lines} \textbf{Прямая конечная и бесконечная.} Если прямую мысленно продолжают без конца в обе стороны, её называют \emph{бесконечной} или \emph{неограниченной} прямой. Конечно, такую прямую изобразить на рисунке
% Печатная с. 6.
невозможно. Изображают только часть прямой и представляют, что она продолжается в обе стороны без конца. Прямую обычно обозначают двумя буквами, соответствующими любым двум её точкам. Так говорят «прямая \(AB\) или \(BA\)» (рис.~\ref{fig:kiselev-intro-line}).

\begin{figure}[htbp]
    \centering
    \captionsetup{name=Рис.,labelsep=period}
    \begin{minipage}[t]{0.46\linewidth}
        \centering
        \captionsetup{type=figure}
        \begin{tikzpicture}
    \draw[thick] (0,0) node[left] {\(A\)} -- (3.3,0) node[right] {\(B\)};
\end{tikzpicture}

        \caption{Прямая \(AB\) или \(BA\)}
        \label{fig:kiselev-intro-line}
    \end{minipage}\hfill
    \begin{minipage}[t]{0.46\linewidth}
        \centering
        \captionsetup{type=figure}
        \begin{tikzpicture}
    \draw[thick] (0,0) node[left] {\(C\)} -- (2.4,0) node[right] {\(D\)};
    \draw[thick] (0,-0.09) -- (0,0.09) (2.4,-0.09) -- (2.4,0.09);
\end{tikzpicture}

        \caption{Отрезок \(CD\)}
        \label{fig:kiselev-intro-segment}
    \end{minipage}
\end{figure}

Часть прямой, ограниченная с обеих сторон, называется \emph{отрезком} прямой или \emph{конечной} прямой; такой отрезок обозначают двумя буквами у его концов (отрезок \(CD\), рис.~\ref{fig:kiselev-intro-segment}). Отрезок прямой, соединяющий две точки, называется иногда \emph{расстоянием} между ними.

Иногда рассматривают прямую, ограниченную только с одной стороны, например, в точке \(A\) (рис.~\ref{fig:kiselev-intro-ray}). О такой прямой говорят, что она \emph{исходит} из точки \(A\); её называют \emph{полупрямой} (или \emph{лучом}).

\begin{figure}[htbp]
    \centering
    \captionsetup{name=Рис.,labelsep=period}
    \begin{tikzpicture}
    \draw[thick] (0,0) node[left] {\(A\)} -- (4.2,0);
    \draw[thick] (0,-0.09) -- (0,0.09);
    \draw[thick,densely dashed] (4.2,0) -- (6,0);
\end{tikzpicture}

    \caption{Полупрямая (или луч)}
    \label{fig:kiselev-intro-ray}
\end{figure}

\kiselevparagraph{para:intro-comparison-of-segments} \textbf{Равенство и неравенство конечных прямых. Два отрезка прямой считаются равными, если они могут быть наложены друг на друга так, что совмещаются.} Допустим, что мы накладываем отрезок \(AB\) на отрезок \(CD\) (рис.~\ref{fig:kiselev-intro-sum}) так, чтобы точка \(A\) совпала с точкой \(C\), а отрезок \(AB\) был направлен от \(C\) к \(D\); если при этом концы \(B\) и \(D\) совпадут, то отрезки \(AB\) и \(CD\) считаются равными; в противном случае отрезки будут неравны, причём меньшим считается тот, который составляет часть другого.

Чтобы отложить на прямой отрезок, равный данному, пользуются \emph{циркулем} --- прибором, знакомым учащимся по опыту.

\kiselevparagraph{para:intro-sum-of-segments} \textbf{Сумма конечных прямых. Суммой нескольких данных отрезков прямой называется отрезок прямой, составленный из частей, равных соответственно данным отрезкам.} Допустим, нужно найти сумму трёх отрезков: \(AB\), \(CD\) и \(EF\) (рис.~\ref{fig:kiselev-intro-sum}). Для
% Печатная с. 7.
этого на какой-нибудь прямой берём произвольную точку \(M\) и откладываем от неё часть \(MN\), равную \(AB\); затем от точки \(N\) в том же направлении откладываем часть \(NP\), равную \(CD\), и часть \(PQ\), равную \(EF\). Отрезок \(MQ\) будет суммой данных отрезков \(AB\), \(CD\) и \(EF\), которые по отношению к этой сумме называются слагаемыми. Подобным образом можно получить сумму какого угодно числа отрезков.

\begin{figure}[htbp]
    \centering
    \captionsetup{name=Рис.,labelsep=period}
    \begin{tikzpicture}
    \foreach \leftlabel/\rightlabel/\first/\last in {A/B/0/3.3,C/D/4.9/7.9,E/F/9.5/11.3} {
        \draw[thick] (\first,1.2) node[left] {\(\leftlabel\)}
            -- (\last,1.2) node[right] {\(\rightlabel\)};
        \draw[thick] (\first,1.11) -- (\first,1.29) (\last,1.11) -- (\last,1.29);
    }
    \draw[thick] (0.2,0) -- (8.3,0);
    \draw[thick,densely dashed] (-0.9,0) -- (0.2,0) (8.3,0) -- (9.5,0);
    \foreach \position/\point in {0.2/M,3.5/N,6.5/P,8.3/Q} {
        \draw[thick] (\position,-0.09) -- (\position,0.09) node[above] {\(\point\)};
    }
\end{tikzpicture}

    \caption{Сумма конечных прямых}
    \label{fig:kiselev-intro-sum}
\end{figure}

Сумма отрезков прямой обладает теми же свойствами, что и любая сумма. Она не зависит от порядка слагаемых (\emph{переместительное свойство}) и не изменяется, если некоторые слагаемые заменить их суммой (\emph{сочетательное свойство}). Например, легко убедиться (рис.~\ref{fig:kiselev-intro-sum}), что
\[
    \begin{gathered}
        AB+CD+EF=CD+EF+AB=EF+AB+CD=\dots \\
        \text{и}\quad AB+CD+EF=AB+(CD+EF).
    \end{gathered}
\]

Из понятия о сумме выводятся понятия о разности, произведении и частном отрезков. Так, \emph{разность} отрезков \(AB\) и \(CD\) (если \(AB>CD\)) --- это отрезок, который в сумме с \(CD\) даёт \(AB\); \emph{произведение} отрезка \(AB\) на число~3 равно сумме трёх отрезков, из которых каждый равен \(AB\), \emph{частное} при делении отрезка \(AB\) на число~3 равно третьей части \(AB\) и т.~п.

Мы считаем очевидным, что любой отрезок прямой можно разделить (хотя бы только мысленно) на~2, на~3, на~4 и т.~д. равные части.

\kiselevparagraph{para:intro-plane} \textbf{Плоскость. Плоскостью называется поверхность, обладающая тем свойством, что прямая, проходящая через любые две точки этой поверхности, целиком лежит на этой поверхности.} Допустим, мы хотим проверить, является ли поверхность стола плоскостью. Для этого берём хорошо выверенную линейку и прикладываем её краем в различных направлениях к поверхности стола так, чтобы какие-нибудь
% Печатная с. 8.
две точки линейки лежали на этой поверхности. Если при этом окажется, что, в каком бы направлении мы ни приложили линейку, все остальные её точки будут лежать на поверхности стола, то эта поверхность является плоскостью.

Существование плоскости в пространстве принимается за аксиому.

Укажем следующее свойство плоскости, которое мы примем здесь без доказательства:

\textbf{Любую часть плоскости можно наложить всеми её точками на другое место этой или другой плоскости, причём накладываемую часть можно предварительно перевернуть на другую сторону.}

\kiselevparagraph{para:intro-branches-of-geometry} \textbf{Разделение геометрии.} Геометрия делится на две части: \emph{планиметрия} и \emph{стереометрия}. Первая рассматривает свойства фигур, все части которых лежат в одной плоскости. Вторая рассматривает свойства фигур, не все части которых лежат в одной плоскости.


%%% Главы I–V
% В главах нумерация рисунков и таблиц следует настройкам Book/setup.tex.
\ifnumequal{\value{contnumfig}}{1}{%
    \counterwithout{figure}{chapter}%
}{%
    \counterwithin{figure}{chapter}%
}
\ifnumequal{\value{contnumtab}}{1}{%
    \counterwithout{table}{chapter}%
}{%
    \counterwithin{table}{chapter}%
}

\begingroup
    \chapterstyle{kiselev}
    \renewcommand*{\thechapter}{\Roman{chapter}}
    \chapter{Определение положения плоскости}\label{ch:ch1}
\setcounter{kiselevparagraph}{343} % Первый пункт главы в издании 1914 года — 344.

% А. П. Киселёв, «Элементарная геометрия», 23-е издание, 1914. Печатные с. 271–275.
% Печатная с. 271.
\kiselevparagraph{para:plane-definition} \textbf{Определение.} Плоскостью (\hyperref[para:intro-plane]{см.~Введение, пункт~\ref*{para:intro-plane}}) называется поверхность, обладающая тем свойством, что прямая, проходящая через две произвольные точки этой поверхности, \emph{целиком лежит на этой поверхности}. Возможность существования такой поверхности принимается за аксиому.

\emph{Стереометрией} (\hyperref[para:intro-branches-of-geometry]{см.~Введение, пункт~\ref*{para:intro-branches-of-geometry}}) называется часть геометрии, в которой рассматриваются свойства фигур, не все части которых лежат в одной плоскости.

\kiselevparagraph{para:plane-axioms} \textbf{Аксиомы плоскости.} Укажем следующие свойства плоскости, которые мы принимаем за очевидные:

\begin{enumerate}[label=\arabic*),ref=\arabic*]
    \item Плоскость --- \emph{незамкнутая} поверхность.
    \item Любая плоскость делит пространство на две части, расположенные по разные стороны от неё.
    \item\label{item:line-crosses-plane} Прямая, имеющая с плоскостью только одну общую точку, \emph{пересекает} плоскость, т.~е. из части пространства, лежащей по одну сторону от плоскости, переходит в часть пространства, лежащую по другую её сторону; поэтому две полупрямые, на которые плоскость делит эту прямую, лежат по разные стороны от неё.
    % Печатная с. 272.
    \item\label{item:segment-crosses-plane} Отрезок прямой, соединяющий две точки пространства, расположенные по разные стороны от плоскости, \emph{пересекает} эту плоскость, тогда как отрезок прямой, соединяющий две точки, расположенные по одну сторону плоскости, не пересекает её.
    \item Через любую прямую можно провести \emph{бесчисленное множество} плоскостей.
    \item Любая прямая на плоскости делит её на две части (называемые \emph{полуплоскостями}).
    \item Плоскость можно вращать вокруг любой прямой, лежащей на ней, причём каждую из частей, на которые плоскость делится этой прямой, можно заставить пройти через любую точку пространства.
\end{enumerate}

\kiselevparagraph{para:plane-representation} \textbf{Изображение и обозначение плоскости.} На рисунке изображают часть плоскости, обычно в форме параллелограмма или прямоугольника (рис.~\ref{fig:kiselev-stereo-301}). Плоскость обычно обозначают одной или двумя буквами, например: плоскость \(P\), плоскость \(MN\).

\begin{figure}[htbp]
    \centering
    \renewcommand{\thefigure}{301}
    \captionsetup{name=Рис.,labelsep=period}
    \begin{tikzpicture}[line join=round]
    \coordinate (front left) at (0,0);
    \coordinate (front right) at (3.6,0);
    \coordinate (back left) at (1.5,1.65);
    \coordinate (back right) at (5.1,1.65);

    \draw[thick] (front left) -- (back left) -- (back right);
    \draw[very thick] (front left) -- (front right) -- (back right);
    \path (front left) -- (back left) node[midway,left] {\(M\)};
    \path (front right) -- (back right) node[midway,right] {\(N\)};
    \node at (2.55,0.825) {\(P\)};
\end{tikzpicture}

    \caption{Изображение и обозначение плоскости}\label{fig:kiselev-stereo-301}
\end{figure}

\kiselevparagraph{para:plane-through-three-points} \textbf{Теорема. Через любые три точки} (\(A\), \(B\) и \(C\), рис.~\ref{fig:kiselev-stereo-302}), \textbf{не лежащие на одной прямой, можно провести ровно одну плоскость.}

\begin{figure}[htbp]
    \centering
    \renewcommand{\thefigure}{302}
    \captionsetup{name=Рис.,labelsep=period}
    \begin{tikzpicture}[scale=1.5,line cap=round,every node/.style={inner sep=3.5pt}]
    \coordinate (A) at (0,0);
    \coordinate (B) at (4,0);
    \coordinate (C) at (3.2,-1.8);
    \coordinate (D) at (1.74,-2.052);
    \coordinate (E) at (2.12,0);
    \coordinate (F) at (1.92,-1.08);
    \coordinate (M) at (2.24,0.648);

    \draw[thick] (-1,0) -- (5,0);
    \draw[thick] (-1,0.5625) -- (4.28,-2.4075);
    \draw[thick] (4.4,0.9) -- (2.9,-2.475);
    \draw[thick] (2.34,1.188) -- (D);

    \foreach \point/\placement in {A/below left,B/above left,E/above left,F/below left,M/right} {
        \filldraw[fill=white] (\point) circle[radius=1pt] node[\placement] {\(\point\)};
    }
    \filldraw[fill=white] (C) circle[radius=1pt] node[above right,xshift=5pt] {\(C\)};
    \node[below] at (D) {\(D\)};
\end{tikzpicture}

    \caption{Плоскость, которая проходит через три точки \(A\), \(B\) и \(C\)}\label{fig:kiselev-stereo-302}
\end{figure}

\begin{enumerate}[label=\arabic*)]
    \item Через две любые из трёх данных точек, например через \(A\) и \(B\), проведём прямую, а через неё --- произвольную плоскость. Будем вращать эту плоскость вокруг прямой \(AB\), пока она не пройдёт через точку \(C\). Получим плоскость, которая проходит через три точки \(A\), \(B\) и \(C\).

    Остаётся доказать, что такая плоскость может быть только одна.

    \item Предположим, что через те же три точки проведены две плоскости. Обозначим одну \(P\), а другую --- \(P_1\).
    % Печатная с. 273.
    Докажем, что эти плоскости сливаются в одну. Сначала заметим, что прямые \(AB\), \(BC\) и \(AC\), проходящие через каждую пару данных точек, принадлежат обеим плоскостям, так как эти прямые имеют по две общие точки и с плоскостью \(P\), и с плоскостью \(P_1\). Теперь возьмём на плоскости \(P\) произвольную точку \(M\) и докажем, что эта точка принадлежит и другой плоскости \(P_1\). Для этого на плоскости \(P\) через точку \(M\) проведём прямую \(MD\) так, чтобы она пересекала \(AB\) и \(AC\) в точках \(E\) и \(F\). Так как прямые \(AB\) и \(AC\) принадлежат и другой плоскости \(P_1\), то их точки \(E\) и \(F\) также принадлежат этой плоскости. Поэтому прямая \(MD\), проходящая через \(E\) и \(F\), целиком лежит в плоскости \(P_1\) (\hyperref[para:plane-definition]{по определению плоскости}), а значит, и её точка \(M\) лежит в этой плоскости. Таким образом, любая точка \(M\) плоскости \(P\) принадлежит и плоскости \(P_1\); значит, эти плоскости сливаются.
\end{enumerate}

\kiselevparagraph{para:plane-uniqueness-corollaries} \textbf{Следствия.}

\begin{enumerate}[label=\arabic*),ref=\arabic*]
    \item\label{item:plane-through-line-and-point} \textbf{Через прямую и точку вне её можно провести ровно одну плоскость}, потому что точка вне прямой и две любые точки этой прямой образуют три точки, через которые \hyperref[para:plane-through-three-points]{по доказанной теореме} можно провести ровно одну плоскость.
    \item\label{item:plane-through-intersecting-lines} \textbf{Через две пересекающиеся прямые можно провести ровно одну плоскость}, потому что, взяв точку пересечения и ещё по одной точке на каждой прямой, мы будем иметь три точки, через которые и т.~д.
    \item \textbf{Через две параллельные прямые можно провести ровно одну плоскость}, потому что параллельные прямые, по определению, лежат в одной плоскости; эта плоскость единственная, так как через одну из параллельных и какую-нибудь точку другой можно провести не более одной плоскости.
    \item \textbf{Любую часть плоскости можно наложить всеми её точками на другое место этой или другой плоскости, причём накладываемую часть можно предварительно перевернуть на другую сторону}, потому что всегда можно наложить одну плоскость на другую так, чтобы у них совпали какие-нибудь три точки, не лежащие на одной прямой, а тогда совпадут и остальные точки.
\end{enumerate}

% Печатная с. 274.
\kiselevparagraph{para:intersecting-planes} \textbf{Теорема. Если две плоскости имеют общую точку, то они имеют и общую прямую, которая проходит через эту точку и по которой плоскости пересекаются.}

Пусть плоскость \(P\) (рис.~\ref{fig:kiselev-stereo-303}) имеет точку \(A\), общую с другой плоскостью \(Q\). Проведём на плоскости \(Q\) через точку \(A\) какие-нибудь две прямые \(BB_1\) и \(CC_1\). Если одна из этих прямых лежит и в плоскости \(P\), то у этих плоскостей уже есть общая прямая. Предположим, что этого нет; тогда обе эти прямые \emph{пересекают} плоскость \(P\) (\hyperref[item:line-crosses-plane]{см.~§~\ref*{para:plane-axioms}, пункт~\ref*{item:line-crosses-plane}}); следовательно, плоскость \(P\) делит каждую из них на две полупрямые, расположенные по разные стороны от неё. Возьмём на полупрямых \(AC\) и \(AB\)\footnote{В тексте издания 1914 года указан луч \(AB\). На рис.~\ref{fig:kiselev-stereo-303} точка \(D_1\) лежит на луче \(AB_1\).} какие-нибудь точки \(D\) и \(D_1\) и проведём прямую \(DD_1\). Эта прямая проходит из части пространства над плоскостью \(P\) в часть пространства под ней и пересекает плоскость \(P\) в некоторой точке \(E\) (\hyperref[item:segment-crosses-plane]{см.~§~\ref*{para:plane-axioms}, пункт~\ref*{item:segment-crosses-plane}}). Поскольку прямая \(DD_1\) имеет с плоскостью \(Q\) две общие точки (\(D\) и \(D_1\)), она целиком лежит в этой плоскости. Поэтому точка \(E\) прямой \(DD_1\) также принадлежит плоскости \(Q\). Итак, плоскости \(P\) и \(Q\) имеют две общие точки \(A\) и \(E\); значит, они имеют и общую прямую \(AE\), проходящую через эти точки.

\begin{figure}[htbp]
    \centering
    \renewcommand{\thefigure}{303}
    \captionsetup{name=Рис.,labelsep=period}
    \begin{tikzpicture}[scale=1.3,line cap=round,line join=round]
    \coordinate (A) at (1.6,2.8);
    \coordinate (B) at (1,5.5);
    \coordinate (B1) at (2.2222,0);
    \coordinate (C) at (3.67,5.5);
    \coordinate (C1) at (0.22,1);
    \coordinate (D) at (3.256,4.96);
    \coordinate (D1) at (2.08,0.64);
    \coordinate (E) at (2.668,2.8);

    % Плоскость Q. Часть левого ребра скрыта плоскостью P.
    \draw[thick] (0,2.63) -- (0,5.9) -- (4.1,5.9) -- (4.1,0) -- (0,0) -- (0,1.45);
    \draw[thick,densely dashed] (0,2.63) -- (0,1.45);

    % Плоскость P. Дальняя часть контура скрыта плоскостью Q.
    \draw[very thick] (0,2.63) -- (-0.45,1.45) -- (3.58,1.45) -- (4.63,4.2) -- (4.1,4.2);
    \draw[thick,densely dashed] (0,2.63) -- (0.6,4.2) -- (4.1,4.2);

    \draw[thick] (B) -- (A) (C) -- (A) (D) -- (E);
    \draw[thick,densely dashed] (A) -- (1.9,1.45) (A) -- (0.565,1.45) (E) -- (2.3,1.45);
    \draw[thick] (1.9,1.45) -- (B1) (0.565,1.45) -- (C1) (2.3,1.45) -- (D1);
    \draw[very thick] (0.62,2.8) -- (3.42,2.8);

    \node[above left,inner sep=3pt,xshift=-2pt] at (A) {\(A\)};
    \node[right,inner sep=3pt] at (B) {\(B\)};
    \node[below right,inner sep=3pt] at (B1) {\(B_1\)};
    \node[left,inner sep=4pt] at (C) {\(C\)};
    \node[below right,inner sep=3pt] at (C1) {\(C_1\)};
    \node[left,inner sep=4pt] at (D) {\(D\)};
    \node[left,inner sep=3pt] at (D1) {\(D_1\)};
    \node[above right,inner sep=3pt,xshift=2pt] at (E) {\(E\)};
    \node[above right,inner sep=3pt] at (0,1.45) {\(P\)};
    \node[below right,inner sep=4pt] at (0,5.9) {\(Q\)};
\end{tikzpicture}

    \caption{Плоскости \(P\) и \(Q\) пересекаются по прямой \(AE\)}\label{fig:kiselev-stereo-303}
\end{figure}

Что плоскости \(P\) и \(Q\) пересекаются по прямой \(AE\) (а не касаются), следует из того, что прямые \(BB_1\) и \(CC_1\), содержащиеся в плоскости \(Q\), пересекаются с \(P\).

\kiselevparagraph{para:intersection-is-a-line} \textbf{Следствие. Пересечение двух плоскостей --- всегда прямая линия.}

Действительно, если плоскости пересекаются, то они имеют общую точку, но в таком случае они имеют и общую прямую.
% Печатная с. 275.
У плоскостей не может быть общей точки вне этой прямой, иначе плоскости слились бы в одну (\hyperref[item:plane-through-line-and-point]{см.~§~\ref*{para:plane-uniqueness-corollaries}, пункт~\ref*{item:plane-through-line-and-point}}).
  % Глава I
    \chapter{Перпендикуляр к плоскости и наклонные к ней}\label{ch:ch2}

\begingroup
% Интервалы вокруг рисунков в главе с десятью схемами.
\setlength{\intextsep}{10pt plus 2pt minus 2pt}
\setlength{\textfloatsep}{12pt plus 2pt minus 2pt}

% А. П. Киселёв, «Элементарная геометрия», 23-е издание, 1914. Печатные с. 275–282.
% Номера страниц DjVu на Wikimedia Commons на 11 больше печатных.
\newcommand*{\kiselevPlanimetryPage}[2]{%
    \href{https://commons.wikimedia.org/wiki/File:\%D0\%90._\%D0\%9F._\%D0\%9A\%D0\%B8\%D1\%81\%D0\%B5\%D0\%BB\%D1\%91\%D0\%B2._\%D0\%AD\%D0\%BB\%D0\%B5\%D0\%BC\%D0\%B5\%D0\%BD\%D1\%82\%D0\%B0\%D1\%80\%D0\%BD\%D0\%B0\%D1\%8F_\%D0\%B3\%D0\%B5\%D0\%BE\%D0\%BC\%D0\%B5\%D1\%82\%D1\%80\%D0\%B8\%D1\%8F_(1914).djvu?page=#1}{#2}%
}

% Печатная с. 275.
\kiselevparagraph{para:constructing-plane-normal} \textbf{Предварительное замечание.} Возьмём прямую \(AB\) (рис.~\ref{fig:kiselev-stereo-304}) и проведём через неё две произвольные плоскости \(P\) и \(Q\). Выберем на \(AB\) точку \(C\) и проведём через неё в плоскостях \(P\) и \(Q\) прямые \(CD\perp AB\) и \(CE\perp AB\) соответственно. Через пересекающиеся прямые \(CD\) и \(CE\) проведём плоскость \(R\) (\hyperref[item:plane-through-intersecting-lines]{см.~§~\ref*{para:plane-uniqueness-corollaries}, пункт~\ref*{item:plane-through-intersecting-lines}}). В результате прямая \(AB\) пересекает плоскость \(R\) и перпендикулярна двум прямым \(CD\) и \(CE\), проведённым в этой плоскости через точку пересечения. Докажем, что такое расположение прямой и плоскости обладает следующим важным свойством.

\begin{figure}[htbp]
    \centering
    \renewcommand{\thefigure}{304}
    \captionsetup{name=Рис.,labelsep=period}
    \begin{tikzpicture}[line cap=round,line join=round,every node/.style={inner sep=3pt}]
    \coordinate (A) at (-1.4,2.5);
    \coordinate (B) at (-1.4,-2.7);
    \coordinate (C) at (-1.4,0);
    \coordinate (D) at (2.05,0.78);
    \coordinate (E) at (1.2,-0.92);
    \coordinate (Pbottom) at (2.05,-1.92);
    \coordinate (Qbottom) at (1.2,-3.62);

    % Плоскость R и скрытая часть её заднего края.
    \draw[thick] (-1.4,1.26) -- (-2.05,1.26) -- (-3,-1.38);
    \draw[thick,densely dashed] (-1.4,1.26) -- (2.05,1.26);
    \draw[very thick] (-3,-1.38) -- (2.7,-1.38) -- (3.65,1.26) -- (2.05,1.26);

    % Плоскости P и Q проходят через общую прямую AB.
    \draw[thick] (A) -- (2.05,3.28) -- (D) -- (C)
        (A) -- (1.2,1.58) -- (E) -- (C) -- (A);
    \draw[thick,densely dashed] (C) -- (-1.4,-1.38)
        (D) -- (2.05,-1.38) (E) -- (1.2,-1.38) (B) -- (Pbottom);
    \draw[thick] (-1.4,-1.38) -- (B) -- (Qbottom) -- (1.2,-1.38)
        (2.05,-1.38) -- (Pbottom);

    \node[above left] at (A) {\(A\)};
    \node[left] at (B) {\(B\)};
    \node[left] at (C) {\(C\)};
    \node[right] at (D) {\(D\)};
    \node[above right] at (E) {\(E\)};
    \node at (1.5,2.85) {\(P\)};
    \node at (0.5,1.49) {\(Q\)};
    \node at (-2.48,-1.02) {\(R\)};
\end{tikzpicture}

    \caption{Расположение прямой \(AB\) и плоскости \(R\)}\label{fig:kiselev-stereo-304}
\end{figure}

\kiselevparagraph{para:perpendicular-plane-criterion} \textbf{Теорема. Если прямая} (\(AA_1\), рис.~\ref{fig:kiselev-stereo-305})\textbf{, пересекающая плоскость} \(P\)\textbf{, перпендикулярна двум прямым} (\(CD\) и \(EF\))\textbf{, проведённым в этой плоскости через точку пересечения} \(B\)\textbf{, то она перпендикулярна и любой третьей прямой} \(MN\)\textbf{, проведённой в плоскости через ту же точку пересечения.}

\begin{figure}[htbp]
    \centering
    \renewcommand{\thefigure}{305}
    \captionsetup{name=Рис.,labelsep=period}
    \begin{tikzpicture}[line cap=round,line join=round,every node/.style={inner sep=3pt}]
    \coordinate (A) at (0,3.65);
    \coordinate (A1) at (0,-3.65);
    \coordinate (B) at (0,0);
    \coordinate (C) at (-2.1,0);
    \coordinate (D) at (2.15,0);
    \coordinate (E) at (0.85,0.88);
    \coordinate (F) at (-0.85,-0.88);
    \coordinate (M) at (-1.9,0.704);
    \path (D) -- (F) coordinate[pos=0.4] (N);

    \draw[thick] (-3.3,-1.2) -- (-2.15,1.75) -- (3.85,1.75);
    \draw[very thick] (-3.3,-1.2) -- (2.7,-1.2) -- (3.85,1.75);
    \draw[thick] (C) -- (D) (E) -- (F) (M) -- (N) (D) -- (F);

    \draw[thick] (B) -- (A) -- (D) (A) -- (N) (A) -- (F);
    % Разрывы линий к A_1 показывают, где они скрыты плоскостью P.
    \path (A1) -- (D) coordinate[pos=0.6712] (frontD);
    \path (A1) -- (N) coordinate[pos=0.7429] (frontN);
    \path (A1) -- (F) coordinate[pos=0.8845] (frontF);
    \coordinate (frontB) at (0,-1.2);
    \draw[thick] (A1) -- (frontB) (A1) -- (frontD) (A1) -- (frontN) (A1) -- (frontF);
    \draw[thick,densely dashed] (frontB) -- (B) (frontD) -- (D)
        (frontN) -- (N) (frontF) -- (F);

    \foreach \point in {A,A1,B,D,F,N} {
        \filldraw[fill=white] (\point) circle[radius=1pt];
    }
    \node[above] at (A) {\(A\)};
    \node[below] at (A1) {\(A_1\)};
    \node[above left,xshift=-2pt,yshift=4pt] at (B) {\(B\)};
    \node[left] at (C) {\(C\)};
    \node[right] at (D) {\(D\)};
    \node[below right,xshift=2pt] at (E) {\(E\)};
    \node[left] at (F) {\(F\)};
    \node[above right] at (M) {\(M\)};
    \node[below right,xshift=1pt,yshift=-1pt] at (N) {\(N\)};
    \node at (-2.55,-0.8) {\(P\)};
\end{tikzpicture}

    \caption{Прямая \(AA_1\), перпендикулярная к \(CD\) и \(EF\)}\label{fig:kiselev-stereo-305}
\end{figure}

Отложим равные отрезки \(BA\) и \(BA_1\) произвольной длины и проведём в плоскости прямую \(DF\), которая пересекает прямые \(CD\), \(MN\) и \(EF\).
% Печатная с. 276.
Точки пересечения \(D\), \(N\) и \(F\) соединим с точками \(A\) и \(A_1\). Получим несколько треугольников и рассмотрим их по порядку.

Сначала возьмём треугольники \(AFD\) и \(A_1FD\). Они равны, потому что \(FD\) --- общая сторона, \(AD=A_1D_1\)\footnote{\emph{Прим.~ред.} На рис.~\ref{fig:kiselev-stereo-305} из точки \(D\) проведены наклонные \(AD\) и \(A_1D\).} как наклонные к прямой \(AA_1\), основания которых одинаково удалены от основания перпендикуляра \(BD\), и по той же причине \(AF=A_1F\) (\kiselevPlanimetryPage{53}{см.~Планиметрия, §~59, пункт~2}). Из равенства этих треугольников следует, что \(\angle ADF=\angle A_1DF\).

Теперь рассмотрим треугольники \(ADN\) и \(A_1DN\). Они равны, так как \(DN\) --- общая сторона, \(AD=A_1D\) и \(\angle ADN=\angle A_1DN\). Следовательно, \(AN=A_1N\).

Треугольники \(ABN\) и \(A_1BN\) также равны, поскольку их стороны соответственно равны. Поэтому \(\angle ABN=\angle A_1BN\). Значит, \(AA_1\perp MN\).

\kiselevparagraph{para:normal-and-oblique-definitions} \textbf{Определение.} Прямая называется \emph{перпендикулярной к плоскости}, если она пересекает эту плоскость и образует \emph{прямые углы} со всеми прямыми, проведёнными в плоскости через точку пересечения. В этом случае говорят также, что плоскость перпендикулярна прямой.

Прямая, которая пересекает плоскость и не перпендикулярна ей, называется \emph{наклонной}. Точка пересечения прямой с плоскостью называется \emph{основанием} перпендикуляра или наклонной.

Возможность провести прямую, перпендикулярную плоскости, следует из \hyperref[para:perpendicular-plane-criterion]{теоремы §~\ref*{para:perpendicular-plane-criterion}} и \hyperref[para:constructing-plane-normal]{построения в §~\ref*{para:constructing-plane-normal}}.

\kiselevparagraph{para:perpendiculars-in-one-plane} \textbf{Теорема. Если через одну и ту же точку прямой провести в пространстве сколько угодно перпендикуляров к ней, то все они лежат в одной плоскости}, перпендикулярной этой прямой.

% Печатная с. 277.
Дана прямая \(AB\) (рис.~\ref{fig:kiselev-stereo-306}) и точка \(B\) на ней. Проведём через \(B\) несколько перпендикуляров к \(AB\). Для этого проведём через \(AB\) плоскости \(P\), \(Q\), \(R\), \(S\) и так далее, а в каждой из них через точку \(B\) проведём прямую, перпендикулярную \(AB\). Пусть это будут прямые \(BC\), \(BD\), \(BE\), \(BF\) и так далее. Нужно доказать, что все они лежат в одной плоскости, перпендикулярной \(AB\).

\begin{figure}[htbp]
    \centering
    \renewcommand{\thefigure}{306}
    \captionsetup{name=Рис.,labelsep=period}
    \begin{tikzpicture}[line cap=round,line join=round,every node/.style={inner sep=3pt}]
    \coordinate (A) at (0,3.45);
    \coordinate (B) at (0,-0.1);
    \coordinate (C) at (2.7,0.66);
    \coordinate (D) at (2.25,-1.6);
    \coordinate (E) at (-1.86,-1.6);
    \coordinate (F) at (-2.28,-0.8);

    \draw[thick] (-3.65,-2.35) -- (-2.37,1.35) -- (-2.28,1.35);
    \draw[thick,densely dashed] (-2.28,1.35) -- (2.7,1.35);
    \draw[very thick] (-3.65,-2.35) -- (2.97,-2.35) -- (4.25,1.35) -- (2.7,1.35);

    \draw[thick] (A) -- (2.7,4.21) -- (C)
        (A) -- (2.25,1.95) -- (D) -- (B) -- (E) -- (-1.86,1.95) -- (A)
        (A) -- (-2.28,2.75) -- (F)
        (A) -- (B);
    \draw[thick,densely dashed] (C) -- (B) -- (F);

    \node[above] at (A) {\(A\)};
    \node[below] at (B) {\(B\)};
    \node[right] at (C) {\(C\)};
    \node[right] at (D) {\(D\)};
    \node[below] at (E) {\(E\)};
    \node[left] at (F) {\(F\)};
    \node at (2.32,3.7) {\(P\)};
    \node at (1.85,1.87) {\(Q\)};
    \node at (-1.52,1.78) {\(R\)};
    \node at (-2.0,2.48) {\(S\)};
    \node at (-3.02,-1.96) {\(M\)};
\end{tikzpicture}

    \caption{Перпендикуляры к прямой \(AB\)}\label{fig:kiselev-stereo-306}
\end{figure}

Через любые два из этих перпендикуляров, например \(BC\) и \(BD\), проведём плоскость \(M\). По \hyperref[para:perpendicular-plane-criterion]{теореме §~\ref*{para:perpendicular-plane-criterion}} она перпендикулярна \(AB\). Предположим, что один из остальных перпендикуляров, например \(BE\), проведённый в плоскости \(R\), не лежит в плоскости \(M\). Тогда плоскости \(R\) и \(M\) пересекаются по прямой, отличной от \(BE\). Получается, что в плоскости \(R\) через точку \(B\) проходят два перпендикуляра к \(AB\): \(BE\) и прямая пересечения плоскостей \(M\) и \(R\) (\hyperref[para:perpendicular-plane-criterion]{см.~§~\ref*{para:perpendicular-plane-criterion}}). Это невозможно. Значит, все проведённые перпендикуляры лежат в плоскости \(R\)\footnote{\emph{Прим.~ред.} В начале доказательства плоскость через \(BC\) и \(BD\) обозначена буквой \(M\).}, которая перпендикулярна \(AB\).

\kiselevparagraph{para:unique-plane-normal-to-line} \textbf{Теорема. Через любую точку пространства можно провести ровно одну плоскость, перпендикулярную данной прямой.}

Рассмотрим два случая: точка лежит на данной прямой или вне её.

\begin{figure}[htbp]
    \centering
    \renewcommand{\thefigure}{307}
    \captionsetup{name=Рис.,labelsep=period}
    \begin{tikzpicture}[line cap=round,line join=round,every node/.style={inner sep=3pt}]
    \coordinate (A) at (-1.05,2.75);
    \coordinate (B) at (-1.05,-2.65);
    \coordinate (C) at (-1.05,0);
    \coordinate (D) at (1.9,0.55);
    \coordinate (E) at (0.92,-0.72);
    \coordinate (Pbottom) at (1.9,-2.1);
    \coordinate (Qbottom) at (0.92,-3.37);

    \draw[thick] (-1.05,1.52) -- (-1.9,1.52) -- (-3.2,-1.58);
    \draw[thick,densely dashed] (-1.05,1.52) -- (1.9,1.52);
    \draw[very thick] (-3.2,-1.58) -- (2.77,-1.58) -- (4.07,1.52) -- (1.9,1.52);

    \draw[thick] (D) -- (1.9,3.3) -- (A) -- (0.92,2.03) -- (E) -- (C) -- (A);
    \draw[thick,densely dashed] (C) -- (D) -- (1.9,-1.58)
        (C) -- (-1.05,-1.58) (E) -- (0.92,-1.58) (B) -- (Pbottom);
    \draw[thick] (-1.05,-1.58) -- (B) -- (Qbottom) -- (0.92,-1.58)
        (1.9,-1.58) -- (Pbottom);

    \node[above left] at (A) {\(A\)};
    \node[left] at (B) {\(B\)};
    \node[left] at (C) {\(C\)};
    \node[below right] at (D) {\(D\)};
    \node[below right] at (E) {\(E\)};
    \node at (1.5,2.85) {\(P\)};
    \node at (0.35,1.85) {\(Q\)};
    \node at (-2.56,-1.14) {\(R\)};
\end{tikzpicture}

    \caption{Плоскость \(R\), перпендикулярная прямой \(AB\)}\label{fig:kiselev-stereo-307}
\end{figure}

\begin{enumerate}[label=\arabic*),ref=\arabic*]
    \item\label{item:normal-plane-point-on-line} Пусть точка \(C\) лежит на данной прямой \(AB\) (рис.~\ref{fig:kiselev-stereo-307}). Проведём через \(AB\) две плоскости \(P\) и \(Q\), а в них --- прямые \(CD\) и \(CE\), перпендикулярные \(AB\). Через эти пересекающиеся прямые проведём плоскость \(R\).
    % Печатная с. 278.
    Эта плоскость перпендикулярна \(AB\) в точке \(C\), поскольку \(CD\) и \(CE\) перпендикулярны \(AC\)\footnote{\emph{Прим.~ред.} На рис.~\ref{fig:kiselev-stereo-307} точки \(A\), \(B\) и \(C\) лежат на одной прямой.}. Следовательно, любая третья прямая плоскости \(R\), проходящая через \(C\), также перпендикулярна \(AB\) (\hyperref[para:perpendicular-plane-criterion]{см.~§~\ref*{para:perpendicular-plane-criterion}}).

    Чтобы доказать единственность этой плоскости, предположим, что через \(C\) проходит другая плоскость \(R_1\), перпендикулярная \(AB\). Она должна пересекать \(P\) по прямой, проходящей через \(C\) и перпендикулярной \(AB\). В плоскости \(P\) такая прямая только одна --- \(CD\). Значит, обе плоскости \(R_1\) и \(R\) пересекают \(P\) по \(CD\). Так же убедимся, что обе они пересекают плоскость \(Q\) по \(CE\). Следовательно, \(R_1\) и \(R\) проходят через одни и те же пересекающиеся прямые \(CD\) и \(CE\) и потому сливаются в одну (\hyperref[item:plane-through-intersecting-lines]{см.~§~\ref*{para:plane-uniqueness-corollaries}, пункт~\ref*{item:plane-through-intersecting-lines}}).

    \item\label{item:normal-plane-point-outside-line} Пусть точка \(D\) лежит вне данной прямой \(AB\) (рис.~\ref{fig:kiselev-stereo-307}). Проведём через \(D\) и \(AB\) плоскость \(P\), а через \(AB\) --- ещё одну плоскость \(Q\). В плоскости \(P\) опустим из точки \(D\) перпендикуляр \(DC\) на \(AB\), а в плоскости \(Q\) проведём из точки \(C\) перпендикуляр \(CE\) к \(AB\). Плоскость \(R\), проходящая через \(DC\) и \(CE\), проходит через \(D\) и перпендикулярна \(AB\) (\hyperref[para:perpendicular-plane-criterion]{см.~§~\ref*{para:perpendicular-plane-criterion}}).

    Предположим, что через \(D\) проходит ещё одна плоскость \(R_1\), перпендикулярная \(AB\). Она должна пересекать \(P\) по прямой, проходящей через \(D\) и перпендикулярной \(AB\). В плоскости \(P\) такая прямая только одна --- \(DC\).
    % Печатная с. 279.
    Значит, обе плоскости \(R_1\) и \(R\) пересекают \(P\) по \(DC\). В частности, \(R_1\) проходит через \(C\) и потому пересекает \(Q\) по \(CE\), единственной прямой плоскости \(Q\), проходящей через \(C\) и перпендикулярной \(AB\). Таким образом, \(R_1\) и \(R\) проходят через одни и те же пересекающиеся прямые \(DC\) и \(CE\) и сливаются в одну.
\end{enumerate}

\kiselevparagraph{para:unique-line-normal-to-plane} \textbf{Теорема. Через любую точку пространства можно провести ровно один перпендикуляр к данной плоскости.}

Рассмотрим два случая: точка лежит в данной плоскости или вне её.

\begin{enumerate}[label=\arabic*),ref=\arabic*]
    \item\label{item:normal-line-point-in-plane} Пусть точка \(O\) лежит в плоскости \(P\) (рис.~\ref{fig:kiselev-stereo-308}). Проведём через неё в этой плоскости произвольную прямую \(MN\), а затем через \(O\) проведём плоскость \(Q\), перпендикулярную \(MN\) (\hyperref[item:normal-plane-point-on-line]{см.~§~\ref*{para:unique-plane-normal-to-line}, пункт~\ref*{item:normal-plane-point-on-line}}). Пусть \(P\) и \(Q\) пересекаются по прямой \(AB\). В плоскости \(Q\) проведём \(OC\perp AB\). Докажем, что \(OC\) --- перпендикуляр к плоскости \(P\).

    \begin{figure}[htbp]
        \centering
        \renewcommand{\thefigure}{308}
        \captionsetup{name=Рис.,labelsep=period}
        \begin{tikzpicture}[line cap=round,line join=round,every node/.style={inner sep=3pt}]
    \coordinate (A) at (-1.3,0.72);
    \coordinate (B) at (1.4,-0.72);
    \path (A) -- (B) coordinate[midway] (O);
    \coordinate (C) at (0.05,2.5);
    \coordinate (M) at (-2.2,-1);
    \coordinate (N) at (2.075,0.9);
    \coordinate (Qedge) at (1.4,0.6);

    \draw[thick] (-3.2,-1.55) -- (-1.72,1.65) -- (-1.3,1.65);
    \draw[thick,densely dashed] (-1.3,1.65) -- (1.4,1.65);
    \draw[very thick] (-3.2,-1.55) -- (2.55,-1.55) -- (4.03,1.65) -- (1.4,1.65);

    \draw[thick] (A) -- (-1.3,3.22) -- (1.4,1.78) -- (B) -- cycle;
    \draw[very thick] (O) -- (C);
    \draw[thick] (M) -- (O) (Qedge) -- (N);
    \draw[thick,densely dashed] (O) -- (Qedge);

    \node[left] at (A) {\(A\)};
    \node[below right] at (B) {\(B\)};
    \node[above] at (C) {\(C\)};
    \node[below left] at (M) {\(M\)};
    \node[right] at (N) {\(N\)};
    \node[below] at (O) {\(O\)};
    \node at (2.97,1.15) {\(P\)};
    \node at (-0.96,2.67) {\(Q\)};
\end{tikzpicture}

        \caption{Прямая \(OC\) --- перпендикуляр к плоскости \(P\)}\label{fig:kiselev-stereo-308}
    \end{figure}

    Поскольку \(MN\) перпендикулярна плоскости \(Q\), она перпендикулярна любой прямой этой плоскости, проходящей через \(O\). В частности, \(OC\perp MN\). По построению \(OC\perp AB\). Значит, \(OC\) перпендикулярна двум прямым плоскости \(P\), проведённым через \(O\), и потому перпендикулярна самой плоскости (\hyperref[para:perpendicular-plane-criterion]{см.~§~\ref*{para:perpendicular-plane-criterion}}).

    Предположим, что через \(O\) проходит ещё один перпендикуляр \(OC_1\) к плоскости \(P\) (рис.~\ref{fig:kiselev-stereo-309}).
    % Печатная с. 280.
    Проведём через \(OC\) и \(OC_1\) плоскость, которая пересечёт \(P\) по прямой \(OH\). Эта плоскость на рисунке не показана. Тогда углы \(COH\) и \(C_1OH\) должны быть прямыми. Это невозможно, поскольку один из этих углов составляет часть другого. Значит, другого перпендикуляра к \(P\) через точку \(O\) провести нельзя.

    \begin{figure}[htbp]
        \centering
        \renewcommand{\thefigure}{309}
        \captionsetup{name=Рис.,labelsep=period}
        \begin{tikzpicture}[line cap=round,line join=round,every node/.style={inner sep=3pt}]
    \coordinate (O) at (-0.55,0);
    \coordinate (C) at (-0.55,3.05);
    \coordinate (C1) at (0.4,3.0);
    \coordinate (H) at (1.8,0.94);

    \draw[thick] (-3,-1.38) -- (-1.6,1.58) -- (3.45,1.58);
    \draw[very thick] (-3,-1.38) -- (2.05,-1.38) -- (3.45,1.58);
    \draw[thick] (C) -- (O) -- (C1) (O) -- (H);

    \node[left] at (O) {\(O\)};
    \node[above] at (C) {\(C\)};
    \node[above] at (C1) {\(C_1\)};
    \node[right] at (H) {\(H\)};
    \node at (-2.3,-0.98) {\(P\)};
\end{tikzpicture}

        \caption{Предположение о другом перпендикуляре \(OC_1\) через точку \(O\)}\label{fig:kiselev-stereo-309}
    \end{figure}

    \item\label{item:normal-line-point-outside-plane} Пусть точка \(O\) лежит вне плоскости \(P\) (рис.~\ref{fig:kiselev-stereo-310}). Проведём в плоскости \(P\) произвольную прямую \(MN\), а через \(O\) --- плоскость \(Q\), перпендикулярную \(MN\) (\hyperref[item:normal-plane-point-outside-line]{см.~§~\ref*{para:unique-plane-normal-to-line}, пункт~\ref*{item:normal-plane-point-outside-line}}). Пусть \(P\) и \(Q\) пересекаются по прямой \(AB\). В плоскости \(Q\) проведём \(OC\perp AB\). Докажем, что \(OC\) --- перпендикуляр к плоскости \(P\).

    \begin{figure}[htbp]
        \centering
        \renewcommand{\thefigure}{310}
        \captionsetup{name=Рис.,labelsep=period}
        \begin{tikzpicture}[line cap=round,line join=round,every node/.style={inner sep=3pt}]
    \coordinate (A) at (1.55,-0.7);
    \coordinate (B) at (-0.9,0.75);
    \path (A) -- (B) coordinate[midway] (C);
    \coordinate (O) at (0.325,3.5);
    \coordinate (O1) at (0.325,-3.45);
    \coordinate (M) at (-2.37,-1.2);
    \coordinate (N) at (3.118,-0.5);
    \path (M) -- (A) coordinate[pos=0.25] (D);

    \draw[thick] (-3.25,-1.55) -- (-1.95,1.6) -- (-0.9,1.6);
    \draw[thick,densely dashed] (-0.9,1.6) -- (1.55,1.6);
    \draw[very thick] (-3.25,-1.55) -- (2.75,-1.55) -- (4.05,1.6) -- (1.55,1.6);

    \draw[thick] (B) -- (-0.9,4.75) -- (1.55,3.3) -- (A)
        (-0.9,-1.55) -- (-0.9,-3.25) -- (1.55,-4.7) -- (1.55,-1.55);
    \draw[thick,densely dashed] (B) -- (-0.9,-1.55) (A) -- (1.55,-1.55);
    \draw[thick] (B) -- (A) (M) -- (N);
    \draw[thick,densely dashed] (O) -- (A) -- (O1) -- (D) -- (O) (D) -- (C);
    \draw[very thick] (O) -- (C);
    \draw[thick,densely dashed] (C) -- (O1);

    \foreach \point in {A,B,C,D,O,O1} {
        \filldraw[fill=white] (\point) circle[radius=1pt];
    }
    \node[right,xshift=2pt,yshift=-5pt] at (A) {\(A\)};
    \node[left] at (B) {\(B\)};
    \node[above right,yshift=2pt] at (C) {\(C\)};
    \node[above left] at (D) {\(D\)};
    \node[above] at (M) {\(M\)};
    \node[above left] at (N) {\(N\)};
    \node[right,inner sep=4pt] at (O) {\(O\)};
    \node[right,inner sep=4pt] at (O1) {\(O_1\)};
    \node at (3,1.16) {\(P\)};
    \node at (1.16,2.85) {\(Q\)};
\end{tikzpicture}

        \caption{Перпендикуляр \(OC\) из точки \(O\) вне плоскости \(P\)}\label{fig:kiselev-stereo-310}
    \end{figure}

    Достаточно показать, что \(OC\) перпендикулярна двум прямым, проведённым в плоскости \(P\) через \(C\). Одной такой прямой служит \(AB\), которой \(OC\) перпендикулярна по построению. Второй возьмём прямую \(CD\), соединяющую \(C\) с произвольной точкой \(D\) прямой \(MN\).

    Чтобы доказать, что \(OC\perp CD\), продолжим \(OC\) на отрезок \(CO_1=OC\) и соединим точки \(O\) и \(O_1\) с \(A\) и \(D\). Так как \(MN\perp Q\), прямая \(MN\) перпендикулярна \(AO\) и \(AO_1\). Поэтому треугольники \(OAD\) и \(O_1AD\) прямоугольные при вершине \(A\). Катет \(AD\) у них общий. Катеты \(AO\) и \(AO_1\) равны как наклонные, проведённые в плоскости \(Q\) из точки \(A\) к прямой \(OO_1\), поскольку расстояния от их оснований до основания перпендикуляра \(AC\) (отрезки \(OC\) и \(O_1C\)) равны (\kiselevPlanimetryPage{53}{см.~Планиметрия, §~59, пункт~2}). Из равенства этих треугольников следует, что \(OD=O_1D\).

    Треугольники \(OCD\) и \(O_1CD\) равны, поскольку их стороны соответственно равны.
    % Печатная с. 281.
    Значит, \(\angle OCD=\angle O_1CD\), и потому эти углы прямые. Итак, \(OC\perp CA\) и \(OC\perp CD\), следовательно, \(OC\perp P\).

    Предположим, что из \(O\) можно опустить на плоскость \(P\) другой перпендикуляр \(OC_1\) (рис.~\ref{fig:kiselev-stereo-311}). Соединим \(C\) и \(C_1\). Тогда у треугольника \(OCC_1\) два угла, при вершинах \(C\) и \(C_1\), должны быть прямыми. Это невозможно, поэтому второго перпендикуляра из \(O\) к плоскости \(P\) провести нельзя.

    \begin{figure}[htbp]
        \centering
        \renewcommand{\thefigure}{311}
        \captionsetup{name=Рис.,labelsep=period}
        \begin{tikzpicture}[line cap=round,line join=round,every node/.style={inner sep=3pt}]
    \coordinate (O) at (0.5,2.92);
    \coordinate (C) at (0.5,0);
    \coordinate (C1) at (-1.53,-0.6);

    \draw[thick] (-3.15,-1.46) -- (-1.65,1.45) -- (3.5,1.45);
    \draw[very thick] (-3.15,-1.46) -- (2,-1.46) -- (3.5,1.45);
    \draw[thick] (O) -- (C1) -- (C) -- cycle;

    \node[above] at (O) {\(O\)};
    \node[right] at (C) {\(C\)};
    \node[below left] at (C1) {\(C_1\)};
    \node at (1.66,0.91) {\(P\)};
\end{tikzpicture}

        \caption{Предположение о другом перпендикуляре \(OC_1\) из точки \(O\)}\label{fig:kiselev-stereo-311}
    \end{figure}
\end{enumerate}

\kiselevparagraph{para:normal-and-oblique-lengths} \textbf{Сравнительная длина перпендикуляра и наклонных.} Если из одной точки \(A\) (рис.~\ref{fig:kiselev-stereo-312}) к плоскости проведены перпендикуляр \(AB\) и наклонная \(AC\), то \emph{проекцией наклонной на плоскость} называют прямую \(BC\), проведённую между основаниями перпендикуляра и наклонной.

\begin{figure}[htbp]
    \centering
    \renewcommand{\thefigure}{312}
    \captionsetup{name=Рис.,labelsep=period}
    \begin{tikzpicture}[line cap=round,line join=round,every node/.style={inner sep=3pt}]
    \coordinate (A) at (0,3.45);
    \coordinate (B) at (0,0);
    \coordinate (C) at (2,0.28);
    \coordinate (D) at (-0.88,-1.1);
    \coordinate (E) at (-2.18,-1.01);

    \draw[thick] (-3.4,-1.91) -- (-1.95,1.38) -- (3.6,1.38);
    \draw[very thick] (-3.4,-1.91) -- (2.15,-1.91) -- (3.6,1.38);
    \draw[very thick] (A) -- (B);
    \draw[thick] (A) -- (C) -- (B) -- (D) -- (A) -- (E) -- (B);

    \node[above] at (A) {\(A\)};
    \node[above right] at (B) {\(B\)};
    \node[right] at (C) {\(C\)};
    \node[below right] at (D) {\(D\)};
    \node[below left] at (E) {\(E\)};
    \node at (2.88,0.97) {\(P\)};
\end{tikzpicture}

    \caption{Сравнительная длина перпендикуляра и наклонных}\label{fig:kiselev-stereo-312}
\end{figure}

\textbf{Теоремы. Если из одной и той же точки} \(A\)\textbf{, взятой вне плоскости} \(P\) \textbf{(рис.~\ref{fig:kiselev-stereo-312}), к этой плоскости проведены перпендикуляр} \(AB\) \textbf{и какие-нибудь наклонные} (\(AC\), \(AD\), \(AE,\ldots\))\textbf{, то:}

\begin{enumerate}[label=\arabic*),ref=\arabic*]
    \item \textbf{перпендикуляр короче любой наклонной;}
    \item\label{item:equal-projections-equal-obliques} \textbf{наклонные с равными проекциями равны;}
    \item \textbf{из двух наклонных больше та, у которой больше проекция.}
\end{enumerate}

Вращая прямоугольные треугольники \(ABC\) и \(ABD\) вокруг катета \(AB\), можно совместить их плоскости с плоскостью треугольника \(ABE\). Тогда все наклонные окажутся в одной плоскости с перпендикуляром, а все проекции будут лежать на одной прямой.
% Печатная с. 282.
Таким образом, эти теоремы сводятся к аналогичным теоремам планиметрии (см.~Планиметрия, §~59, \kiselevPlanimetryPage{52}{пункт~1} и \kiselevPlanimetryPage{53}{пункт~2}).

\textbf{Замечание.} \emph{Расстоянием} от точки \(A\) до плоскости \(P\) называют длину перпендикуляра, опущенного из этой точки на плоскость (рис.~\ref{fig:kiselev-stereo-312}).

\kiselevparagraph{para:obliques-determine-projections} \textbf{Обратные теоремы. Если из одной и той же точки вне плоскости провести к этой плоскости перпендикуляр и наклонные, то у равных наклонных проекции равны. Из двух проекций больше та, которая соответствует большей наклонной.}

Доказательство от противного оставляем читателю.

\kiselevparagraph{para:three-perpendiculars} Приведём ещё одну теорему о перпендикулярах, которая понадобится дальше.\footnote{В некоторых курсах геометрии эту теорему или её вариант называют \emph{теоремой трёх перпендикуляров}. Она связывает три перпендикуляра: \(AB\) к плоскости \(P\), \(BC\) к прямой \(DE\) и \(AC\) к той же прямой \(DE\).}

\textbf{Теорема. Прямая} \(DE\) (рис.~\ref{fig:kiselev-stereo-313})\textbf{, проведённая в плоскости} \(P\) \textbf{через основание наклонной} \(AC\) \textbf{перпендикулярно её проекции} \(BC\)\textbf{, перпендикулярна и самой наклонной.}

\begin{figure}[htbp]
    \centering
    \renewcommand{\thefigure}{313}
    \captionsetup{name=Рис.,labelsep=period}
    \begin{tikzpicture}[line cap=round,line join=round,every node/.style={inner sep=3pt}]
    \coordinate (A) at (-1.26,3.62);
    \coordinate (B) at (-1.26,0.2);
    \coordinate (D) at (0.35,-1.31);
    \coordinate (E) at (1.51,0.95);
    \path (D) -- (E) coordinate[midway] (C);

    \draw[thick] (-3.25,-1.95) -- (-1.92,1.53) -- (3.45,1.53);
    \draw[very thick] (-3.25,-1.95) -- (2.12,-1.95) -- (3.45,1.53);
    \draw[very thick] (A) -- (B);
    \draw[thick] (A) -- (C) -- (B) (D) -- (E);
    \draw[thick,densely dashed] (D) -- (A) -- (E) -- (B) -- cycle;

    \node[above left] at (A) {\(A\)};
    \node[left] at (B) {\(B\)};
    \node[right] at (C) {\(C\)};
    \node[below right] at (D) {\(D\)};
    \node[right] at (E) {\(E\)};
    \node at (-2.57,-1.53) {\(P\)};
\end{tikzpicture}

    \caption{Теорема трёх перпендикуляров}\label{fig:kiselev-stereo-313}
\end{figure}

Отложим произвольные, но равные отрезки \(CD\) и \(CE\) и соединим точки \(A\) и \(B\) с \(D\) и \(E\). Тогда \(BD=BE\) как наклонные к прямой \(DE\), основания которых одинаково удалены от основания перпендикуляра \(BC\) (\kiselevPlanimetryPage{53}{см.~Планиметрия, §~59, пункт~2}). Значит, \(AD=AE\), поскольку у этих наклонных к плоскости \(P\) равные проекции \(BD\) и \(BE\) (\hyperref[item:equal-projections-equal-obliques]{см.~§~\ref*{para:normal-and-oblique-lengths}, пункт~\ref*{item:equal-projections-equal-obliques}}). Поэтому треугольник \(ADE\) равнобедренный, и его медиана \(AC\) перпендикулярна основанию \(DE\) (\kiselevPlanimetryPage{39}{см.~Планиметрия, §~39}).

\clearpage
\endgroup
  % Глава II
    \chapter{Параллельные прямые и плоскости}\label{ch:ch3}

\begingroup
\setlength{\intextsep}{10pt plus 2pt minus 2pt}
\setlength{\textfloatsep}{12pt plus 2pt minus 2pt}

% А. П. Киселёв, «Элементарная геометрия», 23-е издание, 1914. Печатные с. 283–292.
% Страницы скана планиметрии на 11 больше печатных.
\newcommand*{\kiselevPlanimetryPage}[2]{%
    \href{https://commons.wikimedia.org/wiki/File:\%D0\%90._\%D0\%9F._\%D0\%9A\%D0\%B8\%D1\%81\%D0\%B5\%D0\%BB\%D1\%91\%D0\%B2._\%D0\%AD\%D0\%BB\%D0\%B5\%D0\%BC\%D0\%B5\%D0\%BD\%D1\%82\%D0\%B0\%D1\%80\%D0\%BD\%D0\%B0\%D1\%8F_\%D0\%B3\%D0\%B5\%D0\%BE\%D0\%BC\%D0\%B5\%D1\%82\%D1\%80\%D0\%B8\%D1\%8F_(1914).djvu?page=#1}{#2}%
}

\section*{Параллельные прямые}
\addcontentsline{toc}{section}{Параллельные прямые}

% Печатная с. 283.
\kiselevparagraph{para:skew-and-parallel-lines} \textbf{Предварительное замечание.} Две прямые могут быть расположены в пространстве так, что через них нельзя провести плоскость. Возьмём, например, две такие прямые \(AB\) и \(DE\) (рис.~\ref{fig:kiselev-stereo-314}), из которых одна пересекает некоторую плоскость \(P\), а другая лежит в ней, но не проходит через точку пересечения \(C\). Через такие две прямые нельзя провести плоскость, потому что в противном случае через прямую \(DE\) и точку \(C\) проходили бы две различные плоскости: одна --- \(P\), пересекающая прямую \(AB\), и другая, содержащая её. Это невозможно (\hyperref[item:plane-through-line-and-point]{см.~§~\ref*{para:plane-uniqueness-corollaries}, пункт~\ref*{item:plane-through-line-and-point}}).

\begin{figure}[htbp]
    \centering
    \renewcommand{\thefigure}{314}
    \captionsetup{name=Рис.,labelsep=period}
    \begin{tikzpicture}[line cap=round,line join=round]
    \coordinate (A) at (-0.95,2);
    \coordinate (C) at (0.85,0.25);
    \coordinate (B) at (3.25,-2.083);
    \path (C) -- (B) coordinate[pos=0.643] (front);
    \draw[thick] (-3.15,-1.25) -- (-2.25,1.15) -- (3.35,1.15);
    \draw[very thick] (-3.15,-1.25) -- (2.45,-1.25) -- (3.35,1.15);
    \draw[thick] (A) -- (C) (front) -- (B);
    \draw[thick,densely dashed] (C) -- (front);
    \draw[thick] (-1.96,-0.6) node[left] {\(D\)} -- (0.86,-0.6) node[right] {\(E\)};
    \node[above left] at (A) {\(A\)};
    \node[right] at (B) {\(B\)};
    \node[above right] at (C) {\(C\)};
    \node at (-1.96,0.74) {\(P\)};
\end{tikzpicture}

    \caption{Прямые \(AB\) и \(DE\), через которые нельзя провести плоскость}\label{fig:kiselev-stereo-314}
\end{figure}

Две прямые, не лежащие в одной плоскости, конечно, не пересекаются, сколько бы их ни продолжали. Однако их не называют параллельными. Это название оставляют только для тех прямых, которые, \emph{находясь в одной плоскости}, не пересекаются, сколько бы их ни продолжали.

Две прямые, не лежащие в одной плоскости, называются \emph{скрещивающимися}.

В планиметрии мы видели (см.~Планиметрия, \kiselevPlanimetryPage{70}{§~77} и \kiselevPlanimetryPage{71}{§~79}), что через любую точку \emph{плоскости} можно провести прямую, параллельную данной прямой, и притом только одну. То же самое можно сказать о любой точке \emph{пространства}, потому что через точку и данную прямую можно провести плоскость, и только одну.

% Печатная с. 284.
\kiselevparagraph{para:plane-intersects-parallel-lines} \textbf{Теорема. Если плоскость} \(P\) (рис.~\ref{fig:kiselev-stereo-315}) \textbf{пересекает одну из параллельных прямых} \(AB\)\textbf{, то она пересекает и другую} \(CD\)\textbf{.}

\begin{figure}[htbp]
    \centering
    \renewcommand{\thefigure}{315}
    \captionsetup{name=Рис.,labelsep=period}
    \begin{tikzpicture}[line cap=round,line join=round]
    \coordinate (A) at (-0.5,2.9);
    \coordinate (B) at (-1.55,-0.45);
    \coordinate (C) at (1.6,2.9);
    \coordinate (F) at (0.55,-0.45);
    \path (F) -- (C) coordinate[pos=0.5] (D);
    \draw[thick] (-2.9,-1.45) -- (-1.95,1.23) -- (2.9,1.23);
    \draw[very thick] (-2.9,-1.45) -- (1.95,-1.45) -- (2.9,1.23);
    \draw[thick] (A) -- (B) -- (1.53,-0.45) node[right] {\(E\)};
    \draw[thick] (C) -- (D);
    \draw[thick,densely dashed] (D) -- (F);
    \node[above right] at (A) {\(A\)};
    \node[left] at (B) {\(B\)};
    \node[above right] at (C) {\(C\)};
    \node[above right] at (D) {\(D\)};
    \node[above left] at (F) {\(F\)};
    \node at (2.4,0.76) {\(P\)};
\end{tikzpicture}

    \caption{Плоскость \(P\) пересекает параллельные прямые \(AB\) и \(CD\)}\label{fig:kiselev-stereo-315}
\end{figure}

Проведём через \(AB\) и \(CD\) плоскость. Эта плоскость содержит точку \(B\), в которой прямая \(AB\) пересекается с \(P\). Значит, эта плоскость пересекается с \(P\) по некоторой прямой \(BE\) (\hyperref[para:intersecting-planes]{см.~§~\ref*{para:intersecting-planes}}). Эта прямая, находясь в одной плоскости с \(AB\) и \(CD\) и пересекая одну из этих параллельных, должна пересечь и другую (\kiselevPlanimetryPage{71}{см.~Планиметрия, §~80}) в некоторой точке \(F\). Точка \(F\), находясь одновременно на прямой \(BE\) и на прямой \(CD\), должна быть точкой пересечения плоскости \(P\) с прямой \(CD\).

\kiselevparagraph{para:plane-perpendicular-to-parallel-lines} \textbf{Теорема. Если плоскость} \(P\) (рис.~\ref{fig:kiselev-stereo-316}) \textbf{перпендикулярна к одной из параллельных прямых} \(AB\)\textbf{, то она перпендикулярна и к~другой}~\(CD\)\textbf{.}

\begin{figure}[htbp]
    \centering
    \renewcommand{\thefigure}{316}
    \captionsetup{name=Рис.,labelsep=period}
    \begin{tikzpicture}[line cap=round,line join=round]
    \coordinate (A) at (-1.4,3);
    \coordinate (B) at (-1.4,-0.2);
    \coordinate (C) at (0.95,3);
    \coordinate (D) at (0.95,-0.2);
    \coordinate (M) at (-1.4,2.24);
    \draw[thick] (-2.95,-1.63) -- (-1.78,1.4) -- (3.15,1.4);
    \draw[very thick] (-2.95,-1.63) -- (1.98,-1.63) -- (3.15,1.4);
    \draw[thick] (A) -- (B) -- (D) -- (C) (M) -- (D);
    % Верхний край плоскости R на рисунке Киселёва показан линией обрыва.
    \draw[thick] (A) sin (-1.204,3.08) cos (-1.008,3)
        sin (-0.812,2.92) cos (-0.616,3)
        sin (-0.42,3.08) cos (-0.224,3)
        sin (-0.028,2.92) cos (0.168,3)
        sin (0.364,3.08) cos (0.56,3)
        sin (0.755,2.92) cos (C);
    \draw[thick] (0.43,-1.13) node[left,yshift=2pt] {\(E\)} -- (D) -- (1.47,0.73);
    \node[above left] at (A) {\(A\)};
    \node[below left] at (B) {\(B\)};
    \node[above right] at (C) {\(C\)};
    \node[right] at (D) {\(D\)};
    \node[above right] at (M) {\(M\)};
    \node at (-0.25,2.12) {\(R\)};
    \node at (-2.38,-1.25) {\(P\)};
\end{tikzpicture}

    \caption{Плоскость \(P\), перпендикулярная к параллельным прямым \(AB\) и \(CD\)}\label{fig:kiselev-stereo-316}
\end{figure}

Предстоит доказать, во-первых, что плоскость \(P\) пересекает прямую \(CD\), а во-вторых, что эта прямая перпендикулярна к каким-нибудь двум прямым, проведённым в плоскости \(P\) через основание \(CD\).

\begin{enumerate}[label=\arabic*),ref=\arabic*]
    \item Плоскость \(P\) должна пересечь \(CD\), потому что она, по условию, пересекает прямую \(AB\), параллельную \(CD\).
    \item Проведём через \(AB\) и \(CD\) плоскость \(R\) и возьмём её пересечение \(BD\) с плоскостью \(P\). Так как, по условию, \(AB\) перпендикулярна к \(P\), то \(AB\perp BD\). Поэтому и \(CD\perp BD\) (\kiselevPlanimetryPage{72}{см.~Планиметрия, §~82}). Проведём в плоскости \(P\) прямую \(DE\), перпендикулярную к \(BD\), и возьмём какую-нибудь наклонную \(MD\), для которой проекцией служит \(BD\). Прямая \(ED\), перпендикулярная к проекции \(BD\), должна быть перпендикулярной и к наклонной \(MD\) (\hyperref[para:three-perpendiculars]{см.~§~\ref*{para:three-perpendiculars}})
    % Печатная с. 285.
    и, следовательно, перпендикулярна к плоскости \(R\) (\hyperref[para:perpendicular-plane-criterion]{см.~§~\ref*{para:perpendicular-plane-criterion}}), а значит, и к прямой \(CD\). Таким образом, прямая \(CD\) оказывается перпендикулярной к двум прямым плоскости \(P\), именно к \(DB\) и к \(DE\). Следовательно, она перпендикулярна к этой плоскости.
\end{enumerate}

\kiselevparagraph{para:lines-normal-to-same-plane} \textbf{Обратная теорема. Если две прямые} \(AB\) и \(CD\) (рис.~\ref{fig:kiselev-stereo-317}) \textbf{перпендикулярны к одной и той же плоскости} \(P\)\textbf{, то они параллельны.}

\begin{figure}[htbp]
    \centering
    \renewcommand{\thefigure}{317}
    \captionsetup{name=Рис.,labelsep=period}
    \begin{tikzpicture}[line cap=round,line join=round]
    \coordinate (A) at (-1.25,2.83);
    \coordinate (B) at (-1.25,-0.5);
    \coordinate (C) at (0.5,2.83);
    \coordinate (D) at (0.5,-0.5);
    \coordinate (C1) at (1.25,2.83);
    \draw[thick] (-2.85,-1.52) -- (-1.88,1.02) -- (2.97,1.02);
    \draw[very thick] (-2.85,-1.52) -- (2,-1.52) -- (2.97,1.02);
    \draw[very thick] (A) -- (B) (C) -- (D);
    \draw[thick,densely dashed] (D) -- (C1);
    \node[above left] at (A) {\(A\)};
    \node[below right] at (B) {\(B\)};
    \node[above left] at (C) {\(C\)};
    \node[below right] at (D) {\(D\)};
    \node[above right] at (C1) {\(C_1\)};
    \node at (-2.18,-1.05) {\(P\)};
\end{tikzpicture}

    \caption{Прямые \(AB\) и \(CD\), перпендикулярные к одной плоскости \(P\)}\label{fig:kiselev-stereo-317}
\end{figure}

Предположим противное, то есть что прямые \(AB\) и \(CD\) не параллельны. Проведём тогда через точку \(D\) прямую, параллельную \(AB\). При нашем предположении это будет какая-нибудь прямая \(DC_1\), не совпадающая с \(DC\). Согласно \hyperref[para:plane-perpendicular-to-parallel-lines]{прямой теореме} прямая \(DC_1\) будет перпендикулярна к плоскости \(P\). Тогда у нас будут два перпендикуляра к плоскости \(P\), проходящие через одну и ту же точку, \(DC\) и \(DC_1\). Так как это невозможно, то нельзя допустить, чтобы прямые \(AB\) и \(CD\) были непараллельны.

\kiselevparagraph{para:lines-parallel-to-same-line} \textbf{Теорема. Если две прямые} \(A\) и \(B\) (рис.~\ref{fig:kiselev-stereo-318}) \textbf{параллельны третьей прямой} \(C\)\textbf{, то они параллельны между собой.}

\begin{figure}[htbp]
    \centering
    \renewcommand{\thefigure}{318}
    \captionsetup{name=Рис.,labelsep=period}
    \begin{tikzpicture}[line cap=round,line join=round]
    \draw[thick] (-2.95,-1.62) -- (-1.95,0.96) -- (3.1,0.96);
    \draw[very thick] (-2.95,-1.62) -- (2.1,-1.62) -- (3.1,0.96);
    \draw[thick] (-1.4,-0.83) -- (-1.4,2.64) node[above right] {\(A\)};
    \draw[thick] (0.04,-0.18) -- (0.04,3.29) node[above right] {\(B\)};
    \draw[thick] (1.48,-0.59) -- (1.48,2.88) node[above right] {\(C\)};
    \node at (-2.31,-1.13) {\(P\)};
\end{tikzpicture}

    \caption{Прямые \(A\) и \(B\), параллельные третьей прямой \(C\)}\label{fig:kiselev-stereo-318}
\end{figure}

Проведём плоскость \(P\), перпендикулярную к \(C\). Тогда \(A\) и \(B\) будут перпендикулярны к этой плоскости (\hyperref[para:plane-perpendicular-to-parallel-lines]{см.~§~\ref*{para:plane-perpendicular-to-parallel-lines}}) и, следовательно, \(A\parallel B\) (\hyperref[para:lines-normal-to-same-plane]{см.~§~\ref*{para:lines-normal-to-same-plane}}).

\kiselevparagraph{para:plane-intersection-parallel-to-lines} \textbf{Теорема. Если две прямые} \(A\) и \(B\) (рис.~\ref{fig:kiselev-stereo-319}) \textbf{параллельны и через каждую из них проходит плоскость, то линия пересечения этих плоскостей} (если она существует) \textbf{параллельна первым прямым.}

\begin{figure}[htbp]
    \centering
    \renewcommand{\thefigure}{319}
    \captionsetup{name=Рис.,labelsep=period}
    \begin{tikzpicture}[line cap=round,line join=round]
    \coordinate (C) at (-1.91,1.85);
    \coordinate (D) at (-1.91,-1.6);
    \coordinate (E) at (-1.91,1.1);
    \draw[thick] (D) -- (C) -- (1.45,2.26) -- (1.45,-1.19) -- (0.45,-1.312);
    \draw[thick,densely dashed] (0.45,-1.312) -- (D);
    \draw[thick] (C) -- (0.45,0.87) -- (0.45,-2.58) -- (D);
    \draw[very thick] (1.05,1.64) -- (1.05,-0.94);
    \draw[very thick] (0.11,0.44) -- (0.11,-2.14);
    \draw[thick,densely dashed] (E) -- (-1.49,-1.27) node[above right] {\(A_1\)};
    \node[above left] at (C) {\(C\)};
    \node[below left] at (D) {\(D\)};
    \node[left] at (E) {\(E\)};
    \node[above left] at (1.05,1.64) {\(A\)};
    \node[above left] at (0.11,0.44) {\(B\)};
    \node[above] at (0.5,2.14) {\(P\)};
    \node[above] at (-0.52,1.27) {\(Q\)};
\end{tikzpicture}

    \caption{Линия пересечения плоскостей параллельна прямым \(A\) и \(B\)}\label{fig:kiselev-stereo-319}
\end{figure}

% Печатная с. 286.
Пусть через прямую \(A\) проходит плоскость \(P\) и через прямую \(B\) --- плоскость \(Q\), и пусть эти плоскости пересекаются по прямой \(CD\). Требуется доказать, что \(CD\parallel A\) и \(CD\parallel B\). Через какую-нибудь точку \(E\) линии пересечения проведём прямую \(A_1\), параллельную \(A\). Тогда эта прямая должна быть также параллельна \(B\). Так как прямая \(A_1\) параллельна \(A\) и проходит через точку \(E\), то \(A_1\) должна лежать в плоскости, содержащей прямую \(A\) и точку \(E\), то есть в плоскости \(P\). С другой стороны, так как прямая \(A_1\) параллельна \(B\) и проходит через точку \(E\), то \(A_1\) должна лежать в плоскости \(Q\). Если же прямая \(A_1\) лежит одновременно в плоскостях \(P\) и \(Q\), то она должна совпадать с линией пересечения \(CD\). Значит, \(CD\parallel A\) и \(CD\parallel B\).


\section*{Прямая и плоскость, параллельные между собой}
\addcontentsline{toc}{section}{Прямая и плоскость, параллельные между собой}

\kiselevparagraph{para:parallel-line-plane-definition} \textbf{Определение.} Плоскость и прямая, не лежащая в этой плоскости, называются \emph{параллельными}, если они \emph{не пересекаются}, сколько бы их ни продолжали.

Следующие две теоремы выражают \emph{признаки} параллельности прямой с плоскостью.

\kiselevparagraph{para:line-plane-parallel-criteria} \textbf{Теорема 1. Если прямая} \(AB\) (рис.~\ref{fig:kiselev-stereo-320}) \textbf{и плоскость} \(P\) \textbf{перпендикулярны к одной и той же прямой} \(AC\)\textbf{, то они параллельны.}

\begin{figure}[htbp]
    \centering
    \renewcommand{\thefigure}{320}
    \captionsetup{name=Рис.,labelsep=period}
    \begin{tikzpicture}[line cap=round,line join=round]
    \coordinate (A) at (-1.83,1.9);
    \coordinate (B) at (1.02,1.9);
    \coordinate (C) at (-1.83,-0.31);
    \coordinate (M) at (2,-0.31);
    \draw[thick] (-2.95,-1.38) -- (-1.92,1.0) -- (3.22,1.0);
    \draw[very thick] (-2.95,-1.38) -- (2.19,-1.38) -- (3.22,1.0);
    \draw[thick] (-1.83,2.63) -- (C) (-1.83,-1.38) -- (-1.83,-1.93);
    \draw[thick,densely dashed] (C) -- (-1.83,-1.38) (C) -- (M) (B) -- (M);
    \draw[very thick] (A) -- (B);
    \node[left] at (A) {\(A\)};
    \node[above right] at (B) {\(B\)};
    \node[below right] at (C) {\(C\)};
    \node[above right] at (M) {\(M\)};
    \node[left] at (-2.78,-0.97) {\(P\)};
\end{tikzpicture}

    \caption{Прямая \(AB\) и плоскость \(P\), перпендикулярные к прямой \(AC\)}\label{fig:kiselev-stereo-320}
\end{figure}

Предположим, что прямая \(AB\) пересекается с плоскостью \(P\) в некоторой точке \(M\). Тогда, соединив прямой эту точку с точкой, в которой плоскость \(P\) пересекается с данной прямой, мы получим два перпендикуляра \(MC\) и \(MA\) на прямую \(AC\) из одной точки \(M\), что невозможно. Значит, \(AB\) не пересекается с \(P\), то есть \(AB\) параллельна \(P\).

% Печатная с. 287.
\textbf{Теорема 2. Если прямая} \(AB\) (рис.~\ref{fig:kiselev-stereo-321}) \textbf{параллельна какой-нибудь прямой} \(CD\)\textbf{, проведённой в плоскости} \(P\)\textbf{, то она параллельна самой плоскости.}

\begin{figure}[htbp]
    \centering
    \renewcommand{\thefigure}{321}
    \captionsetup{name=Рис.,labelsep=period}
    \begin{tikzpicture}[line cap=round,line join=round]
    \coordinate (A) at (-1.28,2.92);
    \coordinate (B) at (2.17,3.02);
    \coordinate (C) at (-1.98,-0.2);
    \coordinate (D) at (1.47,-0.1);
    \draw[thick] (-2.8,-1.24) -- (-1.88,1.04) -- (-1.72,1.044)
        (1.72,1.111) -- (2.86,1.14);
    \draw[thick,densely dashed] (-1.72,1.044) -- (1.72,1.111);
    \draw[very thick] (-2.8,-1.24) -- (1.94,-1.14) -- (2.86,1.14);
    \draw[thick] (A) -- (C) (B) -- (D) (-2.27,-0.208) -- (1.76,-0.092);
    \draw[very thick] (-1.59,2.911) -- (2.49,3.029);
    \node[above] at (A) {\(A\)};
    \node[above] at (B) {\(B\)};
    \node[below] at (C) {\(C\)};
    \node[below] at (D) {\(D\)};
    \node[left] at (-2.47,-0.42) {\(P\)};
    \node at (0.27,2.08) {\(R\)};
\end{tikzpicture}

    \caption{Прямая \(AB\), параллельная прямой \(CD\) в плоскости \(P\)}\label{fig:kiselev-stereo-321}
\end{figure}

Проведём через \(AB\) и \(CD\) плоскость \(R\) и предположим, что прямая \(AB\) где-нибудь пересекается с плоскостью \(P\). Тогда точка пересечения, находясь на прямой \(AB\), должна принадлежать также и плоскости \(R\), в которой лежит \(AB\). В то же время точка пересечения, конечно, должна принадлежать и плоскости \(P\). Значит, точка пересечения, находясь одновременно и в плоскости \(R\), и в плоскости \(P\), должна лежать на прямой \(CD\), по которой пересекаются эти плоскости. Следовательно, прямая \(AB\) пересекается с прямой \(CD\). Но это невозможно, так как, по условию, \(AB\parallel CD\). Значит, нельзя допустить, чтобы прямая \(AB\) пересекалась с плоскостью \(P\), и потому \(AB\parallel P\).

Следующие теоремы выражают некоторые \emph{свойства прямой и плоскости, параллельных между собой}.

\kiselevparagraph{para:parallel-line-plane-section} \textbf{Теорема. Если плоскость} \(R\) (рис.~\ref{fig:kiselev-stereo-321}) \textbf{проходит через прямую} \(AB\)\textbf{, параллельную другой плоскости} \(P\)\textbf{, и пересекает эту плоскость, то линия пересечения} \(CD\) \textbf{параллельна первой прямой} \(AB\)\textbf{.}

Действительно, во-первых, прямая \(CD\) лежит в одной плоскости с \(AB\). Во-вторых, эта прямая не может пересечься с \(AB\), потому что в противном случае прямая \(AB\) пересекалась бы с плоскостью \(P\), что невозможно.

\kiselevparagraph{para:line-parallel-to-two-intersecting-planes} \textbf{Следствие. Если прямая} \(AB\) (рис.~\ref{fig:kiselev-stereo-322}) \textbf{параллельна двум пересекающимся плоскостям} \(P\) и \(Q\)\textbf{, то она параллельна линии их пересечения} \(CD\)\textbf{.}

\begin{figure}[htbp]
    \centering
    \renewcommand{\thefigure}{322}
    \captionsetup{name=Рис.,labelsep=period}
    \begin{tikzpicture}[line cap=round,line join=round]
    \coordinate (C) at (-2.1,1.71);
    \coordinate (D) at (-2.1,-1.63);
    \coordinate (M) at (-2.1,0.38);
    \draw[thick] (D) -- (C) -- (0.72,2.13) -- (0.72,-1.21) -- (0.33,-1.268);
    \draw[thick,densely dashed] (0.33,-1.268) -- (D);
    \draw[thick] (C) -- (0.33,0.89) -- (0.33,-2.45) -- (D);
    \draw[very thick] (1.7,1.54) node[above right] {\(A\)}
        -- (1.7,-1.98) node[below right] {\(B\)};
    \fill (M) circle[radius=1.1pt];
    \node[left] at (C) {\(C\)};
    \node[left] at (M) {\(M\)};
    \node[left] at (D) {\(D\)};
    \node at (0.35,1.65) {\(P\)};
    \node at (-0.16,0.69) {\(Q\)};
\end{tikzpicture}

    \caption{Прямая \(AB\), параллельная пересекающимся плоскостям \(P\) и \(Q\)}\label{fig:kiselev-stereo-322}
\end{figure}

Представим плоскость, проходящую через \(AB\) и какую-нибудь точку \(M\) прямой \(CD\). Эта плоскость должна пересечься с \(P\) и \(Q\) по прямым, параллельным \(AB\) и проходящим через \(M\).
% Печатная с. 288.
Но через \(M\) можно провести только одну прямую, параллельную \(AB\). Значит, две линии пересечения этой плоскости с плоскостями \(P\) и \(Q\) должны слиться в одну прямую. Эта прямая, находясь одновременно и в плоскости \(P\), и в плоскости \(Q\), должна совпадать с прямой \(CD\), по которой плоскости \(P\) и \(Q\) пересекаются. Значит, \(CD\parallel AB\).

\kiselevparagraph{para:distance-from-parallel-line-to-plane} \textbf{Теорема. Все точки прямой} \(AB\) (рис.~\ref{fig:kiselev-stereo-323})\textbf{, параллельной плоскости} \(P\)\textbf{, одинаково удалены от этой плоскости.}

\begin{figure}[htbp]
    \centering
    \renewcommand{\thefigure}{323}
    \captionsetup{name=Рис.,labelsep=period}
    \begin{tikzpicture}[line cap=round,line join=round]
    \coordinate (M) at (-1.27,2.85);
    \coordinate (N) at (1.49,2.85);
    \coordinate (C) at (-1.27,-0.24);
    \coordinate (D) at (1.49,-0.24);
    \draw[thick] (-2.85,-1.12) -- (-2.05,0.87) -- (3.1,0.87);
    \draw[very thick] (-2.85,-1.12) -- (2.3,-1.12) -- (3.1,0.87);
    \draw[thick] (-2.495,-0.24) -- (2.655,-0.24) (M) -- (C) (N) -- (D);
    \draw[very thick] (-2.73,2.85) node[left] {\(A\)} -- (2.73,2.85) node[right] {\(B\)};
    \node[above] at (M) {\(M\)};
    \node[above] at (N) {\(N\)};
    \node[below] at (C) {\(C\)};
    \node[below] at (D) {\(D\)};
    \node at (-2.31,-0.74) {\(P\)};
\end{tikzpicture}

    \caption{Точки прямой \(AB\) одинаково удалены от плоскости \(P\)}\label{fig:kiselev-stereo-323}
\end{figure}

Из каких-нибудь точек \(M\) и \(N\) прямой \(AB\) опустим на \(P\) перпендикуляры \(MC\) и \(ND\). Так как эти перпендикуляры параллельны (\hyperref[para:lines-normal-to-same-plane]{см.~§~\ref*{para:lines-normal-to-same-plane}}), то через них можно провести плоскость. Эта плоскость пересекается с \(P\) по прямой \(CD\), параллельной \(AB\) (\hyperref[para:parallel-line-plane-section]{см.~§~\ref*{para:parallel-line-plane-section}}). Поэтому фигура \(MNDC\) --- параллелограмм и, следовательно, \(MC=ND\). Но длину перпендикуляра, опущенного из точки на плоскость, принимают за меру расстояния этой точки от плоскости. Значит, любые точки \(M\) и \(N\) прямой \(AB\) одинаково удалены от плоскости \(P\).


\section*{Параллельные плоскости}
\addcontentsline{toc}{section}{Параллельные плоскости}

\kiselevparagraph{para:parallel-planes-definition} \textbf{Определение.} Две плоскости называются \emph{параллельными}, если они \emph{не пересекаются, сколько бы их ни продолжали}.

Следующие две теоремы выражают \emph{признаки} параллельности двух плоскостей.

\kiselevparagraph{para:parallel-planes-criteria} \textbf{Теорема 1. Если две плоскости} \(P\) и \(Q\) (рис.~\ref{fig:kiselev-stereo-324}) \textbf{перпендикулярны к одной и той же прямой} \(AB\)\textbf{, то они параллельны.}

\begin{figure}[htbp]
    \centering
    \renewcommand{\thefigure}{324}
    \captionsetup{name=Рис.,labelsep=period}
    \begin{tikzpicture}[line cap=round,line join=round]
    \foreach \height/\plane in {0.6/P,-2.7/Q}{
        \begin{scope}[yshift=\height cm]
            \draw[thick] (-2.62,0) -- (-1.73,2) -- (2.58,2);
            \draw[very thick] (-2.62,0) -- (1.69,0) -- (2.58,2);
            \node at (1.35,1.36) {\(\plane\)};
        \end{scope}
    }
    \draw[thick] (-0.53,3.68) node[above right] {\(A\)} -- (-0.53,1.48)
        (-0.53,0.6) -- (-0.53,-1.82) (-0.53,-2.7) -- (-0.53,-3.8) node[below right] {\(B\)};
    \draw[thick,densely dashed] (-0.53,1.48) -- (-0.53,0.6)
        (-0.53,-1.82) -- (-0.53,-2.7);
    \fill (-0.53,1.48) circle[radius=1.35pt];
    \fill (-0.53,-1.82) circle[radius=1.35pt];
\end{tikzpicture}

    \caption{Плоскости \(P\) и \(Q\), перпендикулярные к одной прямой \(AB\)}\label{fig:kiselev-stereo-324}
\end{figure}

Действительно, если бы плоскости \(P\) и \(Q\) пересекались, то через любую точку их пересечения проходили бы две плоскости \(P\) и \(Q\), перпендикулярные к прямой \(AB\), что невозможно.

% Печатная с. 289.
\phantomsection\label{thm:parallel-planes-two-lines}
\textbf{Теорема 2. Если две пересекающиеся прямые} \(AB\) и \(AC\) (рис.~\ref{fig:kiselev-stereo-325}) \textbf{одной плоскости} \(P\) \textbf{соответственно параллельны двум прямым} \(A_1B_1\) и \(A_1C_1\) \textbf{другой плоскости} \(Q\)\textbf{, то эти плоскости параллельны.}

\begin{figure}[htbp]
    \centering
    \renewcommand{\thefigure}{325}
    \captionsetup{name=Рис.,labelsep=period}
    \begin{tikzpicture}[line cap=round,line join=round]
    \coordinate (A) at (-1.75,1.69);
    \coordinate (A11) at (-1.75,-2.15);
    \coordinate (A1) at (0.38,-2.15);
    \foreach \height/\plane in {0.74/P,-3.1/Q}{
        \begin{scope}[yshift=\height cm]
            \draw[thick] (-3.27,0) -- (-2.43,2.04) -- (3.3,2.04);
            \draw[very thick] (-3.27,0) -- (2.46,0) -- (3.3,2.04);
            \node[left] at (-2.89,0.92) {\(\plane\)};
        \end{scope}
    }
    \draw[thick,densely dashed] (A) -- (-1.75,0.74);
    \draw[thick] (-1.75,0.74) -- (A11);
    \draw[very thick] (-0.48,2.26) node[right] {\(C\)} -- (A) -- (-0.48,1.12) node[right] {\(B\)};
    \draw[thick] (-0.48,-1.58) node[right] {\(C_{11}\)} -- (A11)
        -- (-0.48,-2.72) node[right] {\(B_{11}\)};
    \draw[very thick] (1.65,-1.58) node[right] {\(C_1\)} -- (A1)
        -- (1.65,-2.72) node[right] {\(B_1\)};
    \node[left] at (A) {\(A\)};
    \node[below left] at (A11) {\(A_{11}\)};
    \node[left] at (A1) {\(A_1\)};
\end{tikzpicture}

    \caption{Соответственно параллельные прямые в плоскостях \(P\) и \(Q\)}\label{fig:kiselev-stereo-325}
\end{figure}

Из точки \(A\) опустим на плоскость \(Q\) перпендикуляр \(AA_{11}\) и проведём прямые \(A_{11}B_{11}\) и \(A_{11}C_{11}\), соответственно параллельные прямым \(A_1B_1\) и \(A_1C_1\). Эти прямые также параллельны прямым \(AB\) и \(AC\) (\hyperref[para:lines-parallel-to-same-line]{см.~§~\ref*{para:lines-parallel-to-same-line}}). Так как \(AA_{11}\perp A_{11}B_{11}\) и \(AB\parallel A_{11}B_{11}\), то \(AA_{11}\perp AB\). Так же доказывается, что \(AA_{11}\perp AC\). Следовательно, прямая \(AA_{11}\) перпендикулярна к плоскости \(P\) (\hyperref[para:perpendicular-plane-criterion]{см.~§~\ref*{para:perpendicular-plane-criterion}}). Таким образом, плоскости \(P\) и \(Q\) перпендикулярны к одной и той же прямой \(AA_{11}\) и потому параллельны.

Следующие теоремы выражают некоторые \emph{свойства параллельных плоскостей}.

\kiselevparagraph{para:sections-of-parallel-planes} \textbf{Теорема. Если две параллельные плоскости} \(P\) и \(Q\) (рис.~\ref{fig:kiselev-stereo-326}) \textbf{пересекаются третьей плоскостью} \(R\)\textbf{, то линии пересечения} \(AC\) и \(BD\) \textbf{параллельны.}

\begin{figure}[htbp]
    \centering
    \renewcommand{\thefigure}{326}
    \captionsetup{name=Рис.,labelsep=period}
    \begin{tikzpicture}[line cap=round,line join=round]
    \coordinate (A) at (-1.32,1.45);
    \coordinate (C) at (1.04,2.19);
    \coordinate (B) at (-0.69,-2.09);
    \coordinate (D) at (1.67,-1.35);
    \path (A) -- (B) coordinate[pos=0.2288] (frontA);
    \path (C) -- (D) coordinate[pos=0.4379] (frontC);
    \foreach \height/\plane in {0.64/P,-3.04/Q}{
        \begin{scope}[yshift=\height cm]
            \draw[thick] (-3.12,0) -- (-2.25,2.22) -- (3.07,2.22);
            \draw[very thick] (-3.12,0) -- (2.2,0) -- (3.07,2.22);
            \node[left] at (-2.68,1.14) {\(\plane\)};
        \end{scope}
    }
    \draw[thick] (A) -- (C) (frontA) -- (B) -- (D) -- (frontC);
    \draw[thick,densely dashed] (A) -- (frontA) (C) -- (frontC);
    \node[left] at (A) {\(A\)};
    \node[above right] at (C) {\(C\)};
    \node[below left] at (B) {\(B\)};
    \node[right] at (D) {\(D\)};
    \node at (0.31,-0.35) {\(R\)};
\end{tikzpicture}

    \caption{Линии пересечения параллельных плоскостей \(P\) и \(Q\) с плоскостью \(R\)}\label{fig:kiselev-stereo-326}
\end{figure}

Действительно, во-первых, прямые \(AC\) и \(BD\) находятся в одной плоскости \(R\). Во-вторых, они не могут пересечься, так как в противном случае пересекались бы плоскости \(P\) и \(Q\), что противоречит условию.

\kiselevparagraph{para:transversals-of-parallel-planes} \textbf{Теоремы.}

\begin{enumerate}[label=\arabic*),ref=\arabic*]
    \item \textbf{Если прямая пересекает одну из параллельных плоскостей, то она пересекает и другую.}
    \item \textbf{Если плоскость пересекает одну из параллельных плоскостей, то она пересекает и другую.}
\end{enumerate}

\begin{figure}[htbp]
    \centering
    \renewcommand{\thefigure}{327}
    \captionsetup{name=Рис.,labelsep=period}
    \begin{tikzpicture}[line cap=round,line join=round]
    \coordinate (B) at (-0.72,1.47);
    \coordinate (C) at (-0.72,-2.12);
    \foreach \height/\plane in {0.56/P,-3.03/Q}{
        \begin{scope}[yshift=\height cm]
            \draw[thick] (-2.91,0) -- (-2.02,2.13) -- (2.83,2.13);
            \draw[very thick] (-2.91,0) -- (1.94,0) -- (2.83,2.13);
            \node at (2.21,1.69) {\(\plane\)};
        \end{scope}
    }
    \draw[thick] (-1.98,3.8) node[above left] {\(A\)} -- (B);
    \draw[thick,densely dashed] (B) -- (-0.72,0.56);
    \draw[thick] (-0.72,0.56) -- (C);
    \draw[thick] (-1.89,1.47) node[left] {\(D\)} -- (1.4,1.47) node[right] {\(E\)};
    \draw[thick] (-1.89,-2.12) node[left] {\(F\)} -- (1.4,-2.12) node[right] {\(H\)};
    \node[above right] at (B) {\(B\)};
    \node[above right] at (C) {\(C\)};
\end{tikzpicture}

    \caption{Пересечение параллельных плоскостей \(P\) и \(Q\)}\label{fig:kiselev-stereo-327}
\end{figure}

\begin{enumerate}[label=\arabic*),ref=\arabic*]
    \item\label{item:line-intersects-both-parallel-planes} Пусть прямая \(AB\) (рис.~\ref{fig:kiselev-stereo-327}) пересекает в точке \(B\) плоскость \(P\), параллельную \(Q\). Соединим прямой точку \(B\) с какой-нибудь точкой плоскости \(Q\)
    % Печатная с. 290.
    и через \(AB\) и \(BC\) проведём плоскость, не указанную на рисунке. Эта плоскость, содержащая точки \(B\) и \(C\), пересекается с \(P\) и \(Q\) по некоторым прямым \(DE\) и \(FH\), которые параллельны (\hyperref[para:sections-of-parallel-planes]{см.~§~\ref*{para:sections-of-parallel-planes}}). Прямая \(AB\) лежит в одной плоскости с \(DE\) и \(FH\) и пересекает одну из этих параллельных. Следовательно, как мы знаем из планиметрии (\kiselevPlanimetryPage{71}{см.~§~80, пункт~1}), она пересечёт и другую, а значит, пересечёт и плоскость \(Q\).
    \item Пусть какая-нибудь плоскость пересекает плоскость \(P\) (рис.~\ref{fig:kiselev-stereo-327}), параллельную \(Q\). Тогда на ней можно взять прямую \(AB\), которая тоже пересекает плоскость \(P\). \hyperref[item:line-intersects-both-parallel-planes]{По доказанному} эта прямая пересекает и плоскость \(Q\). Значит, с этой плоскостью пересечётся и та плоскость, в которой взята \(AB\).
\end{enumerate}

\kiselevparagraph{para:line-normal-to-parallel-planes} \textbf{Теорема. Если прямая} \(AB\) (рис.~\ref{fig:kiselev-stereo-328}) \textbf{перпендикулярна к одной из параллельных плоскостей} \(P\)\textbf{, то она перпендикулярна и к~другой}~\(Q\)\textbf{.}

\begin{figure}[htbp]
    \centering
    \renewcommand{\thefigure}{328}
    \captionsetup{name=Рис.,labelsep=period}
    \begin{tikzpicture}[line cap=round,line join=round]
    \coordinate (A) at (-1.47,3.53);
    \coordinate (B) at (-1.47,1.15);
    \coordinate (B1) at (-1.47,-2.27);
    \coordinate (C) at (1.53,1.7);
    \coordinate (C1) at (1.53,-1.72);
    \coordinate (D) at (0.76,0.54);
    \coordinate (D1) at (0.76,-2.88);
    \foreach \height/\plane in {0.06/P,-3.36/Q}{
        \begin{scope}[yshift=\height cm]
            \draw[thick] (-3.34,0) -- (-2.47,2.2) -- (-1.47,2.2);
            \draw[thick,densely dashed] (-1.47,2.2) -- (1.53,2.2);
            \draw[thick] (1.53,2.2) -- (3.21,2.2);
            \draw[very thick] (-3.34,0) -- (2.34,0) -- (3.21,2.2);
            \node at (-2.84,0.45) {\(\plane\)};
        \end{scope}
    }
    \draw[thick] (A) -- (1.53,4.14) -- (C)
        (A) -- (0.76,3.05) -- (D) -- (B) -- (A);
    \draw[thick,densely dashed] (B) -- (C)
        (B) -- (-1.47,0.06) (C) -- (1.53,0.06) (D) -- (0.76,0.06)
        (B1) -- (C1);
    \draw[thick] (-1.47,0.06) -- (B1) -- (D1) -- (0.76,0.06)
        (1.53,0.06) -- (C1);
    \node[above left] at (A) {\(A\)};
    \node[left] at (B) {\(B\)};
    \node[below left] at (B1) {\(B_1\)};
    \node[right] at (C) {\(C\)};
    \node[right] at (C1) {\(C_1\)};
    \node[right,yshift=3pt] at (D) {\(D\)};
    \node[right] at (D1) {\(D_1\)};
    \node at (1.02,3.72) {\(M\)};
    \node at (0.23,2.84) {\(N\)};
\end{tikzpicture}

    \caption{Прямая \(AB\), перпендикулярная к параллельным плоскостям \(P\) и \(Q\)}\label{fig:kiselev-stereo-328}
\end{figure}

Прямая \(AB\), пересекая одну из параллельных плоскостей, пересечёт и другую в некоторой точке \(B_1\). Проведём через \(AB\) какие-нибудь две плоскости \(M\) и \(N\), которые пересекутся с \(P\) и \(Q\) по параллельным прямым (\hyperref[para:sections-of-parallel-planes]{см.~§~\ref*{para:sections-of-parallel-planes}}): одна по \(BC\) и \(B_1C_1\), другая по \(BD\) и \(B_1D_1\). Согласно условию, прямая \(AB\) перпендикулярна к \(BC\) и \(BD\). Следовательно, она также перпендикулярна к \(B_1C_1\) и \(B_1D_1\), а потому перпендикулярна и к плоскости \(Q\).

\kiselevparagraph{para:unique-parallel-plane-through-point} \textbf{Теорема. Через любую точку} \(B\) (рис.~\ref{fig:kiselev-stereo-328}) \textbf{пространства можно провести плоскость} \(P\)\textbf{, параллельную данной плоскости} \(Q\)\textbf{, и притом только одну.}

Всегда можно из точки \(B\) опустить на плоскость \(Q\) перпендикуляр \(BB_1\) и затем через точку \(B\) провести плоскость \(P\), перпендикулярную к \(BB_1\). Эта плоскость будет параллельна \(Q\) (\hyperref[para:parallel-planes-criteria]{см.~§~\ref*{para:parallel-planes-criteria}, теорема~1}).

% Печатная с. 291.
Другой плоскости, параллельной \(Q\), через точку \(B\) провести нельзя, так как, если бы это было возможно, то через точку \(B\) проходили бы две плоскости, перпендикулярные к \(BB_1\) (\hyperref[para:line-normal-to-parallel-planes]{см.~§~\ref*{para:line-normal-to-parallel-planes}}), что невозможно.

\kiselevparagraph{para:equal-segments-between-parallel-planes} \textbf{Теорема. Отрезки параллельных прямых} \(AC\) и \(BD\) (рис.~\ref{fig:kiselev-stereo-329})\textbf{, заключённые между параллельными плоскостями} \(P\) и \(Q\)\textbf{, равны.}

\begin{figure}[htbp]
    \centering
    \renewcommand{\thefigure}{329}
    \captionsetup{name=Рис.,labelsep=period}
    \begin{tikzpicture}[line cap=round,line join=round]
    \coordinate (A) at (-1.0,0.57);
    \coordinate (B) at (1.16,2.8);
    \coordinate (C) at (-1.0,-2.89);
    \coordinate (D) at (1.16,-0.66);
    \foreach \height/\plane in {0.57/P,-2.89/Q}{
        \begin{scope}[yshift=\height cm]
            \draw[thick] (-3.05,0) -- (-2.23,2.23) -- (3.02,2.23);
            \draw[very thick] (-3.05,0) -- (2.2,0) -- (3.02,2.23);
            \node at (-2.04,1.02) {\(\plane\)};
        \end{scope}
    }
    \draw[thick] (A) -- (B) (C) -- (D) (A) -- (C) (1.16,0.57) -- (D);
    \draw[thick,densely dashed] (B) -- (1.16,0.57);
    \node[above left] at (A) {\(A\)};
    \node[above] at (B) {\(B\)};
    \node[above left] at (C) {\(C\)};
    \node[below right] at (D) {\(D\)};
    \node at (0.11,-0.42) {\(R\)};
\end{tikzpicture}

    \caption{Равные отрезки \(AC\) и \(BD\) между параллельными плоскостями}\label{fig:kiselev-stereo-329}
\end{figure}

Через параллельные прямые \(AC\) и \(BD\) проведём плоскость \(R\). Она пересечёт \(P\) и \(Q\) по параллельным прямым \(AB\) и \(CD\). Следовательно, фигура \(ABCD\)\footnote{В издании 1914 года указан порядок \(ABCD\). На рис.~\ref{fig:kiselev-stereo-329} при обходе по сторонам вершины идут в порядке \(A,B,D,C\).} будет параллелограммом, и потому \(AC=BD\).

\kiselevparagraph{para:parallel-planes-equidistant} \textbf{Следствие. Параллельные плоскости везде одинаково удалены одна от другой}, потому что, если параллельные прямые \(AC\) и \(BD\) (рис.~\ref{fig:kiselev-stereo-329}) перпендикулярны к \(P\), то они также будут перпендикулярны к \(Q\) и в то же время равны.

\kiselevparagraph{para:angles-with-parallel-sides} \textbf{Теорема. Два угла} \(BAC\) и \(B_1A_1C_1\) (рис.~\ref{fig:kiselev-stereo-330}) \textbf{с соответственно параллельными и одинаково направленными сторонами равны и лежат в параллельных плоскостях} \(P\) и \(Q\)\textbf{.}

\begin{figure}[htbp]
    \centering
    \renewcommand{\thefigure}{330}
    \captionsetup{name=Рис.,labelsep=period}
    \begin{tikzpicture}[line cap=round,line join=round]
    \coordinate (A) at (-1.93,1.58);
    \coordinate (B) at (0.16,2.5);
    \coordinate (C) at (-0.2,1.0);
    \coordinate (A1) at (-0.84,-2.28);
    \coordinate (B1) at (1.25,-1.36);
    \coordinate (C1) at (0.89,-2.86);
    \foreach \height/\plane in {0.52/P,-3.34/Q}{
        \begin{scope}[yshift=\height cm]
            \draw[thick] (-3.08,0) -- (-2.25,2.5) -- (3.16,2.5);
            \draw[very thick] (-3.08,0) -- (2.33,0) -- (3.16,2.5);
            \node[left,xshift=-3pt] at (-2.18,2.17) {\(\plane\)};
        \end{scope}
    }
    \path (A) -- (A1) coordinate[pos=0.2746] (frontA);
    \path (B) -- (B1) coordinate[pos=0.5130] (frontB);
    \path (C) -- (C1) coordinate[pos=0.1244] (frontC);
    \draw[thick,densely dashed] (A) -- (frontA) (B) -- (frontB) (C) -- (frontC);
    \draw[thick] (frontA) -- (A1) (frontB) -- (B1) (frontC) -- (C1);
    \draw[very thick] (A) -- (B) -- (C) -- cycle (A1) -- (B1) -- (C1) -- cycle;
    \node[left] at (A) {\(A\)};
    \node[right] at (B) {\(B\)};
    \node[right] at (C) {\(C\)};
    \node[left] at (A1) {\(A_1\)};
    \node[right] at (B1) {\(B_1\)};
    \node[right] at (C1) {\(C_1\)};
\end{tikzpicture}

    \caption{Углы \(BAC\) и \(B_1A_1C_1\) с соответственно параллельными сторонами}\label{fig:kiselev-stereo-330}
\end{figure}

То, что плоскости \(P\) и \(Q\) параллельны, было доказано прежде (\hyperref[thm:parallel-planes-two-lines]{см.~§~\ref*{para:parallel-planes-criteria}, теорема~2}). Остаётся доказать, что углы \(A\) и \(A_1\) равны.

Отложим \(AB=A_1B_1\), \(AC=A_1C_1\) и проведём \(AA_1\), \(BB_1\), \(CC_1\), \(BC\) и \(B_1C_1\). Так как отрезки \(AB\) и \(A_1B_1\) равны и параллельны, то фигура \(ABB_1A_1\) является параллелограммом (\kiselevPlanimetryPage{85}{см.~Планиметрия, §~99, пункт~2}). Поэтому отрезки \(AA_1\) и \(BB_1\) равны и параллельны. По той же причине равны и параллельны отрезки \(AA_1\) и \(CC_1\). Следовательно, \(BB_1\parallel CC_1\) и \(BB_1=CC_1\). Поэтому \(BC=B_1C_1\) и \(\triangle ABC=\triangle A_1B_1C_1\) (по трём сторонам). Значит, \(\angle A=\angle A_1\).

% Печатная с. 292.
\kiselevparagraph{para:parallel-planes-proportional-segments} \textbf{Теорема. Если две прямые} \(AB\) и \(CD\) (рис.~\ref{fig:kiselev-stereo-331}) \textbf{пересекаются рядом параллельных плоскостей} \(P,Q,R,\ldots\)\textbf{, то эти плоскости делят их на пропорциональные части.}

\begin{figure}[htbp]
    \centering
    \renewcommand{\thefigure}{331}
    \captionsetup{name=Рис.,labelsep=period}
    \begin{tikzpicture}[line cap=round,line join=round,scale=0.82]
    \coordinate (E) at (0.05,4.2);
    \coordinate (F) at (0.7325,0.95);
    \coordinate (H) at (1.436,-2.4);
    \coordinate (K) at (-1.95,4.2);
    \coordinate (L) at (-2.8275,0.95);
    \coordinate (M) at (-3.732,-2.4);
    \coordinate (L1) at (-0.8275,0.95);
    \coordinate (M1) at (-1.732,-2.4);
    \foreach \height/\plane in {3.17/P,-0.08/Q,-3.43/R}{
        \begin{scope}[yshift=\height cm]
            \draw[thick] (-4.94,0) -- (-3.98,2.05) -- (3.05,2.05);
            \draw[very thick] (-4.94,0) -- (2.09,0) -- (3.05,2.05);
            \node[right] at (2.6,1.1) {\(\plane\)};
        \end{scope}
    }
    \path (E) -- (F) coordinate[pos=0.3169] (frontE);
    \path (F) -- (H) coordinate[pos=0.3075] (frontF);
    \path (H) -- (1.8056,-4.16) coordinate[pos=0.5852] (frontH);
    \path (K) -- (L) coordinate[pos=0.3169] (frontK);
    \path (L) -- (M) coordinate[pos=0.3075] (frontL);
    \path (M) -- (-4.2072,-4.16) coordinate[pos=0.5852] (frontM);
    \path (E) -- (L1) coordinate[pos=0.3169] (frontE1);
    \path (L1) -- (M1) coordinate[pos=0.3075] (frontL1);
    \path (M1) -- (-2.2072,-4.16) coordinate[pos=0.5852] (frontM1);

    \draw[thick,densely dashed] (E) -- (frontE) (F) -- (frontF) (H) -- (frontH)
        (K) -- (frontK) (L) -- (frontL) (M) -- (frontM)
        (E) -- (frontE1) (L1) -- (frontL1) (M1) -- (frontM1);
    \draw[thick] (-0.2755,5.75) node[above left] {\(A\)} -- (E)
        (frontE) -- (F) (frontF) -- (H)
        (frontH) -- (1.8056,-4.16) node[below right] {\(B\)};
    \draw[thick] (-1.5315,5.75) node[above left] {\(C\)} -- (K)
        (frontK) -- (L) (frontL) -- (M)
        (frontM) -- (-4.2072,-4.16) node[below] {\(D\)};
    \draw[thick] (0.4685,5.75) node[above right] {\(C_1\)} -- (E)
        (frontE1) -- (L1) (frontL1) -- (M1)
        (frontM1) -- (-2.2072,-4.16) node[below] {\(D_1\)};
    \draw[thick] (L1) -- (F) (M1) -- (H);
    \node[right] at (E) {\(E\)};
    \node[right] at (F) {\(F\)};
    \node[right] at (H) {\(H\)};
    \node[left] at (K) {\(K\)};
    \node[left] at (L) {\(L\)};
    \node[right] at (M) {\(M\)};
    \node[left] at (L1) {\(L_1\)};
    \node[left] at (M1) {\(M_1\)};
\end{tikzpicture}

    \caption{Пропорциональные отрезки между параллельными плоскостями}\label{fig:kiselev-stereo-331}
\end{figure}

Пусть \(E,F,H,\ldots\) --- точки пересечения прямой \(AB\) с плоскостями \(P,Q,R,\ldots\), а \(K,L,M,\ldots\) --- точки пересечения другой прямой \(CD\) с теми же плоскостями. Требуется доказать, что части \(EF,FH,\ldots\) одной прямой пропорциональны частям \(KL,LM,\ldots\) другой прямой, то есть
\[
    \frac{EF}{KL}=\frac{FH}{LM}=\cdots.
\]

Проведём через точку \(E\) прямую \(C_1D_1\), параллельную \(CD\), и представим плоскость, проходящую через две пересекающиеся прямые \(AB\) и \(C_1D_1\). Прямые \(L_1F,M_1H,\ldots\), по которым эта плоскость пересекается с данными плоскостями, должны быть параллельны между собой (\hyperref[para:sections-of-parallel-planes]{см.~§~\ref*{para:sections-of-parallel-planes}}). Поэтому (\kiselevPlanimetryPage{171}{см.~Планиметрия, §~219})\footnote{В первом знаменателе этой формулы в издании 1914 года напечатано \(E_1L_1\). На рис.~\ref{fig:kiselev-stereo-331} отмечена точка \(E\), а далее в доказательстве используется отрезок \(EL_1\).}
\[
    \frac{EF}{E_1L_1}=\frac{FH}{L_1M_1}=\cdots.
\]

Но \(EL_1=KL\), \(L_1M_1=LM,\ldots\) (\hyperref[para:equal-segments-between-parallel-planes]{см.~§~\ref*{para:equal-segments-between-parallel-planes}}). Подставив в полученный ряд равных отношений отрезки \(KL,LM,\ldots\) вместо \(EL_1,L_1M_1,\ldots\), получим то, что требовалось доказать.

\clearpage
\endgroup
  % Глава III
    \chapter{Двугранные углы}\label{ch:ch4}

\begingroup
\setlength{\intextsep}{10pt plus 2pt minus 2pt}
\setlength{\textfloatsep}{12pt plus 2pt minus 2pt}

% А. П. Киселёв, «Элементарная геометрия», 23-е издание, 1914. Печатные с. 292–300.
% Страницы скана планиметрии на 11 больше печатных.
\newcommand*{\kiselevPlanimetryPage}[2]{%
    \href{https://commons.wikimedia.org/wiki/File:\%D0\%90._\%D0\%9F._\%D0\%9A\%D0\%B8\%D1\%81\%D0\%B5\%D0\%BB\%D1\%91\%D0\%B2._\%D0\%AD\%D0\%BB\%D0\%B5\%D0\%BC\%D0\%B5\%D0\%BD\%D1\%82\%D0\%B0\%D1\%80\%D0\%BD\%D0\%B0\%D1\%8F_\%D0\%B3\%D0\%B5\%D0\%BE\%D0\%BC\%D0\%B5\%D1\%82\%D1\%80\%D0\%B8\%D1\%8F_(1914).djvu?page=#1}{#2}%
}

% Печатная с. 292.
\kiselevparagraph{para:dihedral-angle-definition} \textbf{Определение.} Фигура, образованная двумя полуплоскостями \(P\) и \(Q\) (рис.~\ref{fig:kiselev-stereo-332}), исходящими из одной прямой \(AB\), вместе с частью пространства, ограниченной ими, называется \emph{двугранным углом}.

\begin{figure}[htbp]
    \centering
    \renewcommand{\thefigure}{332}
    \captionsetup{name=Рис.,labelsep=period}
    \begin{tikzpicture}[line cap=round,line join=round]
    \begin{scope}
        \coordinate (A) at (0,1.97);
        \coordinate (B) at (0,-1.87);
        \draw[thick] (B) -- (A) -- (2.59,1.97) -- (2.59,-1.87);
        \draw[thick,densely dashed] (B) -- (2.59,-1.87);
        \draw[very thick] (A) -- (2.13,1.03) -- (2.13,-2.81) -- (B);
        \node[above left] at (A) {\(A\)};
        \node[left] at (B) {\(B\)};
        \node[below left] at (2.59,1.97) {\(P\)};
        \node[below left] at (2.13,1.03) {\(Q\)};
    \end{scope}
    \begin{scope}[xshift=4.36cm]
        \coordinate (C) at (0,1.97);
        \coordinate (D) at (0,-1.87);
        \draw[thick] (D) -- (C);
        \foreach \x/\top/\bottom/\face in {3.6/2.27/-1.57/P,3.08/1.49/-2.35/Q,2.54/0.77/-3.07/R}{
            \draw[thick] (C) -- (\x,\top) -- (\x,\bottom);
            \draw[thick,densely dashed] (D) -- (\x,\bottom);
            \node[below left] at (\x,\top) {\(\face\)};
        }
        \draw[very thick] (C) -- (1.91,0.13) -- (1.91,-3.71) -- (D);
        \node[below left] at (1.91,0.13) {\(S\)};
        \node[above left] at (C) {\(C\)};
        \node[left] at (D) {\(D\)};
    \end{scope}
\end{tikzpicture}

    \caption{Двугранные углы с рёбрами \(AB\) и \(CD\)}\label{fig:kiselev-stereo-332}
\end{figure}

Прямая \(AB\) называется \emph{ребром}, а полуплоскости \(P\) и \(Q\) --- сторонами, или \emph{гранями}, двугранного угла. Такой угол обычно обозначают двумя буквами, поставленными у его ребра (двугранный угол \(AB\)).
% Печатная с. 293.
Но если при одном ребре лежат несколько двугранных углов, то каждый из них обозначают четырьмя буквами, из которых две средние стоят при ребре, а две крайние --- у граней (например, двугранный угол \(SCDR\)).

Если через ребро \(CD\) двугранного угла \(PCDS\) (рис.~\ref{fig:kiselev-stereo-332}) проведём внутри него, то есть в той части пространства, которая принадлежит углу, какие-нибудь полуплоскости \(R,Q,\ldots\), то образовавшиеся при этом двугранные углы \(RCDS,QCDR,\ldots\) рассматриваются как части первого двугранного угла.

Если из произвольной точки \(D\) ребра \(AB\) (рис.~\ref{fig:kiselev-stereo-333}) проведём на каждой грани по перпендикуляру к ребру, то образованный ими угол \(CDE\) называется \emph{линейным углом двугранного угла}. Величина линейного угла не зависит от положения точки \(D\) на ребре. Так, линейные углы \(CDE\) и \(C_1D_1E_1\) равны, потому что их стороны соответственно параллельны и одинаково направлены.

\begin{figure}[htbp]
    \centering
    \renewcommand{\thefigure}{333}
    \captionsetup{name=Рис.,labelsep=period}
    \begin{tikzpicture}[line cap=round,line join=round]
    \coordinate (A) at (0,2.55);
    \coordinate (B) at (0,-2.13);
    \coordinate (D) at (0,0.89);
    \coordinate (D1) at (0,-0.7);
    \coordinate (C) at (3.24,1.67);
    \coordinate (C1) at (3.24,0.08);
    \coordinate (E) at (2.17,0.12);
    \coordinate (E1) at (2.17,-1.47);

    \draw[thick] (B) -- (A) -- (3.24,3.33) -- (3.24,-1.35);
    \draw[thick,densely dashed] (B) -- (3.24,-1.35) (D) -- (C) (D1) -- (C1);
    \draw[very thick] (A) -- (2.17,1.78) -- (2.17,-2.90) -- (B);
    \draw[thick] (D) -- (E) (D1) -- (E1);
    \node[above left] at (A) {\(A\)};
    \node[below left] at (B) {\(B\)};
    \node[left] at (D) {\(D\)};
    \node[left] at (D1) {\(D_1\)};
    \node[right] at (C) {\(C\)};
    \node[right] at (C1) {\(C_1\)};
    \node[above right,yshift=1pt] at (E) {\(E\)};
    \node[above right,yshift=1pt] at (E1) {\(E_1\)};
    \node at (2.84,2.88) {\(P\)};
    \node at (1.53,1.59) {\(Q\)};
\end{tikzpicture}

    \caption{Линейные углы \(CDE\) и \(C_1D_1E_1\) двугранного угла}\label{fig:kiselev-stereo-333}
\end{figure}

Плоскость линейного угла перпендикулярна к ребру, так как она содержит две прямые, перпендикулярные к нему (\hyperref[para:perpendicular-plane-criterion]{см.~§~\ref*{para:perpendicular-plane-criterion}}).

\kiselevparagraph{para:dihedral-angles-comparison} \textbf{Равенство и неравенство двугранных углов.} Два двугранных угла считаются \emph{равными}, если при наложении они могут совместиться. В противном случае тот из углов считается меньшим, который может составить часть другого угла.

Можно рассматривать \emph{сумму, разность, произведение и частное двугранных углов} в том же смысле, что и для углов планиметрии. Подобно этим углам, двугранные углы могут быть \emph{смежными, прямыми, вертикальными} и т.~д.

\kiselevparagraph{para:dihedral-and-linear-angle-comparison} \textbf{Теоремы.}

\begin{enumerate}[label=\arabic*),ref=\arabic*]
    \item \textbf{Равным двугранным углам соответствуют равные линейные углы.}
    \item \textbf{Большему двугранному углу соответствует больший линейный угол.}
\end{enumerate}

% Печатная с. 294.
Пусть \(PABQ\) и \(P_1A_1B_1Q_1\) (рис.~\ref{fig:kiselev-stereo-334}) будут два двугранных угла. Наложим угол \(A_1B_1\) на угол \(AB\) так, чтобы ребро \(A_1B_1\) совпало с ребром \(AB\), а грань \(P_1\) --- с гранью \(P\). Тогда, если эти двугранные углы равны, грань \(Q_1\) совпадёт с \(Q\). Если же двугранные углы не равны, грань \(Q_1\) не совпадёт с \(Q\). Например, она займёт положение \(Q_{11}\), если угол \(A_1B_1\) меньше угла \(AB\).

\begin{figure}[htbp]
    \centering
    \renewcommand{\thefigure}{334}
    \captionsetup{name=Рис.,labelsep=period}
    \begin{tikzpicture}[line cap=round,line join=round]
    \coordinate (A) at (0,2.45);
    \coordinate (B) at (0,-0.9);
    \coordinate (C) at (3.36,-0.6);
    \coordinate (C11) at (2.92,-1.35);
    \coordinate (D) at (2.02,-2.11);
    \draw[thick] (-1.14,-2.75) -- (-0.27,-0.26) -- (0,-0.26)
        (3.36,-0.26) -- (4.10,-0.26);
    \draw[thick,densely dashed] (0,-0.26) -- (3.36,-0.26);
    \draw[very thick] (-1.14,-2.75) -- (3.19,-2.75) -- (4.10,-0.26);
    \draw[thick] (B) -- (A) -- (3.36,2.75) -- (C)
        (A) -- (2.92,2.00) -- (C11);
    \draw[thick,densely dashed] (B) -- (C) (B) -- (C11);
    \draw[very thick] (A) -- (2.02,1.24) -- (D) -- (B);
    \node[above left] at (A) {\(A\)};
    \node[below left] at (B) {\(B\)};
    \node[right] at (C) {\(C\)};
    \node[below,xshift=3pt] at (C11) {\(C_{11}\)};
    \node[right] at (D) {\(D\)};
    \node[below left] at (3.36,2.75) {\(Q\)};
    \node[below left] at (2.92,2.00) {\(Q_{11}\)};
    \node[below left] at (2.02,1.24) {\(P\)};
    \node at (-0.56,-2.34) {\(R\)};

    \begin{scope}[xshift=5.12cm]
        \coordinate (A1) at (0,2.45);
        \coordinate (B1) at (0,-0.9);
        \coordinate (C1) at (2.82,-0.62);
        \coordinate (D1) at (2.44,-1.93);
        \draw[thick] (B1) -- (A1) -- (2.82,2.73) -- (C1);
        \draw[thick,densely dashed] (B1) -- (C1);
        \draw[very thick] (A1) -- (2.44,1.42) -- (D1) -- (B1);
        \node[above left] at (A1) {\(A_1\)};
        \node[left] at (B1) {\(B_1\)};
        \node[right] at (C1) {\(C_1\)};
        \node[right] at (D1) {\(D_1\)};
        \node[below left] at (2.82,2.73) {\(Q_1\)};
        \node[below left] at (2.44,1.42) {\(P_1\)};
    \end{scope}
\end{tikzpicture}

    \caption{Наложение двугранных углов \(PABQ\) и \(P_1A_1B_1Q_1\)}\label{fig:kiselev-stereo-334}
\end{figure}

Заметив это, возьмём на общем ребре какую-нибудь точку \(B\) и проведём через неё плоскость \(R\), перпендикулярную к ребру. От пересечения этой плоскости с гранями двугранных углов получатся линейные углы. Ясно, что если двугранные углы совпадут, то у них окажется один и тот же линейный угол, а именно \(CBD\). Если же двугранные углы не совпадут, например, если грань \(Q_1\) займёт положение \(Q_{11}\), то у большего двугранного угла окажется больший линейный угол, а именно \(CBD>C_{11}BD\).

\needspace{3\baselineskip}
\kiselevparagraph{para:linear-and-dihedral-angle-comparison} \textbf{Обратные теоремы.}

\begin{enumerate}[label=\arabic*),ref=\arabic*]
    \item \textbf{Равным линейным углам соответствуют равные двугранные углы.}
    \item \textbf{Большему линейному углу соответствует больший двугранный угол.}
\end{enumerate}

Эти теоремы легко доказываются от противного (\kiselevPlanimetryPage{46}{см.~Планиметрия, §~51}).

\kiselevparagraph{para:superposition-of-spatial-figures} \textbf{Замечание.} Мы принимаем за очевидное, что \emph{наложение одной фигуры на другую}, часто употребляемое в стереометрии, всегда можно выполнить в такой последовательности:

\begin{enumerate}[label=\arabic*),ref=\arabic*]
    % Печатная с. 295 начинается со слова «фигур» в пункте 1.
    \item совмещаем какие-нибудь две точки фигур;
    \item совмещаем какие-нибудь две полупрямые, исходящие из совпавших точек;
    \item совмещаем какие-нибудь две полуплоскости, исходящие из совпавших прямых.
\end{enumerate}

Совместятся ли при этом другие элементы фигур, зависит от их свойств.

\kiselevparagraph{para:right-and-equal-dihedral-angles} \textbf{Следствия.}

\begin{enumerate}[label=\arabic*),ref=\arabic*]
    \item\label{item:right-dihedral-and-linear-angle} \textbf{Прямому двугранному углу соответствует прямой линейный угол и обратно.}

    Пусть двугранный угол \(PABQ\) (рис.~\ref{fig:kiselev-stereo-335}) прямой. Это значит, что он равен смежному углу \(QABP_1\). Но в таком случае линейные углы \(CDE\) и \(CDE_1\) также равны. А так как они смежные, то каждый из них должен быть прямым. Обратно, если равны смежные линейные углы \(CDE\) и \(CDE_1\), то равны и смежные двугранные углы, то есть каждый из них должен быть прямым.

    \item \textbf{Прямые двугранные углы равны}, потому что у них равны линейные углы. По той же причине:
    \item \textbf{вертикальные двугранные углы равны;}
    \item \textbf{двугранные углы с соответственно параллельными и одинаково направленными гранями равны.}
\end{enumerate}

\begin{figure}[htbp]
    \centering
    \renewcommand{\thefigure}{335}
    \captionsetup{name=Рис.,labelsep=period}
    \begin{tikzpicture}[line cap=round,line join=round]
    \coordinate (A) at (-0.97,-1.3);
    \coordinate (B) at (1.44,1.05);
    \path (A) -- (B) coordinate[pos=0.477] (D);
    \coordinate (C) at (0.18,2.68);
    \draw[thick] (-3.20,-1.3) -- (-2.29,1.05) -- (-0.97,1.05)
        (1.44,1.05) -- (3.39,1.05);
    \draw[thick,densely dashed] (-0.97,1.05) -- (B)
        (-1.95,-0.18) -- (D);
    \draw[very thick] (-3.20,-1.3) -- (2.48,-1.3) -- (3.39,1.05);
    \draw[thick] (A) -- (-0.97,1.56) -- (1.44,3.91) -- (B) -- (A)
        (D) -- (C);
    \draw[very thick] (D) -- (1.98,-0.18);
    \node[above left] at (A) {\(A\)};
    \node[above right] at (B) {\(B\)};
    \node[left] at (C) {\(C\)};
    \node[below right,xshift=2pt] at (D) {\(D\)};
    \node[right] at (1.98,-0.18) {\(E\)};
    \node[above] at (-1.95,-0.18) {\(E_1\)};
    \node[right] at (1.44,3.64) {\(Q\)};
    \node[right] at (2.9,-0.22) {\(P\)};
    \node[left] at (-2.81,-0.38) {\(P_1\)};
\end{tikzpicture}

    \caption{Прямой двугранный угол и его линейный угол}\label{fig:kiselev-stereo-335}
\end{figure}

\kiselevparagraph{para:ratio-of-dihedral-angles} \textbf{Теорема. Двугранные углы относятся, как их линейные углы.}

При доказательстве рассмотрим отдельно два случая:

\begin{figure}[htbp]
    \centering
    \renewcommand{\thefigure}{336}
    \captionsetup{name=Рис.,labelsep=period}
    \begin{tikzpicture}[line cap=round,line join=round]
    \begin{scope}
        \coordinate (A) at (0,2.08);
        \coordinate (B) at (0,-1.49);
        \draw[thick] (A) -- (B);
        % Три части угла AB.
        \foreach \x/\top/\bottom in {3.28/2.43/-1.14,3.03/1.73/-1.84,2.69/1.02/-2.55}{
            \draw[thick] (A) -- (\x,\top) -- (\x,\bottom);
            \draw[thick,densely dashed] (B) -- (\x,\bottom);
        }
        \draw[very thick] (A) -- (2.19,0.32) -- (2.19,-3.25) -- (B);
        \node[above left] at (A) {\(A\)};
        \node[left] at (B) {\(B\)};
        \node[right] at (3.28,-1.14) {\(C\)};
        \node[below right] at (2.19,-3.25) {\(D\)};
        \node[above right] at (3.28,2.43) {\(P\)};
        \node[below left] at (2.19,0.32) {\(Q\)};
    \end{scope}
    \begin{scope}[xshift=4.96cm]
        \coordinate (A1) at (0,2.08);
        \coordinate (B1) at (0,-1.49);
        \draw[thick] (A1) -- (B1);
        % Четыре части угла A1B1.
        \foreach \x/\top/\bottom in {3.9/3.21/-0.36,3.54/2.47/-1.10,3.20/1.71/-1.86,2.85/0.96/-2.61}{
            \draw[thick] (A1) -- (\x,\top) -- (\x,\bottom);
            \draw[thick,densely dashed] (B1) -- (\x,\bottom);
        }
        \draw[very thick] (A1) -- (2.41,0.21) -- (2.41,-3.36) -- (B1);
        \node[above left] at (A1) {\(A_1\)};
        \node[left] at (B1) {\(B_1\)};
        \node[right] at (3.9,-0.36) {\(C_1\)};
        \node[below right] at (2.41,-3.36) {\(D_1\)};
        \node[above right] at (3.9,3.21) {\(P_1\)};
        \node[below left] at (2.41,0.21) {\(Q_1\)};
    \end{scope}
\end{tikzpicture}

    \caption{Деление двугранных углов на равные части}\label{fig:kiselev-stereo-336}
\end{figure}

\begin{enumerate}[label=\arabic*),ref=\arabic*]
    \item \emph{Линейные углы \(CBD\) и \(C_1B_1D_1\) соизмеримы} (рис.~\ref{fig:kiselev-stereo-336}). Пусть их общая мера содержится в первом угле \(m\) раз, а во втором \(n\) раз. Проведём ряд плоскостей через рёбра и прямые, делящие линейные углы на части, равные общей мере. Тогда мы разделим угол \(AB\) на \(m\), а угол \(A_1B_1\) на
    % Печатная с. 296.
    \(n\) частей, которые все равны между собой вследствие равенства линейных углов. Поэтому
    \[
        \frac{\angle CBD}{\angle C_1B_1D_1}=\frac{m}{n}
        \quad\text{и}\quad
        \frac{\text{дв. уг. } AB}{\text{дв. уг. } A_1B_1}=\frac{m}{n}.
    \]
    Откуда
    \[
        \frac{\text{дв. уг. } AB}{\text{дв. уг. } A_1B_1}=
        \frac{\angle CBD}{\angle C_1B_1D_1}.
    \]

    \item \emph{Линейные углы несоизмеримы.} Разделим (рис.~\ref{fig:kiselev-stereo-336}) угол \(C_1B_1D_1\) на \(n\) равных частей. Пусть \(1/n\) этого угла содержится в угле \(CBD\) больше \(m\), но меньше \(m+1\) раз. Тогда приближённое отношение углов \(CBD\) и \(C_1B_1D_1\) с точностью до \(1/n\) (с недостатком) равно \(m/n\). Проведя плоскости так же, как и в первом случае, найдём, что приближённое отношение двугранных углов \(AB\) и \(A_1B_1\) с точностью до \(1/n\) (с недостатком) также равно \(m/n\). Таким образом, приближённые отношения равны при любом \(n\), а в этом случае мы условились считать несоизмеримые отношения \emph{равными}.
\end{enumerate}

\kiselevparagraph{para:measure-of-dihedral-angle} \textbf{Следствие.} Если за единицу измерения двугранных углов возьмём такой угол, который соответствует единице измерения линейных углов, то можно сказать, что \emph{двугранный угол измеряется его линейным углом}.


\section*{Перпендикулярные плоскости}
\addcontentsline{toc}{section}{Перпендикулярные плоскости}

\kiselevparagraph{para:perpendicular-planes-definition} \textbf{Определение.} Две плоскости называются \emph{взаимно перпендикулярными}, если, пересекаясь, они образуют \emph{прямые двугранные углы}. Возможность существования таких плоскостей показывает следующая теорема, указывающая \emph{признак перпендикулярности двух плоскостей}.

\kiselevparagraph{para:perpendicular-planes-criterion} \textbf{Теорема. Если плоскость} \(P\) (рис.~\ref{fig:kiselev-stereo-337}) \textbf{проходит через перпендикуляр} \(AB\) \textbf{к другой плоскости} \(Q\)\textbf{, то она перпендикулярна к этой плоскости.}

\begin{figure}[htbp]
    \centering
    \renewcommand{\thefigure}{337}
    \captionsetup{name=Рис.,labelsep=period}
    \begin{tikzpicture}[line cap=round,line join=round]
    \coordinate (D) at (-1.42,-1.38);
    \coordinate (E) at (1.41,1.08);
    \path (D) -- (E) coordinate[pos=0.49] (B);
    \coordinate (A) at (-0.033,2.38);
    \coordinate (C) at (1.97,-0.175);
    \draw[thick] (-3.27,-1.38) -- (-2.29,1.08) -- (-1.42,1.08)
        (E) -- (3.21,1.08);
    \draw[thick,densely dashed] (-1.42,1.08) -- (E);
    \draw[very thick] (-3.27,-1.38) -- (2.23,-1.38) -- (3.21,1.08);
    \draw[thick] (D) -- (-1.42,1.17) -- (1.41,3.63) -- (E) -- (D);
    \draw[very thick] (A) -- (B) -- (C);
    \node[above left] at (A) {\(A\)};
    \node[left] at (B) {\(B\)};
    \node[right] at (C) {\(C\)};
    \node[below] at (D) {\(D\)};
    \node[above right] at (E) {\(E\)};
    \node[above left] at (-1.42,1.17) {\(P\)};
    \node[right] at (2.89,0.29) {\(Q\)};
\end{tikzpicture}

    \caption{Плоскость \(P\) проходит через перпендикуляр \(AB\) к плоскости \(Q\)}\label{fig:kiselev-stereo-337}
\end{figure}

% Печатная с. 297.
В плоскости \(Q\) проведём \(BC\perp DE\). Тогда угол \(ABC\) будет линейным углом двугранного угла \(PDEQ\). Так как прямая \(AB\), по условию, перпендикулярна к \(Q\), то \(AB\perp BC\). Значит, угол \(ABC\) прямой, а потому и двугранный угол прямой (\hyperref[item:right-dihedral-and-linear-angle]{см.~§~\ref*{para:right-and-equal-dihedral-angles}, пункт~\ref*{item:right-dihedral-and-linear-angle}}), то есть плоскость \(P\) перпендикулярна к \(Q\).

Следующие две теоремы выражают \emph{главные свойства перпендикулярных плоскостей}.

\kiselevparagraph{para:normal-in-perpendicular-plane} \textbf{Теорема. Если две плоскости} \(P\) и \(Q\) (рис.~\ref{fig:kiselev-stereo-338}) \textbf{взаимно перпендикулярны и к одной из них} (к \(Q\)) \textbf{проведён перпендикуляр} \(AB\)\textbf{, имеющий общую точку} \(A\) \textbf{с другой плоскостью} (с \(P\))\textbf{, то он весь лежит в этой плоскости.}

\begin{figure}[htbp]
    \centering
    \renewcommand{\thefigure}{338}
    \captionsetup{name=Рис.,labelsep=period}
    \begin{tikzpicture}[line cap=round,line join=round]
    \coordinate (D) at (-0.81,-1.50);
    \coordinate (E) at (0.96,1.05);
    \path (D) -- (E) coordinate[pos=0.408] (C);
    \coordinate (A) at (-0.088,2.12);
    \coordinate (B) at (0.52,-1.03);
    \coordinate (F) at (1.95,-0.46);
    \draw[thick] (-3.12,-1.50) -- (-2.18,1.05) -- (-0.81,1.05)
        (E) -- (3.54,1.05);
    \draw[thick,densely dashed] (-0.81,1.05) -- (E);
    \draw[very thick] (-3.12,-1.50) -- (2.60,-1.50) -- (3.54,1.05);
    \draw[thick] (D) -- (-0.81,1.08) -- (0.96,3.63) -- (E) -- (D);
    \draw[very thick] (A) -- (B);
    \draw[thick,densely dashed] (A) -- (C) -- (F);
    \node[above left] at (A) {\(A\)};
    \node[right,xshift=2pt] at (B) {\(B\)};
    \node[left] at (C) {\(C\)};
    \node[below] at (D) {\(D\)};
    \node[above right] at (E) {\(E\)};
    \node[right] at (F) {\(F\)};
    \node[above right] at (0.96,3.63) {\(P\)};
    \node at (2.67,0.55) {\(Q\)};
\end{tikzpicture}

    \caption{Предположение, что \(AB\) не лежит в плоскости \(P\)}\label{fig:kiselev-stereo-338}
\end{figure}

Предположим, что \(AB\) не лежит в плоскости \(P\), как изображено на рисунке. Проведём в этой плоскости \(AC\perp DE\), а в плоскости \(Q\) проведём \(CF\perp DE\). Тогда угол \(ACF\), как линейный угол прямого двугранного угла, будет прямым. Поэтому прямая \(AC\), образуя прямые углы с \(DE\) и \(CF\), будет перпендикуляром к плоскости \(Q\). У нас тогда будет два перпендикуляра, опущенных из одной и той же точки \(A\) на плоскость \(Q\), именно \(AB\) и \(AC\). Так как это невозможно, то нельзя допустить, чтобы перпендикуляр \(AB\) не лежал в плоскости \(P\).

\kiselevparagraph{para:intersection-normal-to-third-plane} \textbf{Следствие. Пересечение} \(AB\) (рис.~\ref{fig:kiselev-stereo-339}) \textbf{двух плоскостей} \(P\) и \(Q\)\textbf{, перпендикулярных к третьей плоскости} \(R\)\textbf{, есть перпендикуляр к этой плоскости.}

\begin{figure}[htbp]
    \centering
    \renewcommand{\thefigure}{339}
    \captionsetup{name=Рис.,labelsep=period}
    \begin{tikzpicture}[line cap=round,line join=round]
    \coordinate (A) at (-0.60,3.29);
    \coordinate (B) at (-0.60,-0.50);
    \draw[thick] (-3.0,-1.58) -- (-2.08,0.99) -- (-0.60,0.99)
        (1.75,0.99) -- (3.14,0.99);
    \draw[thick,densely dashed] (-0.60,0.99) -- (1.75,0.99);
    \draw[very thick] (-3.0,-1.58) -- (2.22,-1.58) -- (3.14,0.99);
    \draw[thick] (B) -- (A) -- (1.75,3.72) -- (1.75,-0.07);
    \draw[thick,densely dashed] (B) -- (1.75,-0.07);
    \draw[very thick] (A) -- (0.96,2.86) -- (0.96,-0.93) -- (B);
    \node[above left] at (A) {\(A\)};
    \node[left] at (B) {\(B\)};
    \node[below left] at (1.75,3.72) {\(P\)};
    \node[below left] at (0.96,2.86) {\(Q\)};
    \node at (2.51,0.51) {\(R\)};
\end{tikzpicture}

    \caption{Линия пересечения \(AB\) перпендикулярна к плоскости \(R\)}\label{fig:kiselev-stereo-339}
\end{figure}

% Доказательство продолжается на печатной с. 298.
Действительно, если через какую-нибудь точку \(A\) линии пересечения проведём перпендикуляр к плоскости \(R\), то этот перпендикуляр, согласно \hyperref[para:normal-in-perpendicular-plane]{предыдущей теореме}, должен лежать и в плоскости \(Q\), и в плоскости \(P\). Значит, он сольётся с \(AB\).

\section*{Угол двух скрещивающихся прямых}
\addcontentsline{toc}{section}{Угол двух скрещивающихся прямых}

\kiselevparagraph{para:angle-between-skew-lines} Мы видели (\hyperref[para:skew-and-parallel-lines]{см.~§~\ref*{para:skew-and-parallel-lines}}), что в пространстве можно взять такие прямые, которые не пересекаются и в то же время не параллельны, так как не лежат в одной плоскости. Для таких \emph{скрещивающихся} прямых дадим следующее определение.

\textbf{Определение.} \emph{Углом двух скрещивающихся прямых} \(AB\) и \(CD\) (рис.~\ref{fig:kiselev-stereo-340}), положение и направление которых заданы, называется угол \(MON\), который получится, если из произвольной точки пространства \(O\) проведём полупрямые \(OM\) и \(ON\), \emph{соответственно параллельные данным прямым} \(AB\) и \(CD\) \emph{и одинаково с ними направленные}.

\begin{figure}[htbp]
    \centering
    \renewcommand{\thefigure}{340}
    \captionsetup{name=Рис.,labelsep=period}
    \begin{tikzpicture}[line cap=round,line join=round]
    \coordinate (A) at (-3.55,1.29);
    \coordinate (B) at (-0.61,1.92);
    \coordinate (C) at (-3.55,0.40);
    \coordinate (D) at (-0.61,0.21);
    \coordinate (O) at (0.48,0.82);
    \coordinate (M) at (3.02,1.3643);
    \coordinate (N) at (3.02,0.6559);
    \coordinate (O1) at (-1.93,-1.28);
    \coordinate (M1) at (0.63,-0.7314);
    \coordinate (N1) at (0.63,-1.4454);

    \foreach \first/\last in {A/B,C/D,O/M,O/N,O1/M1,O1/N1}{
        \path (\first) -- (\last) coordinate[pos=0.79] (arrow);
        \draw[very thick,->] (\first) -- (arrow);
        \draw[very thick] (arrow) -- (\last);
    }
    \node[left] at (A) {\(A\)};
    \node[right] at (B) {\(B\)};
    \node[left] at (C) {\(C\)};
    \node[right] at (D) {\(D\)};
    \node[left] at (O) {\(O\)};
    \node[above right] at (M) {\(M\)};
    \node[below right] at (N) {\(N\)};
    \node[left] at (O1) {\(O_1\)};
    \node[above right] at (M1) {\(M_1\)};
    \node[below right] at (N1) {\(N_1\)};
\end{tikzpicture}

    \caption{Угол \(MON\) двух скрещивающихся прямых \(AB\) и \(CD\)}\label{fig:kiselev-stereo-340}
\end{figure}

Величина угла \(MON\) не зависит от положения точки \(O\). В самом деле, если построим указанным способом угол \(M_1O_1N_1\) при какой-нибудь другой точке \(O_1\), то \(MON=M_1O_1N_1\), так как эти углы имеют соответственно параллельные и одинаково направленные стороны (\hyperref[para:lines-parallel-to-same-line]{см.~§~\ref*{para:lines-parallel-to-same-line}} и \hyperref[para:angles-with-parallel-sides]{§~\ref*{para:angles-with-parallel-sides}}).

\section*{Угол, образуемый прямой с плоскостью}
\addcontentsline{toc}{section}{Угол, образуемый прямой с плоскостью}

\kiselevparagraph{para:orthogonal-projection-on-plane} \textbf{Проекция прямой на плоскость.} Мы говорили прежде (\hyperref[para:normal-and-oblique-lengths]{см.~§~\ref*{para:normal-and-oblique-lengths}}), что, когда из одной точки проведены к плоскости перпендикуляр и наклонная, то проекцией этой наклонной на плоскость называется прямая, соединяющая основание
% Печатная с. 299.
перпендикуляра с основанием наклонной. Дадим теперь более общее определение проекции.

\begin{enumerate}[label=\arabic*),ref=\arabic*]
    \item \emph{Проекцией какой-нибудь точки на плоскость} (например, точки \(M\) на плоскость \(P\), рис.~\ref{fig:kiselev-stereo-341}) называется \emph{основание} \(m\) \emph{перпендикуляра, опущенного на эту плоскость из взятой точки}.
    \item \emph{Проекцией какой-нибудь линии на плоскость} называется \emph{геометрическое место проекций всех точек этой линии}.
\end{enumerate}

\begin{figure}[htbp]
    \centering
    \renewcommand{\thefigure}{341}
    \captionsetup{name=Рис.,labelsep=period}
    \begin{tikzpicture}[line cap=round,line join=round]
    \coordinate (A) at (-2.00,1.35);
    \coordinate (B) at (1.88,2.96);
    \coordinate (a) at (-2.00,-0.54);
    \coordinate (b) at (1.88,-0.19);
    \path (A) -- (B) coordinate[pos=0.30] (M) coordinate[pos=0.66] (N);
    \path (a) -- (b) coordinate[pos=0.30] (m) coordinate[pos=0.66] (n);
    \draw[thick] (-3.31,-1.45) -- (-2.30,0.87) -- (3.31,0.87);
    \draw[very thick] (-3.31,-1.45) -- (2.30,-1.45) -- (3.31,0.87);
    \draw[thick] (a) -- (A) (m) -- (M) (n) -- (N) (b) -- (B) (a) -- (b);
    \draw[very thick] (A) -- (B);
    \foreach \point in {M,N} {
        \fill (\point) circle[radius=1pt];
    }
    \node[left] at (A) {\(A\)};
    \node[above] at (B) {\(B\)};
    \node[above] at (M) {\(M\)};
    \node[above] at (N) {\(N\)};
    \foreach \point in {a,m,n,b} {
        \node[below] at (\point) {\(\point\)};
    }
    \node at (-2.60,-1.06) {\(P\)};
    \node at (1.19,1.48) {\(Q\)};
\end{tikzpicture}

    \caption{Проекции точек \(A\), \(M\), \(N\), \(B\) прямой \(AB\) на плоскость \(P\)}\label{fig:kiselev-stereo-341}
\end{figure}

В частности, если проецируемая линия --- \emph{прямая} (например, \(AB\), рис.~\ref{fig:kiselev-stereo-341}), то её проекция на плоскость \(P\) --- также \emph{прямая}. В самом деле, если через прямую \(AB\) и перпендикуляр \(Mm\), опущенный на плоскость проекций из какой-нибудь одной точки \(M\) этой прямой, проведём плоскость \(Q\), то эта плоскость должна быть перпендикулярна к плоскости \(P\) (\hyperref[para:perpendicular-planes-criterion]{см.~§~\ref*{para:perpendicular-planes-criterion}}). Поэтому перпендикуляр, опущенный на плоскость \(P\) из любой точки прямой \(AB\) (например, из точки \(N\)), должен лежать в этой плоскости \(Q\) (\hyperref[para:normal-in-perpendicular-plane]{см.~§~\ref*{para:normal-in-perpendicular-plane}}). Следовательно, проекции всех точек прямой \(AB\) должны лежать на прямой \(ab\), по которой пересекаются плоскости \(P\) и \(Q\). Таким образом, эта прямая \(ab\) представляет собой геометрическое место проекций всех точек данной прямой \(AB\) и, следовательно, является её проекцией.

Существует, впрочем, один частный случай, когда \emph{проекция прямой становится точкой}. Это бывает тогда, когда прямая перпендикулярна к плоскости проекций.

\kiselevparagraph{para:angle-between-line-and-plane} \textbf{Угол прямой с плоскостью.} \emph{Углом прямой} \(AB\) (рис.~\ref{fig:kiselev-stereo-342}) \emph{с плоскостью} \(P\) в том случае, когда прямая наклонна к плоскости, называется \emph{угол} \(ABC\)\emph{, образованный этой прямой и её проекцией на плоскость}.

\begin{figure}[htbp]
    \centering
    \renewcommand{\thefigure}{342}
    \captionsetup{name=Рис.,labelsep=period}
    \begin{tikzpicture}[line cap=round,line join=round]
    \coordinate (A) at (1.1,2.9);
    \coordinate (B) at (-1.61,-0.18);
    \coordinate (C) at (1.1,0.63);
    \coordinate (D) at (0.56,-1.03);
    \draw[thick] (-3.05,-1.62) -- (-1.98,1.28) -- (3.14,1.28);
    \draw[very thick] (-3.05,-1.62) -- (2.07,-1.62) -- (3.14,1.28);
    \draw[thick,densely dashed] (A) -- (C) (A) -- (D);
    \draw[very thick] (A) -- (B) -- (C);
    \draw[thick] (B) -- (D);
    \draw[thick] (B) ++(16.64:0.85) arc[start angle=16.64,end angle=48.65,radius=0.85];
    \path (B) -- (C) coordinate[pos=0.57] (CtickA) coordinate[pos=0.61] (CtickB);
    \path (B) -- (D) coordinate[pos=0.57] (DtickA) coordinate[pos=0.61] (DtickB);
    \draw[thick] (CtickA) ++(0.046,-0.146) -- ++(-0.092,0.292)
        (CtickB) ++(0.046,-0.146) -- ++(-0.092,0.292)
        (DtickA) ++(-0.052,-0.142) -- ++(0.104,0.284)
        (DtickB) ++(-0.052,-0.142) -- ++(0.104,0.284);
    \node[above] at (A) {\(A\)};
    \node[left] at (B) {\(B\)};
    \node[right] at (C) {\(C\)};
    \node[right] at (D) {\(D\)};
    \node at (-2.36,-0.87) {\(P\)};
\end{tikzpicture}

    \caption{Угол \(ABC\), образованный наклонной \(AB\) и её проекцией \(BC\)}\label{fig:kiselev-stereo-342}
\end{figure}

% Печатная с. 300.
Этот угол обладает тем свойством, что он \emph{наименьший} из всех углов, которые наклонная образует с прямыми, проведёнными в плоскости \(P\) через основание наклонной. Докажем, например, что \(\angle ABC\) меньше \(\angle ABD\). Для этого отложим \(BD=BC\) и соединим \(D\) с \(A\). У треугольников \(ABC\) и \(ABD\) две стороны одного равны соответственно двум сторонам другого, но третьи стороны не равны, а именно \(AD>AC\) (\hyperref[para:normal-and-oblique-lengths]{см.~§~\ref*{para:normal-and-oblique-lengths}}). Вследствие этого \(\angle ABD\) больше \(\angle ABC\) (\kiselevPlanimetryPage{51}{см.~Планиметрия, §~58, пункт~2}).

\clearpage
\endgroup
  % Глава IV
    \chapter{Многогранные углы}\label{ch:ch5}

\begingroup
\setlength{\intextsep}{10pt plus 2pt minus 2pt}
\setlength{\textfloatsep}{12pt plus 2pt minus 2pt}

% А. П. Киселёв, «Элементарная геометрия», 23-е издание, 1914. Печатные с. 300–306.
% Страницы скана планиметрии на 11 больше печатных.
\newcommand*{\kiselevPlanimetryPage}[2]{%
    \href{https://commons.wikimedia.org/wiki/File:\%D0\%90._\%D0\%9F._\%D0\%9A\%D0\%B8\%D1\%81\%D0\%B5\%D0\%BB\%D1\%91\%D0\%B2._\%D0\%AD\%D0\%BB\%D0\%B5\%D0\%BC\%D0\%B5\%D0\%BD\%D1\%82\%D0\%B0\%D1\%80\%D0\%BD\%D0\%B0\%D1\%8F_\%D0\%B3\%D0\%B5\%D0\%BE\%D0\%BC\%D0\%B5\%D1\%82\%D1\%80\%D0\%B8\%D1\%8F_(1914).djvu?page=#1}{#2}%
}

% Печатная с. 300.
\kiselevparagraph{para:polyhedral-angle-definition} \textbf{Определение.} Возьмём несколько углов (рис.~\ref{fig:kiselev-stereo-343}): \(ASB\), \(BSC\), \(CSD\), \(\ldots\), которые, примыкая последовательно друг к другу, расположены в одной плоскости вокруг общей вершины \(S\).

\begin{figure}[htbp]
    \centering
    \renewcommand{\thefigure}{343}
    \captionsetup{name=Рис.,labelsep=period}
    \begin{tikzpicture}[line cap=round,line join=round]
    \coordinate (S) at (0,0);
    \draw[thick] (S) -- (-2.94,2.41) node[above left] {\(A\)}
        (S) -- (-3.58,-0.25) node[left] {\(B\)}
        (S) -- (-1.96,-2.84) node[below left] {\(C\)}
        (S) -- (0.27,-3.29) node[below] {\(D\)}
        (S) -- (2.24,-2.76) node[below right] {\(E\)}
        (S) -- (2.69,2.57) node[above right] {\(F\)};
    \node[above] at (S) {\(S\)};
\end{tikzpicture}

    \caption{Плоские углы вокруг общей вершины \(S\)}\label{fig:kiselev-stereo-343}
\end{figure}

Повернём плоскость угла \(ASB\) вокруг общей стороны \(SB\) так, чтобы эта плоскость образовала некоторый двугранный угол с плоскостью \(BSC\). Затем, сохраняя получившийся двугранный угол, повернём его вокруг прямой
% Печатная с. 301.
\(SC\) так, чтобы плоскость \(BSC\) образовала некоторый двугранный угол с плоскостью \(CSD\). Продолжим такое последовательное вращение вокруг каждой общей стороны. Если при этом последняя сторона \(SF\) совместится с первой стороной \(SA\), то образуется фигура (рис.~\ref{fig:kiselev-stereo-344}), которая вместе с частью пространства, ограниченной плоскостями, называется \emph{многогранным углом}.

\begin{figure}[htbp]
    \centering
    \renewcommand{\thefigure}{344}
    \captionsetup{name=Рис.,labelsep=period}
    \begin{tikzpicture}[line cap=round,line join=round]
    \coordinate (S) at (0,3.33);
    \coordinate (a) at (-2.22,-0.22);
    \coordinate (b) at (-0.98,-1.39);
    \coordinate (c) at (0.80,-1.52);
    \coordinate (d) at (2.35,-0.46);
    \coordinate (e) at (0.05,0.50);

    \foreach \vertex/\tip in {a/A,b/B,c/C,d/D}{
        \path (S) -- (\vertex) coordinate[pos=1.20] (\tip);
        \draw[thick] (S) -- (\tip);
    }
    \path (S) -- (e) coordinate[pos=1.33] (E);
    \draw[thick,densely dashed] (S) -- (E) (a) -- (e) -- (d);
    \draw[very thick] (a) -- (b) -- (c) -- (d);
    \node[above] at (S) {\(S\)};
    \node[left] at (a) {\(a\)};
    \node[below left,xshift=-1pt] at (b) {\(b\)};
    \node[above right] at (c) {\(c\)};
    \node[right] at (d) {\(d\)};
    \node[above left,yshift=1pt] at (e) {\(e\)};
    \node[below left] at (A) {\(A\)};
    \node[below left] at (B) {\(B\)};
    \node[below right] at (C) {\(C\)};
    \node[below right] at (D) {\(D\)};
    \node[below] at (E) {\(E\)};
\end{tikzpicture}

    \caption{Выпуклый многогранный угол \(SABCDE\)}\label{fig:kiselev-stereo-344}
\end{figure}

Углы \(ASB,BSC,\ldots\) называются \emph{плоскими углами}, или \emph{гранями}, их стороны \(SA,SB,\ldots\) называются \emph{рёбрами}, а общая вершина \(S\) --- \emph{вершиной} многогранного угла. Каждому ребру соответствует свой двугранный угол. Поэтому в многогранном угле столько двугранных углов и столько плоских, сколько в нём всего рёбер. Наименьшее число граней в многогранном угле --- три. Такой угол называется \emph{трёхгранным}. Могут быть углы четырёхгранные, пятигранные и т.~д.

Многогранный угол (рис.~\ref{fig:kiselev-stereo-344}) обозначается или одной буквой \(S\), поставленной у вершины, или рядом букв \(SABCDE\), из которых первая обозначает вершину, а остальные --- рёбра по порядку их расположения.

Многогранный угол называется \emph{выпуклым}, если он весь расположен по одну сторону от каждой своей грани. Таков, например, угол, изображённый на рис.~\ref{fig:kiselev-stereo-344}. Наоборот, угол на рис.~\ref{fig:kiselev-stereo-345} нельзя назвать выпуклым, так как он расположен по обе стороны от грани \(ASB\) или от грани \(BSC\).

\begin{figure}[htbp]
    \centering
    \renewcommand{\thefigure}{345}
    \captionsetup{name=Рис.,labelsep=period}
    \begin{tikzpicture}[line cap=round,line join=round]
    \coordinate (S) at (0,3.40);
    \coordinate (a) at (-2.28,-0.30);
    \coordinate (b) at (-0.62,-0.30);
    \coordinate (c) at (0.14,-1.31);
    \coordinate (d) at (2.01,-0.75);
    \coordinate (e) at (0.45,0.81);

    \foreach \vertex/\tip in {a/A,b/B,c/C,d/D}{
        \path (S) -- (\vertex) coordinate[pos=1.23] (\tip);
        \draw[thick] (S) -- (\tip);
    }
    \path (S) -- (e) coordinate[pos=1.31] (E);
    \draw[thick,densely dashed] (S) -- (E) (a) -- (e) -- (d);
    \draw[very thick] (a) -- (b) -- (c) -- (d);
    \node[above] at (S) {\(S\)};
    \node[left] at (a) {\(a\)};
    \node[above right] at (b) {\(b\)};
    \node[above right] at (c) {\(c\)};
    \node[right] at (d) {\(d\)};
    \node[above right] at (e) {\(e\)};
    \node[below left] at (A) {\(A\)};
    \node[below] at (B) {\(B\)};
    \node[below] at (C) {\(C\)};
    \node[below right] at (D) {\(D\)};
    \node[right] at (E) {\(E\)};
\end{tikzpicture}

    \caption{Многогранный угол по обе стороны от грани \(ASB\)}\label{fig:kiselev-stereo-345}
\end{figure}

Если две грани\footnote{В издании 1914 года напечатано «две грани». На рис.~\ref{fig:kiselev-stereo-344} и \ref{fig:kiselev-stereo-345} секущая плоскость пересекает все пять граней.} многогранного угла пересечём плоскостью, то в сечении образуется многоугольник \(abcde\) (рис.~\ref{fig:kiselev-stereo-344} и \ref{fig:kiselev-stereo-345}). У выпуклого угла этот многоугольник тоже выпуклый.

Мы будем рассматривать только выпуклые многогранные углы.

\kiselevparagraph{para:trihedral-plane-angle-inequality} \textbf{Теорема. В каждом трёхгранном угле каждый плоский угол меньше суммы двух других плоских углов.}

Очевидно, эта теорема нуждается в доказательстве только в том случае, когда она применяется к плоскому углу,
% Печатная с. 302.
\emph{наибольшему из трёх}. Пусть в угле \(SABC\) (рис.~\ref{fig:kiselev-stereo-346}) наибольший из плоских углов --- \(ASC\). Докажем, что даже этот наибольший угол меньше суммы двух остальных.

\begin{figure}[htbp]
    \centering
    \renewcommand{\thefigure}{346}
    \captionsetup{name=Рис.,labelsep=period}
    \begin{tikzpicture}[line cap=round,line join=round]
    \coordinate (S) at (0,3.4);
    \coordinate (A) at (-2.00,-0.35);
    \coordinate (B) at (-0.23,-1.15);
    \coordinate (C) at (1.95,-0.24);
    \path (A) -- (C) coordinate[pos=0.618] (D);
    \foreach \point in {A,B,C}{
        \path (S) -- (\point) coordinate[pos=1.17] (tip);
        \draw[thick] (S) -- (tip);
    }
    \path (S) -- (D) coordinate[pos=1.32] (tip);
    \draw[very thick] (A) -- (B) -- (C);
    \draw[thick,densely dashed] (A) -- (C) (S) -- (tip);
    \draw[thick] (S) ++(-118.07:0.85) arc[start angle=-118.07,end angle=-92.90,radius=0.85];
    \draw[thick] (S) ++(-118.07:1.02) arc[start angle=-118.07,end angle=-83.18,radius=1.02];
    \path (S) -- (B) coordinate[pos=0.59] (BmarkA) coordinate[pos=0.63] (BmarkB);
    \path (S) -- (D) coordinate[pos=0.59] (DmarkA) coordinate[pos=0.63] (DmarkB);
    \draw[thick] (BmarkA) ++(-0.15,0.01) -- ++(0.30,-0.02)
        (BmarkB) ++(-0.15,0.01) -- ++(0.30,-0.02)
        (DmarkA) ++(-0.15,-0.018) -- ++(0.30,0.036)
        (DmarkB) ++(-0.15,-0.018) -- ++(0.30,0.036);
    \node[above] at (S) {\(S\)};
    \node[left] at (A) {\(A\)};
    \node[below left] at (B) {\(B\)};
    \node[right] at (C) {\(C\)};
    \node[above right] at (D) {\(D\)};
\end{tikzpicture}

    \caption{Плоские углы трёхгранного угла \(SABC\)}\label{fig:kiselev-stereo-346}
\end{figure}

Отложим на угле \(ASC\) часть \(ASD\), равную \(ASB\). Проведём к плоскости\footnote{В издании 1914 года стоит «к плоскости». На рис.~\ref{fig:kiselev-stereo-346} прямая \(AC\) лежит в плоскости угла \(ASC\).} угла \(ASC\) какую-нибудь прямую \(AC\), пересекающую \(SD\) в некоторой точке \(D\). Отложим \(SB=SD\). Соединив \(B\) с \(A\) и \(C\), получим треугольник \(ACB\), в котором
\[
    AD+DC<AB+BC.
\]

Треугольники \(ASD\) и \(ASB\) равны, так как в каждом из них равный угол заключён между равными сторонами. Следовательно, \(AD=AB\). Поэтому в полученном неравенстве можно отбросить равные части \(AD\) и \(AB\), после чего получим
\[
    DC<BC.
\]

Теперь заметим, что у треугольников \(SCD\) и \(SCB\) две стороны одного равны двум сторонам другого, а третьи стороны не равны. В таком случае против большей из этих сторон лежит больший угол (\kiselevPlanimetryPage{51}{см.~Планиметрия, §~58, пункт~2}). Значит,
\[
    \angle CSD<\angle CSB.
\]

Прибавив к левой части этого неравенства угол \(ASD\), а к правой равный ему угол \(ASB\), получим неравенство, которое требовалось доказать:
\[
    \angle ASC<\angle CSB+\angle ASB.
\]

\textbf{Следствие.} Вычтя из обеих частей последнего неравенства по углу \(ASB\) или по углу \(CSB\), получим\footnote{Обе строки воспроизводят обозначения печатного издания. На рис.~\ref{fig:kiselev-stereo-346} точки \(E\) нет. В предыдущем неравенстве участвуют углы \(ASC\), \(CSB\) и \(ASB\).}:
\[
    \begin{aligned}
        \angle ASC-\angle ASE &< \angle CSB;\\
        \angle ASB-\angle CSE &< \angle ASB.
    \end{aligned}
\]

Так как, кроме того, угол \(ASC\), наибольший из трёх углов, конечно, больше разности двух других углов, то заключаем:

\textbf{В трёхгранном угле каждый плоский угол больше разности двух других углов.}

% Печатная с. 303.
\kiselevparagraph{para:convex-polyhedral-angle-sum} \textbf{Теорема. В каждом выпуклом многогранном угле сумма всех плоских углов меньше} \(4d\)\footnote{Буквой \(d\) обозначен прямой угол.}.

\begin{figure}[htbp]
    \centering
    \renewcommand{\thefigure}{347}
    \captionsetup{name=Рис.,labelsep=period}
    \begin{tikzpicture}[line cap=round,line join=round]
    \coordinate (S) at (0,2.91);
    \coordinate (A) at (-2.70,-0.50);
    \coordinate (B) at (-1.37,-1.17);
    \coordinate (C) at (1.35,-1.19);
    \coordinate (D) at (2.82,-0.39);
    \coordinate (E) at (0,0.18);

    \foreach \point in {A,B,C,D}{
        \path (S) -- (\point) coordinate[pos=1.16] (tip);
        \draw[thick] (S) -- (tip);
    }
    \draw[thick,densely dashed] (S) -- (0,-0.65) (A) -- (E) -- (D);
    \draw[very thick] (A) -- (B) -- (C) -- (D);
    \node[above] at (S) {\(S\)};
    \node[left] at (A) {\(A\)};
    \node[below left] at (B) {\(B\)};
    \node[below right] at (C) {\(C\)};
    \node[right] at (D) {\(D\)};
    \node[above right] at (E) {\(E\)};
\end{tikzpicture}

    \caption{Сечение выпуклого многогранного угла \(SABCDE\)}\label{fig:kiselev-stereo-347}
\end{figure}

Пересечём грани (рис.~\ref{fig:kiselev-stereo-347}) выпуклого угла \(SABCDE\) какой-нибудь плоскостью. В сечении получим выпуклый \(n\)-угольник \(ABCDE\). Применяя \hyperref[para:trihedral-plane-angle-inequality]{теорему предыдущего параграфа} к каждому из трёхгранных углов, образовавшихся при точках \(A\), \(B\), \(C\), \(D\) и \(E\), находим:
\[
    \angle ABC<\angle ABS+\angle SBC;\quad
    \angle BCD<\angle BCS+\angle SCD;\quad \ldots
\]
Сложим почленно все эти неравенства. Тогда в левой части получим сумму всех углов многоугольника \(ABCDE\), которая равна \(2dn-4d\) (\kiselevPlanimetryPage{78}{см.~Планиметрия,~§~89}), а в правой --- сумму углов треугольников \(ASB,BSC,\ldots\), кроме тех углов, которые лежат при вершине \(S\). Обозначив сумму этих последних углов буквой \(x\), получим после сложения:
\[
    2dn-4d<2dn-x.
\]
Откуда
\[
    x<4d.
\]


\section*{Равенство трёхгранных углов}
\addcontentsline{toc}{section}{Равенство трёхгранных углов}

\kiselevparagraph{para:supplementary-trihedral-angle} \textbf{Дополнительный угол.} Из вершины \(S\) (рис.~\ref{fig:kiselev-stereo-348}) трёхгранного угла \(SABC\) проведём к грани \(ASB\) перпендикуляр \(SC_1\), направив его в ту сторону от этой грани, в которой расположено противоположное ребро \(SC\). Подобным образом проведём перпендикуляр \(SA_1\) к грани \(BSC\) и \(SB_1\) к грани \(ASC\). Трёхгранный угол, у которого рёбрами служат полупрямые \(SA_1,SB_1\) и \(SC_1\), называется \emph{дополнительным} для угла \(SABC\).

\begin{figure}[htbp]
    \centering
    \renewcommand{\thefigure}{348}
    \captionsetup{name=Рис.,labelsep=period}
    \begin{tikzpicture}[line cap=round,line join=round]
    \coordinate (S) at (0,2.31);
    \draw[very thick] (S) -- (-1.93,-1.02) node[below left] {\(A\)}
        (S) -- (-0.06,-1.70) node[below] {\(B\)}
        (S) -- (1.08,-1.51) node[below right] {\(C\)};
    \draw[thick,densely dashed] (S) -- (-3.26,-0.12) node[below left] {\(A_1\)}
        (S) -- (-0.61,-0.24) node[below] {\(B_1\)}
        (S) -- (3.25,-0.72) node[below right] {\(C_1\)};
    \node[above] at (S) {\(S\)};
\end{tikzpicture}

    \caption{Трёхгранный угол \(SABC\) и дополнительный угол \(SA_1B_1C_1\)}\label{fig:kiselev-stereo-348}
\end{figure}

Заметим, что если для угла \(SABC\) дополнительным углом служит угол \(SA_1B_1C_1\), то и наоборот: для угла \(SA_1B_1C_1\) дополнительным будет \(SABC\). Действительно, плоскость \(SA_1B_1\), проходящая через перпендикуляры к плоскостям \(BSC\) и \(ASC\), перпендикулярна
% Печатная с. 304.
к ним обеим, а следовательно, и к линии их пересечения \(SC\). Значит, прямая \(SC\) перпендикулярна к грани \(SA_1B_1\) и, кроме того, расположена по ту же сторону от этой грани, по которую лежит противоположное ребро \(SC_1\). Подобным образом убедимся, что прямые \(SB\) и \(SA\) соответственно перпендикулярны к граням \(SA_1C_1\) и \(SB_1C_1\) и расположены по ту сторону от них, по которую лежат рёбра \(SB_1\) и \(SA_1\). Значит, углы \(SABC\) и \(SA_1B_1C_1\) \emph{взаимно дополнительны}.

\kiselevparagraph{para:supplementary-trihedral-angle-lemma} \textbf{Лемма 1. Если два трёхгранных угла взаимно дополнительны, то плоские углы одного служат дополнением до} \(2d\) \textbf{к противоположным двугранным углам другого.}

Каждый плоский угол одного из взаимно дополнительных трёхгранных углов образован двумя \emph{перпендикулярами}, проведёнными к граням противоположного двугранного угла другого трёхгранного угла из одной точки его ребра. Заметив это и приняв во внимание направление перпендикуляров, возьмём какой-нибудь двугранный угол \(AB\) (рис.~\ref{fig:kiselev-stereo-349}) и из произвольной точки \(B\) его ребра проведём перпендикуляры: \(BE\) к грани \(AD\) и \(BF\) к грани \(AC\).

\begin{figure}[htbp]
    \centering
    \renewcommand{\thefigure}{349}
    \captionsetup{name=Рис.,labelsep=period}
    \begin{tikzpicture}[line cap=round,line join=round]
    \coordinate (A) at (0,2.08);
    \coordinate (B) at (0,-1.62);
    \coordinate (C) at (2.14,-2.86);
    \coordinate (D) at (2.95,-1.27);
    \coordinate (E) at (0.54,-3.66);
    \coordinate (F) at (3.56,0.85);

    \draw[thick] (B) -- (A) -- (2.95,2.43) -- (D);
    \draw[very thick] (A) -- (2.14,0.84) -- (C) -- (B);
    \draw[thick,densely dashed] (B) -- (D) (B) -- (E) (B) -- (F);
    \node[above left] at (A) {\(A\)};
    \node[left] at (B) {\(B\)};
    \node[below right] at (C) {\(C\)};
    \node[right] at (D) {\(D\)};
    \node[below] at (E) {\(E\)};
    \node[right] at (F) {\(F\)};
\end{tikzpicture}

    \caption{Перпендикуляры \(BE\) и \(BF\) к граням двугранного угла \(AB\)}\label{fig:kiselev-stereo-349}
\end{figure}

Затем проведём через \(BE\) и \(BF\) плоскость, которая должна быть перпендикулярна к ребру \(AB\) (\hyperref[para:perpendicular-planes-criterion]{см.~§~\ref*{para:perpendicular-planes-criterion}} и \hyperref[para:intersection-normal-to-third-plane]{§~\ref*{para:intersection-normal-to-third-plane}}). Пусть пересечения этой плоскости с гранями угла \(AB\) будут прямые \(BC\) и \(BD\). Тогда угол \(CBD\) должен быть линейным углом двугранного угла \(AB\). Так как стороны угла \(EBF\) соответственно перпендикулярны к сторонам угла \(CBD\) и эти углы не равны, то их сумма равна \(2d\) (\kiselevPlanimetryPage{76}{см.~Планиметрия, §~86}). Это и требовалось доказать.

\kiselevparagraph{para:equal-supplementary-trihedral-angles} \textbf{Лемма 2. Равным трёхгранным углам соответствуют равные дополнительные углы и обратно.}

Равные трёхгранные углы при наложении совмещаются. Поэтому совмещаются и те перпендикуляры, которые образуют рёбра дополнительных углов. Значит, дополнительные углы также совмещаются.

Обратно, если совмещаются дополнительные углы, то совмещаются и данные углы.

\kiselevparagraph{para:trihedral-angle-congruence} \textbf{Теоремы. Трёхгранные углы равны, если они имеют:}

\begin{enumerate}[label=\arabic*),ref=\arabic*]
    \item \textbf{по равному двугранному углу, заключённому между двумя соответственно равными и одинаково расположенными плоскими углами;}
    \item \textbf{или по равному плоскому углу, заключённому между двумя соответственно равными и одинаково расположенными двугранными углами;}
    \item \textbf{или по три соответственно равных и одинаково расположенных плоских угла;}
    \item \textbf{или по три соответственно равных и одинаково расположенных двугранных угла.}
\end{enumerate}

% Печатная с. 305.
\begin{enumerate}[label=\arabic*),ref=\arabic*]
    \item\label{item:trihedral-first-criterion} Пусть \(S\) и \(S_1\) --- два трёхгранных угла (рис.~\ref{fig:kiselev-stereo-350}), у которых
    \[
        \begin{gathered}
            \angle ASB=\angle A_1S_1B_1,\qquad \angle ASC=\angle A_1S_1C_1,\\
            \text{дв. уг. } AS=\text{дв. уг. } A_1S_1.
        \end{gathered}
    \]
    Наложим угол \(S_1\) на угол \(S\) так, чтобы у них совпали точка \(S_1\) с \(S\), прямая \(S_1A_1\) с \(SA\) и плоскость \(A_1S_1B_1\) с \(ASB\). Тогда ребро \(S_1B_1\) пойдёт по \(SB\) (по равенству углов \(A_1S_1B_1\) и \(ASB\)), плоскость \(A_1S_1C_1\) пойдёт по \(ASC\) (по равенству двугранных углов), а ребро \(S_1C_1\) --- по \(SC\) (по равенству углов \(A_1S_1C_1\) и \(ASC\)). Таким образом, трёхгранные углы совместятся всеми своими рёбрами. Значит, они будут равны.

    \begin{figure}[htbp]
        \centering
        \renewcommand{\thefigure}{350}
        \captionsetup{name=Рис.,labelsep=period}
        \begin{tikzpicture}[line cap=round,line join=round]
    \foreach \shift/\apex/\lefttag/\fronttag/\righttag in {0/S/A/B/C,4.95/S_1/A_1/B_1/C_1}{
        \begin{scope}[xshift=\shift cm]
            \coordinate (S) at (0,2.27);
            \coordinate (A) at (-1.45,-1.44);
            \coordinate (B) at (0,-1.85);
            \coordinate (C) at (1.46,-1.44);
            \draw[thick] (A) -- (S) -- (C);
            \draw[very thick] (S) -- (B);
            \draw[thick] (S) ++(-111.35:0.86) arc[start angle=-111.35,end angle=-90,radius=0.86];
            \draw[thick,densely dashed] (S) ++(-111.35:1.12)
                arc[start angle=-111.35,end angle=-68.52,radius=1.12];
            \path (S) -- (A) coordinate[pos=0.63] (mark);
            \draw[thick] (mark) -- ++(0.51,-0.04) (mark) -- ++(0.39,-0.30);
            \node[above] at (S) {\(\apex\)};
            \node[below left] at (A) {\(\lefttag\)};
            \node[below] at (B) {\(\fronttag\)};
            \node[below right] at (C) {\(\righttag\)};
        \end{scope}
    }
\end{tikzpicture}

        \caption{Первый признак равенства трёхгранных углов}\label{fig:kiselev-stereo-350}
    \end{figure}

    \item Второй признак доказывается наложением подобно \hyperref[item:trihedral-first-criterion]{первому}.

    \item Пусть \(S\) и \(S_1\) (рис.~\ref{fig:kiselev-stereo-351}) --- два трёхгранных угла, у которых плоские углы одного равны соответственно плоским углам другого и, кроме того, равные углы одинаково расположены.

    \begin{figure}[htbp]
        \centering
        \renewcommand{\thefigure}{351}
        \captionsetup{name=Рис.,labelsep=period}
        \begin{tikzpicture}[line cap=round,line join=round]
    \foreach \shift/\apex/\fronttag/\lefttag/\righttag/\centre in {0/S/A/B/C/O,5.02/S_1/A_1/B_1/C_1/O_1}{
        \begin{scope}[xshift=\shift cm]
            \coordinate (S) at (0,2.70);
            \coordinate (O) at (0,-1.20);
            \path (O) ++({1.38*cos(247)},{0.82*sin(247)}) coordinate (A);
            \path (O) ++({1.38*cos(164)},{0.82*sin(164)}) coordinate (B);
            \path (O) ++({1.38*cos(21)},{0.82*sin(21)}) coordinate (C);
            \draw[thick,densely dashed] (O) ++(1.38,0)
                arc[start angle=0,end angle=180,x radius=1.38,y radius=0.82];
            \draw[very thick] (O) ++(-1.38,0)
                arc[start angle=180,end angle=360,x radius=1.38,y radius=0.82];
            \draw[thick,densely dashed] (S) -- (O) (O) -- (A) (O) -- (B) (O) -- (C) (B) -- (C);

            \foreach \vertex in {A,B,C}{
                \path (S) -- (\vertex) coordinate[pos=1.14] (tip);
                \draw[thick] (S) -- (tip);
            }
            \draw[thick] (B) -- (A) -- (C);
            \draw[thick] (S) ++(-109.87:0.87) arc[start angle=-109.87,end angle=-96.58,radius=0.87];
            \draw[thick] (S) ++(-96.58:1.03) arc[start angle=-96.58,end angle=-70.30,radius=1.03];
            \draw[thick] (S) ++(-96.58:1.16) arc[start angle=-96.58,end angle=-70.30,radius=1.16];
            \draw[thick,densely dashed] (S) ++(-109.87:1.46)
                arc[start angle=-109.87,end angle=-70.30,radius=1.46];

            \node[above] at (S) {\(\apex\)};
            \node[below right] at (A) {\(\fronttag\)};
            \node[left] at (B) {\(\lefttag\)};
            \node[right] at (C) {\(\righttag\)};
            \node[fill=white,inner sep=1.3pt,right,xshift=2pt] at (O) {\(\centre\)};
        \end{scope}
    }
\end{tikzpicture}

        \caption{Трёхгранные углы с соответственно равными плоскими углами}\label{fig:kiselev-stereo-351}
    \end{figure}

    Отложим на всех рёбрах произвольные, но равные отрезки \(SA=SB=SC=S_1A_1=\ldots\) и построим треугольники \(ABC\) и \(A_1B_1C_1\). Из равенства треугольников \(ABS\) и \(A_1B_1S_1\) находим \(AB=A_1B_1\). Подобно этому из равенства других боковых треугольников выводим \(AC=A_1C_1\) и \(BC=B_1C_1\). Следовательно, \(\triangle ABC=\triangle A_1B_1C_1\). Опустим на плоскости этих треугольников перпендикуляры \(SO\) и \(S_1O_1\). Так как наклонные \(SA,SB\) и \(SC\) равны, то должны быть равны их проекции \(OA,OB\) и \(OC\). Значит, точка \(O\) --- центр окружности, описанной около треугольника \(ABC\). Точно
    % Печатная с. 306.
    так же точка \(O_1\) --- центр окружности, описанной около треугольника \(A_1B_1C_1\). У равных треугольников радиусы описанных окружностей равны. Значит, \(OB=O_1B_1\). Поэтому \(\triangle SBO=\triangle S_1B_1O_1\) (по гипотенузе и катету), и, следовательно, \(OS=O_1S_1\). Наложим теперь фигуру \(S_1A_1B_1C_1\) на фигуру \(SABC\) так, чтобы равные треугольники \(A_1B_1C_1\) и \(ABC\) совместились. Тогда совместятся описанные окружности и, следовательно, их центры \(O_1\) и \(O\). Вследствие этого перпендикуляр \(O_1S_1\) пойдёт по \(OS\), и точка \(S_1\) совпадёт с \(S\). Таким образом, трёхгранные углы совместятся всеми своими рёбрами. Значит, они равны.

    \item Четвёртый признак легко доказывается при помощи дополнительных углов. Если у двух трёхгранных углов соответственно равны и одинаково расположены \emph{двугранные углы}, то у дополнительных углов соответственно равны и одинаково расположены \emph{плоские углы} (\hyperref[para:supplementary-trihedral-angle-lemma]{см.~§~\ref*{para:supplementary-trihedral-angle-lemma}}). Следовательно, дополнительные углы равны. А если равны дополнительные углы, то равны и данные углы (\hyperref[para:equal-supplementary-trihedral-angles]{см.~§~\ref*{para:equal-supplementary-trihedral-angles}}).
\end{enumerate}

\kiselevparagraph{para:symmetric-polyhedral-angles} \textbf{Симметричные многогранные углы.} Как известно, вертикальные углы равны, если речь идёт об углах, образованных прямыми или плоскостями. Посмотрим, применимо ли это к многогранным углам.

Продолжим все рёбра угла \(SABCDE\) за вершину (рис.~\ref{fig:kiselev-stereo-352}). Тогда образуется другой многогранный угол \(SA_1B_1C_1D_1E_1\), который можно назвать \emph{вертикальным} по отношению к первому углу.

\begin{figure}[htbp]
    \centering
    \renewcommand{\thefigure}{352}
    \captionsetup{name=Рис.,labelsep=period}
    \begin{tikzpicture}[line cap=round,line join=round]
    \coordinate (S) at (0,0);
    \foreach \ray/\x/\y in {A/-2.12/-3.48,B/-0.84/-3.82,C/1.17/-3.72,D/2.29/-3.24,E/-0.09/-2.56}{
        \coordinate (\ray) at (\x,\y);
        \coordinate (\ray-opp) at ({-(\x)},{-(\y)});
    }
    \draw[very thick] (A) -- (A-opp) (B) -- (S) -- (C) (D) -- (D-opp) (S) -- (E-opp);
    \draw[thick,densely dashed] (E) -- (S) -- (B-opp) (S) -- (C-opp);
    \node[left,inner sep=3pt,xshift=-3pt] at (S) {\(S\)};
    \node[below left] at (A) {\(A\)};
    \node[below] at (B) {\(B\)};
    \node[below] at (C) {\(C\)};
    \node[below right] at (D) {\(D\)};
    \node[below right] at (E) {\(E\)};
    \node[above right] at (A-opp) {\(A_1\)};
    \node[above right] at (B-opp) {\(B_1\)};
    \node[above left] at (C-opp) {\(C_1\)};
    \node[above left] at (D-opp) {\(D_1\)};
    \node[above] at (E-opp) {\(E_1\)};
\end{tikzpicture}

    \caption{Симметричные углы \(SABCDE\) и \(SA_1B_1C_1D_1E_1\)}\label{fig:kiselev-stereo-352}
\end{figure}

Нетрудно видеть, что у обоих углов равны соответственно и плоские углы, и двугранные. Но и те и другие расположены \emph{в обратном порядке}. Действительно, если представим наблюдателя, который смотрит изнутри многогранного угла на его вершину, то рёбра \(SA,SB,SC,SD,SE\) будут казаться ему расположенными против движения часовой стрелки. Глядя же на угол \(SA_1B_1C_1D_1E_1\), он увидит рёбра \(SA_1,SB_1,\ldots\) расположенными по движению часовой стрелки.

Многогранные углы с соответственно равными плоскими и двугранными углами, расположенными в обратном порядке, в общем случае не могут совместиться при наложении. Значит, они не равны. Такие углы называются \emph{симметричными} (относительно вершины \(S\)).

\clearpage
\endgroup
  % Глава V
\endgroup

%%% Заключительные разделы
\chapter*{Заключение}                       % Заголовок
\addcontentsline{toc}{chapter}{Заключение}  % Добавляем его в оглавление

Благодарю авторов шаблона
\href{https://github.com/AndreyAkinshin/Russian-Phd-LaTeX-Dissertation-Template}{\textit{Russian-Phd-LaTeX-Dissertation-Template}}
за основу для вёрстки этой книги. Спасибо Сергею Седунову за то, что дал мне
эту идею, и нейросетям за помощь в её реализации.

Благодарю своего отца за любопытство и за образование, которое он помог мне получить.
  % Заключение
\printnomenclature[3.5cm] % Значение ширины столбца с обозначениями стоит подбирать вручную
    % Список сокращений и условных обозначений
\include{Book/dictionary}  % Словарь терминов
\clearpage                                  % В том числе гарантирует, что список литературы в оглавлении будет с правильным номером страницы
%\hypersetup{ urlcolor=black }               % Ссылки делаем чёрными
%\providecommand*{\BibDash}{}                % В стилях ugost2008 отключаем использование тире как разделителя
\urlstyle{rm}                               % ссылки URL обычным шрифтом
\ifdefmacro{\microtypesetup}{\microtypesetup{protrusion=false}}{} % не рекомендуется применять пакет микротипографики к автоматически генерируемому списку литературы
\insertbibliofull                           % Подключаем Bib-базы: все статьи единым списком
% Режим с подсписками
%\insertbiblioexternal                      % Подключаем Bib-базы: статьи, не являющиеся статьями автора по теме книги
% Для вывода выберите и расскомментируйте одно из двух
%\insertbiblioauthor                        % Подключаем Bib-базы: работы автора единым списком 
%\insertbiblioauthorgrouped                 % Подключаем Bib-базы: работы автора сгруппированные (ВАК, WoS, Scopus и т.д.)
\ifdefmacro{\microtypesetup}{\microtypesetup{protrusion=true}}{}
\urlstyle{tt}                               % возвращаем установки шрифта ссылок URL
%\hypersetup{ urlcolor={urlcolor} }          % Восстанавливаем цвет ссылок
  % Список литературы
\clearpage
\ifdefmacro{\microtypesetup}{\microtypesetup{protrusion=false}}{} % не рекомендуется применять пакет микротипографики к автоматически генерируемым спискам
\listoffigures  % Список изображений

%%% Список таблиц (ГОСТ Р 7.0.11-2011, 5.3.10)
\ifnum\totvalue{totalcount@table}>0\relax
    \clearpage
    \listoftables
\fi
\ifdefmacro{\microtypesetup}{\microtypesetup{protrusion=true}}{}
\newpage
       % Списки рисунков и таблиц

\setcounter{totalchapter}{\value{chapter}}  % Количество глав

\end{document}
